% Options for packages loaded elsewhere
%DIF LATEXDIFF DIFFERENCE FILE
%DIF DEL e419efab43344d7374f1827012e93d4a83f86b78
%DIF ADD 00ce2ac
% Baseline source: paper/manuscript_tmlr.md at e419efab43344d7374f1827012e93d4a83f86b78
% Revised source: paper/manuscript_tmlr_revisions.md at 00ce2ac
% Deletions are suppressed; additions and replacement text are shown in blue.
\PassOptionsToPackage{unicode}{hyperref}
\PassOptionsToPackage{hyphens}{url}
\documentclass[
]{article}
\usepackage{xcolor}
\usepackage[top=18mm,bottom=18mm,left=18mm,right=18mm]{geometry}
\usepackage{amsmath,amssymb}
\setcounter{secnumdepth}{-\maxdimen} % remove section numbering
\usepackage{iftex}
\ifPDFTeX
  \usepackage[T1]{fontenc}
  \usepackage[utf8]{inputenc}
  \usepackage{textcomp} % provide euro and other symbols
\else % if luatex or xetex
  \usepackage{unicode-math} % this also loads fontspec
  \defaultfontfeatures{Scale=MatchLowercase}
  \defaultfontfeatures[\rmfamily]{Ligatures=TeX,Scale=1}
\fi
\usepackage{lmodern}
\ifPDFTeX\else
  % xetex/luatex font selection
\fi
% Use upquote if available, for straight quotes in verbatim environments
\IfFileExists{upquote.sty}{\usepackage{upquote}}{}
\IfFileExists{microtype.sty}{% use microtype if available
  \usepackage[]{microtype}
  \UseMicrotypeSet[protrusion]{basicmath} % disable protrusion for tt fonts
}{}
\makeatletter
\@ifundefined{KOMAClassName}{% if non-KOMA class
  \IfFileExists{parskip.sty}{%
    \usepackage{parskip}
  }{% else
    \setlength{\parindent}{0pt}
    \setlength{\parskip}{6pt plus 2pt minus 1pt}}
}{% if KOMA class
  \KOMAoptions{parskip=half}}
\makeatother
\usepackage{color}
\usepackage{fancyvrb}
\newcommand{\VerbBar}{|}
\newcommand{\VERB}{\Verb[commandchars=\\\{\}]}
\DefineVerbatimEnvironment{Highlighting}{Verbatim}{commandchars=\\\{\}}
% Add ',fontsize=\small' for more characters per line
\newenvironment{Shaded}{}{}
\newcommand{\AlertTok}[1]{\textcolor[rgb]{1.00,0.00,0.00}{\textbf{#1}}}
\newcommand{\AnnotationTok}[1]{\textcolor[rgb]{0.38,0.63,0.69}{\textbf{\textit{#1}}}}
\newcommand{\AttributeTok}[1]{\textcolor[rgb]{0.49,0.56,0.16}{#1}}
\newcommand{\BaseNTok}[1]{\textcolor[rgb]{0.25,0.63,0.44}{#1}}
\newcommand{\BuiltInTok}[1]{\textcolor[rgb]{0.00,0.50,0.00}{#1}}
\newcommand{\CharTok}[1]{\textcolor[rgb]{0.25,0.44,0.63}{#1}}
\newcommand{\CommentTok}[1]{\textcolor[rgb]{0.38,0.63,0.69}{\textit{#1}}}
\newcommand{\CommentVarTok}[1]{\textcolor[rgb]{0.38,0.63,0.69}{\textbf{\textit{#1}}}}
\newcommand{\ConstantTok}[1]{\textcolor[rgb]{0.53,0.00,0.00}{#1}}
\newcommand{\ControlFlowTok}[1]{\textcolor[rgb]{0.00,0.44,0.13}{\textbf{#1}}}
\newcommand{\DataTypeTok}[1]{\textcolor[rgb]{0.56,0.13,0.00}{#1}}
\newcommand{\DecValTok}[1]{\textcolor[rgb]{0.25,0.63,0.44}{#1}}
\newcommand{\DocumentationTok}[1]{\textcolor[rgb]{0.73,0.13,0.13}{\textit{#1}}}
\newcommand{\ErrorTok}[1]{\textcolor[rgb]{1.00,0.00,0.00}{\textbf{#1}}}
\newcommand{\ExtensionTok}[1]{#1}
\newcommand{\FloatTok}[1]{\textcolor[rgb]{0.25,0.63,0.44}{#1}}
\newcommand{\FunctionTok}[1]{\textcolor[rgb]{0.02,0.16,0.49}{#1}}
\newcommand{\ImportTok}[1]{\textcolor[rgb]{0.00,0.50,0.00}{\textbf{#1}}}
\newcommand{\InformationTok}[1]{\textcolor[rgb]{0.38,0.63,0.69}{\textbf{\textit{#1}}}}
\newcommand{\KeywordTok}[1]{\textcolor[rgb]{0.00,0.44,0.13}{\textbf{#1}}}
\newcommand{\NormalTok}[1]{#1}
\newcommand{\OperatorTok}[1]{\textcolor[rgb]{0.40,0.40,0.40}{#1}}
\newcommand{\OtherTok}[1]{\textcolor[rgb]{0.00,0.44,0.13}{#1}}
\newcommand{\PreprocessorTok}[1]{\textcolor[rgb]{0.74,0.48,0.00}{#1}}
\newcommand{\RegionMarkerTok}[1]{#1}
\newcommand{\SpecialCharTok}[1]{\textcolor[rgb]{0.25,0.44,0.63}{#1}}
\newcommand{\SpecialStringTok}[1]{\textcolor[rgb]{0.73,0.40,0.53}{#1}}
\newcommand{\StringTok}[1]{\textcolor[rgb]{0.25,0.44,0.63}{#1}}
\newcommand{\VariableTok}[1]{\textcolor[rgb]{0.10,0.09,0.49}{#1}}
\newcommand{\VerbatimStringTok}[1]{\textcolor[rgb]{0.25,0.44,0.63}{#1}}
\newcommand{\WarningTok}[1]{\textcolor[rgb]{0.38,0.63,0.69}{\textbf{\textit{#1}}}}
\usepackage{longtable,booktabs,array}
\newcounter{none} % for unnumbered tables
\usepackage{calc} % for calculating minipage widths
% Correct order of tables after \paragraph or \subparagraph
\usepackage{etoolbox}
\makeatletter
\patchcmd\longtable{\par}{\if@noskipsec\mbox{}\fi\par}{}{}
\makeatother
% Allow footnotes in longtable head/foot
\IfFileExists{footnotehyper.sty}{\usepackage{footnotehyper}}{\usepackage{footnote}}
\makesavenoteenv{longtable}
\usepackage{graphicx}
\makeatletter
\newsavebox\pandoc@box
\newcommand*\pandocbounded[1]{% scales image to fit in text height/width
  \sbox\pandoc@box{#1}%
  \Gscale@div\@tempa{\textheight}{\dimexpr\ht\pandoc@box+\dp\pandoc@box\relax}%
  \Gscale@div\@tempb{\linewidth}{\wd\pandoc@box}%
  \ifdim\@tempb\p@<\@tempa\p@\let\@tempa\@tempb\fi% select the smaller of both
  \ifdim\@tempa\p@<\p@\scalebox{\@tempa}{\usebox\pandoc@box}%
  \else\usebox{\pandoc@box}%
  \fi%
}
% Set default figure placement to htbp
\def\fps@figure{htbp}
\makeatother
\usepackage{svg}
% definitions for citeproc citations
\NewDocumentCommand\citeproctext{}{}
\NewDocumentCommand\citeproc{mm}{%
  \begingroup\def\citeproctext{#2}\cite{#1}\endgroup}
\makeatletter
 % allow citations to break across lines
 \let\@cite@ofmt\@firstofone
 % avoid brackets around text for \cite:
 \def\@biblabel#1{}
 \def\@cite#1#2{{#1\if@tempswa , #2\fi}}
\makeatother
\newlength{\cslhangindent}
\setlength{\cslhangindent}{1.5em}
\newlength{\csllabelwidth}
\setlength{\csllabelwidth}{3em}
\newenvironment{CSLReferences}[2] % #1 hanging-indent, #2 entry-spacing
 {\begin{list}{}{%
  \setlength{\itemindent}{0pt}
  \setlength{\leftmargin}{0pt}
  \setlength{\parsep}{0pt}
  % turn on hanging indent if param 1 is 1
  \ifodd #1
   \setlength{\leftmargin}{\cslhangindent}
   \setlength{\itemindent}{-1\cslhangindent}
  \fi
  % set entry spacing
  \setlength{\itemsep}{#2\baselineskip}}}
 {\end{list}}
\usepackage{calc}
\newcommand{\CSLBlock}[1]{\hfill\break\parbox[t]{\linewidth}{\strut\ignorespaces#1\strut}}
\newcommand{\CSLLeftMargin}[1]{\parbox[t]{\csllabelwidth}{\strut#1\strut}}
\newcommand{\CSLRightInline}[1]{\parbox[t]{\linewidth - \csllabelwidth}{\strut#1\strut}}
\newcommand{\CSLIndent}[1]{\hspace{\cslhangindent}#1}
\setlength{\emergencystretch}{3em} % prevent overfull lines
\providecommand{\tightlist}{%
  \setlength{\itemsep}{0pt}\setlength{\parskip}{0pt}}
% Mirrors the base page geometry used by the HTML/CSS export path.
% Figure bleed remains an HTML/CSS-only styling rule for now; if LaTeX
% export later needs the same effect, extend this header or a custom
% Pandoc template rather than trying to reuse the CSS file directly.
\usepackage{graphicx}
\setkeys{Gin}{width=\linewidth,keepaspectratio}
\usepackage{bookmark}
\IfFileExists{xurl.sty}{\usepackage{xurl}}{} % add URL line breaks if available
\urlstyle{same}
\hypersetup{
  hidelinks,
  pdfcreator={LaTeX via pandoc}}

\author{}
\date{}
%DIF PREAMBLE EXTENSION ADDED BY LATEXDIFF
%DIF CFONT PREAMBLE %DIF PREAMBLE
\RequirePackage{color}\definecolor{RED}{rgb}{1,0,0}\definecolor{BLUE}{rgb}{0,0,1} %DIF PREAMBLE
\providecommand{\DIFadd}[1]{{\protect\color{blue} #1}} %DIF PREAMBLE
\providecommand{\DIFdel}[1]{{\protect\color{red} \scriptsize #1}} %DIF PREAMBLE
%DIF SAFE PREAMBLE %DIF PREAMBLE
\providecommand{\DIFaddbegin}{} %DIF PREAMBLE
\providecommand{\DIFaddend}{} %DIF PREAMBLE
\providecommand{\DIFdelbegin}{} %DIF PREAMBLE
\providecommand{\DIFdelend}{} %DIF PREAMBLE
\providecommand{\DIFmodbegin}{} %DIF PREAMBLE
\providecommand{\DIFmodend}{} %DIF PREAMBLE
%DIF FLOATSAFE PREAMBLE %DIF PREAMBLE
\providecommand{\DIFaddFL}[1]{\DIFadd{#1}} %DIF PREAMBLE
\providecommand{\DIFdelFL}[1]{\DIFdel{#1}} %DIF PREAMBLE
\providecommand{\DIFaddbeginFL}{} %DIF PREAMBLE
\providecommand{\DIFaddendFL}{} %DIF PREAMBLE
\providecommand{\DIFdelbeginFL}{} %DIF PREAMBLE
\providecommand{\DIFdelendFL}{} %DIF PREAMBLE
%DIF END PREAMBLE EXTENSION ADDED BY LATEXDIFF

\begin{document}

\section{Floatsom Paper}\label{floatsom-paper}

\subsection{Title}\label{title}

FloatSOM: GPU Accelerated, Distributed, Topology-Flexible
Self-Organizing Maps

\subsection{Authors}\label{authors}

Anonymous Authors

\subsection{Abstract}\label{abstract}

GPU-accelerated Self-Organizing Map (SOM) implementations are among the
most competitive options for large-scale SOM analysis, but growing
dataset sizes increasingly challenge their practical use because
workloads no longer fit cleanly within device-memory limits. We
introduce FloatSOM, a SOM framework for scalable training and deployment
that supports multi-GPU execution, out-of-memory disk-backed streaming,
and  \DIFaddbegin \DIFadd{scalable training-time graph }\DIFaddend topologies beyond regular lattices. We
evaluate FloatSOM on 14 synthetic and real benchmark datasets together
with controlled speed scaling benchmarks, and show that these
 \DIFaddbegin \DIFadd{graph-based }\DIFaddend topologies, combined with topology-aware hyperparameter
fine-tuning, yield lower quantization error than current
state-of-the-art SOM baselines. FloatSOM also sustains this performance
at large scale with high-throughput distributed execution; in the
largest benchmark, it trains a 1024-node SOM network on 1,000,000,000
samples with 50 features in 6.16 minutes on 8 GPUs across two separate
high-performance-computing nodes.

\subsection{1. Introduction}\label{introduction}

Self-Organizing Map (SOM), originally introduced by Kohonen (Kohonen
1990), is an unsupervised machine-learning method that uses competitive
learning to organize nodes such that they capture the topology of the
data. In practice, this topology-preserving representation means SOMs
are commonly used to map dataset topology, produce
dimensionality-reduced visualizations, and conduct clustering at scale
(Kangas et al. 1990). As dataset size and heterogeneity increase, the
computational requirements of SOMs grow in both time and memory.  \DIFaddbegin \DIFadd{Recent
tools such as aweSOM have improved single-node serial-online SOM
execution through CPU/GPU acceleration and ensemble stacking (Ha et al.
2025). However, many current implementations }\DIFaddend remain constrained to
single-device workloads that must fit within video random-access memory
(VRAM, GPU memory), with limited support for distributed compute,
out-of-core execution, and modern GPU orchestration.

SOM topology presents a second limitation. Classical SOMs are
traditionally trained on regular rectangular or hexagonal lattices
because fixed grids make neighborhood definition, visualization, and
optimization straightforward. However, these regular lattices also
impose a strong geometric prior on the learned representation.
Accordingly, prior work on dynamic, growing, and graph-structured SOM
variants reflects a long-standing recognition that fixed lattices are
not always the best match for irregular data geometry (Alahakoon et al.
2000; Vasighi and Amini 2017; Kangas et al. 1990). However, these
alternatives have generally not been developed or evaluated in the
high-throughput, large-sample regime targeted by modern GPU-enabled
applications and are broadly impractical for deployment at scale.

FloatSOM is designed to address these combined systems and topology
limitations. Within FloatSOM, we implement distributed multi-GPU
execution, out-of-memory disk-backed streaming, and scalable
topology-flexible training beyond standard fixed lattices. Specifically,
we implement highly scalable minimum-spanning-tree (MST) and
relative-neighborhood-graph (RNG) topologies \DIFaddbegin \DIFadd{that define the SOM
neighborhood during training rather than being overlaid on a completed
map}\DIFaddend . We also evaluate multiple sampling strategies as a route to further
computational acceleration and derive fine-tuned hyperparameter
configurations across diverse datasets and recommended operating
regimes. FloatSOM is suitable for both high-performance-computing (HPC)
operation and consumer-grade desktop GPUs.

\subsection{2. Related Work}\label{related-work}

Related work on practical SOM deployment spans software implementations,
sampling-efficient training, topology design, and model selection.

\subsubsection{2.1 Open-Source SOM Implementations and
Systems}\label{open-source-som-implementations-and-systems}

 \DIFaddbegin \DIFadd{Open-source }\DIFaddend SOM libraries range from lightweight  \DIFaddbegin \DIFadd{implementations to
systems-oriented packages. MiniSom implements classical serial-online
training }\DIFaddend (Vettigli 2018),  \DIFaddbegin \DIFadd{while aweSOM adds CPU/GPU acceleration and
ensemble stacking for large single-node workloads (Ha et al. 2025).
XPySOM instead provides GPU-accelerated batch training and is the
closest executable comparator to FloatSOM's Python/GPU training pathway
}\DIFaddend (Mancini et al. 2020).
 \DIFaddbegin

\DIFaddend Somoclu and GigaSOM provide  \DIFaddbegin \DIFadd{additional parallel-systems context }\DIFaddend (Wittek
et al. 2017; Kratochvíl et al. 2020).  \DIFaddbegin \DIFadd{GigaSOM.jl, for example, trained a
\(32 \times 32\) SOM on 1,167,129,317 cells within a larger Julia
analysis completed in under 25 minutes on an 11-node, 256-core CPU
cluster (Kratochvíl et al. 2020). These systems differ from FloatSOM in
training regime, execution model, language, or hardware, so we treat
their published results as systems context rather than controlled
benchmarks.
}\DIFaddend

\subsubsection{2.2 Sampling Methodologies for SOM
Training}\label{sampling-methodologies-for-som-training}

Classical online and batch SOM training schemes traditionally sample
every data-point in every training iteration (Kohonen 2001; Liu et al.
2023). Consequently, one obvious route to improving SOM speed and
scalability is subsampling, which reduces the amount of data used to
train the final SOM.

To date, numerous subsampling strategies have been proposed. Beyond
naive random sampling, guided subsampling methods include hierarchical
dynamic subset selection SOM (HDSSSOM), which concentrates computation
on difficult or stale regions of the data (Wetmore et al. 2005), and
adaptive sampling strategies for design-space exploration ({Ito et al.}
2016). However, systematically benchmarked and openly maintained
implementations that compare full-data, random, and guided sampling
within a modern high-performance SOM workflow remain limited.

\subsubsection{2.3 SOM Lattice and Adaptive
Topologies}\label{som-lattice-and-adaptive-topologies}

Most practical SOM implementations retain regular rectangular or
hexagonal lattices because they simplify neighborhood indexing,
visualization, and vectorized updates (Kohonen 1990; Kohonen 2013).
Hexagonal lattices are often preferred in the literature and serve as
the regular-topology baseline in this manuscript (White and Kiester
2008; Kohonen 2013; Forest et al. 2020).

The SOM literature has also explored alternatives to fixed lattices,
including dynamic maps and graph-structured neighborhoods (Vasighi and
Amini 2017; Spanakis and Weiss 2016; Kangas et al. 1990; Jang et al.
2009).  \DIFaddbegin \DIFadd{MSTs have previously appeared in SOM analyses, but prior uses
generally treat the MST as an interpretive structure over an already
trained map rather than as the neighborhood relation that drives SOM
learning. For example, Jang et al.~use MSTs for interpretation, subnode
embedding, and map-shape assessment, while FlowSOM overlays an MST on
trained SOM codes to visualize relationships among metaclusters (Jang et
al. 2009; Van Gassen et al. 2015). This is distinct from the
training-time role used here: in FloatSOM, the MST is the operative
neighborhood graph during SOM updates, is recalculated from the evolving
node-weight geometry, and directly changes the update influence matrix
used during learning. Thus, current deployed MST uses in SOM workflows
generally configure post hoc connections on trained maps, whereas
FloatSOM uses MST/RNG graphs as the training neighborhood itself. A post
hoc MST overlay would also leave the quantization error of the
underlying trained SOM unchanged, because \(QE\) is determined by
distances between samples and the fixed learned prototypes; the \(QE\)
improvements reported here therefore require graph-mediated training
that changes those prototype locations. Earlier SOM variants also
discussed MST-defined neighborhoods during learning (Kangas et al.
1990), but }\DIFaddend these alternatives have not been assessed on large-scale
datasets and do not have implementations that are either publicly
available or suitable for distributed GPU computation.  \DIFaddbegin \DIFadd{We have not
identified prior work implementing dynamically refreshed MST or RNG
topologies as scalable GPU-compatible SOM training neighborhoods.
FloatSOM's topology contribution is therefore a scalable GPU-compatible
implementation and large-scale quantification of refreshed graph-based
SOM training, rather than a post hoc MST overlay on a conventional
trained SOM.
}\DIFaddend

Relative Neighborhood Graphs (RNGs) (Toussaint 1980) are of particular
interest  \DIFaddbegin \DIFadd{here. To our knowledge, dynamically refreshed RNG neighborhoods
have not previously been used as a SOM training topology}\DIFaddend ; we return to
the full rationale and implementation for RNG in Section 3.2.2.

\subsubsection{2.4 Hyperparameter Optimization and Fair
Comparison}\label{hyperparameter-optimization-and-fair-comparison}

SOM performance depends strongly on parameters such as map size,
initialization, learning-rate schedule, and neighborhood schedule. Prior
work shows that these choices can substantially affect observed
performance (Akinduko et al. 2016; Forest et al. 2020). Because recent
SOM system comparisons often emphasize implementation-level performance
against existing tools or reference configurations (Mancini et al. 2020;
Kratochvíl et al. 2020), comparisons based only on untuned  \DIFaddbegin \DIFadd{settings }\DIFaddend can
be difficult to interpret.

More generally, modern machine-learning workflows increasingly rely on
automated hyperparameter optimization rather than manual tuning alone.
Tools such as Optuna provide bounded search over large hyperparameter
spaces and can support multi-objective optimization, allowing parameter
settings to be selected with respect to several benchmark criteria
simultaneously rather than collapsed into a single score (Akiba et al.
2019). Accordingly, we test both tuned and untuned reference and
FloatSOM implementations to distinguish performance under untuned
 \DIFaddbegin \DIFadd{settings }\DIFaddend from performance obtained under more suitable hyperparameters.

\subsection{3. Methods}\label{methods}

FloatSOM combines four methodological components: sample selection,
topology definition, SOM updates, and the compute architecture used to
execute them. The framework supports random, full, and HDSSSOM sampling;
rectangular, hexagonal, MST, and RNG topologies; and execution modes
ranging from local GPU training to single-node and multi-node multi-GPU
operation. Hyperparameter optimization is treated as a separate layer
applied across these configurations. Figure 1 summarizes this design.

\begin{figure}
\centering
\includesvg[width=0.5\linewidth,height=\textheight,keepaspectratio]{assets_manual/figures/fig_1.svg}
\caption{Figure 1}
\end{figure}

\emph{Figure 1. Schematic overview of the FloatSOM methods framing used
in this manuscript. Sample selection and topology definition have
configurable components that feed into the standard SOM training step,
while the compute architecture determines how that same training
procedure is executed in practice. The options shown here summarize the
FloatSOM configurations discussed in the following subsections.}

\subsubsection{3.1 Sampling Selector
Mathematics}\label{sampling-selector-mathematics}

This section formalizes the sampling policies evaluated in this
manuscript.

Let the full dataset be \(X=\{x_i\}_{i=1}^{N}\), where \(N\) is the
total number of samples in the dataset. Let \(m\) denote the number of
samples presented to the SOM in a given training iteration, or the
sampling budget. This budget is either fixed directly or determined as a
proportion \(\rho\) of the dataset:

\[
m=
\begin{cases}
m_0, & \text{if a fixed budget is specified},\\[4pt]
\max\!\left(1,\left\lfloor N\rho \right\rfloor\right), & \text{if a dataset proportion is used}.
\end{cases}
\tag{1}
\]

For index sets \(\mathcal{I}_t^{(s)}\) at iteration \(t\), the
implemented full and random selectors are:

\[
\mathcal{I}_t^{\mathrm{full}}=\{1,\dots,N\},
\qquad
\mathcal{I}_t^{\mathrm{random}} \sim \mathrm{Unif}\!\left(\left\{I\subseteq\{1,\dots,N\}:|I|=m\right\}\right),\ m<N.
\tag{2}
\]

If \(m\ge N\), random returns the full dataset (no subsampling).
Otherwise, to avoid biasing training toward a single fixed subsample,
the random selector redraws \(\mathcal{I}_t^{\mathrm{random}}\) from the
entire dataset at every training iteration. The selected training subset
is \(X_t^{(s)}=\{x_i:i\in\mathcal{I}_t^{(s)}\}\). Because random
subsampling draws uniformly from size-\(m\) subsets, we are able to
better obtain coverage over the full dataset.

HDSSSOM is implemented in multi-GPU-compatible form as an adaptive
sampler that preferentially revisits difficult, under-trained, or stale
samples (Wetmore et al. 2005).

\subsubsection{3.2 Topology Definition}\label{topology-definition}

After initialization, neighborhood relations are defined by the selected
topology. For regular-lattice baselines, we support both grid and
hexagonal layouts. Consistent with prior SOM guidance, we treat
hexagonal as the standard reference topology in this manuscript. We
implement the hexagonal lattice topology in accordance with (Vettigli
2018). For MST and RNG, neighborhood structure is derived from the
current node-weight geometry using the methodologies below.

\DIFaddbegin \DIFadd{These alternatives change which prototypes are coupled during learning.
A fixed hexagonal lattice assigns the same predetermined neighborhood
pattern independently of the learned data-space geometry. Moving one
prototype can therefore influence lattice neighbors that are not locally
related in data space, potentially displacing otherwise useful
prototypes. An MST minimizes this coupling and gives prototypes greater
freedom to redistribute along irregular or elongated structures, but its
tree constraint can omit additional locally appropriate connections in
dense regions. RNG is less free than MST because a node can influence
several neighbors; unlike the fixed lattice, however, those additional
connections are supported by the evolving prototype geometry. It can
therefore remain sparse where appropriate and form a more mesh-like
neighborhood in concentrated regions.
}

\DIFaddend Note that regardless of the selected topology, FloatSOM uses the same
node-weight initialization options. Initialization determines only the
starting node weights; neighborhood relations are applied afterward
according to the selected topology to establish node connections,
keeping the initial state comparable across regular-lattice and
graph-based runs.

\paragraph{3.2.1 MST Topology
Implementation}\label{mst-topology-implementation}

The MST topology replaces fixed lattice neighborhood distance with graph
hop distance on a minimum spanning tree built from current SOM nodes.
\DIFaddbegin \DIFadd{The MST is therefore part of the SOM training rule: it determines which
nodes receive neighborhood influence during each topology-refresh
interval, rather than serving as a visualization or analysis layer after
training. }\DIFaddend For \(P\) nodes in feature dimension \(d\), the pairwise
node-distance matrix is formed with the standard squared-distance Gram
identity, which avoids 3D broadcast tensors and preserves \(O(P^2 d)\)
dense linear-algebra structure.

In a standard lattice SOM, neighborhood distance is determined by fixed
node grid coordinates, which are assigned before training and are
independent of the learned node-weight geometry in data space. MST
instead uses the node weight vectors to locate the nodes in data space,
then calculates the minimum spanning tree between those node locations
(Kruskal 1956). The node distances therefore determine which tree edges
are selected, while the neighborhood distance used by the Gaussian
update is the graph hop distance along the resulting tree.
Topology-derived influence matrices are cached and refreshed according
to the topology-refresh schedule.

At the beginning of training, the topology is refreshed every iteration:
the previous tree edges are discarded, pairwise distances are
recalculated from the updated node weights, and the MST is recomputed
from those updated node locations. The refresh rate then decays as
training proceeds, so topology updates become less frequent later in
training. If the current iteration does not trigger recomputation, the
previous graph state and cached influences are reused. Otherwise, the
MST is recalculated and the influence cache is rebuilt.

\begin{Shaded}
\begin{Highlighting}[]
\NormalTok{Input: node weights W\_t, iteration t, topology{-}refresh policy}
\NormalTok{Output: topology state (E\_t, g\_t, cached influences)}
\NormalTok{ 1: Query the refresh policy for the current iteration}
\NormalTok{ 2: if no topology refresh is due then}
\NormalTok{ 3:         return previous topology state}
\NormalTok{ 4: end if}
\NormalTok{ 5: Compute pairwise squared node distances on GPU}
\NormalTok{ 6: Transfer distances to CPU and run Kruskal to obtain MST edges E\_t}
\NormalTok{ 7: Build adjacency from E\_t}
\NormalTok{ 8: Compute all{-}pairs graph hop distances g\_t with chunked GPU Floyd{-}Warshall}
\NormalTok{ 9: Deduplicate active radii and rebuild/update cached influence maps}
\NormalTok{10: Commit E\_t, g\_t, and cache state}
\NormalTok{11: return topology state}
\end{Highlighting}
\end{Shaded}

\emph{Algorithm 1. Dynamic MST topology update with refresh-triggered
recomputation and cached influence reuse.}

\paragraph{3.2.2 RNG Topology
Implementation}\label{rng-topology-implementation}

Our RNG topology constructs a Relative Neighborhood Graph over current
node distances using the standard RNG criterion (Toussaint 1980). Unlike
MST, RNG is not restricted to a single spanning-tree backbone, so
multiple locally supported neighborhood relations can be retained.
Conceptually, the absence of this restriction permits the RNG topology
to represent both tree-like and denser mesh-like local structures
simultaneously throughout the training process. Consequently, we
hypothesise that this greater freedom in topology will allow RNG-based
SOM maps to better represent a variety of distributions.

\begin{Shaded}
\begin{Highlighting}[]
\NormalTok{Input: node weights W\_t, iteration t, topology{-}refresh policy}
\NormalTok{Output: topology state (E\_t, g\_t, cached influences)}
\NormalTok{ 1: Query the refresh policy for the current iteration}
\NormalTok{ 2: if no topology refresh is due then}
\NormalTok{ 3:         return previous topology state}
\NormalTok{ 4: end if}
\NormalTok{ 5: Compute pairwise squared node distances on GPU}
\NormalTok{ 6: Evaluate RNG candidate elimination in chunks using the blocker test}
\NormalTok{ 7: Retain surviving RNG edges E\_t and build adjacency}
\NormalTok{ 8: Compute all{-}pairs graph hop distances g\_t with chunked GPU Floyd{-}Warshall}
\NormalTok{ 9: Deduplicate active radii and rebuild/update cached influence maps}
\NormalTok{10: Commit E\_t, g\_t, and cache state}
\NormalTok{11: return topology state}
\end{Highlighting}
\end{Shaded}

\emph{Algorithm 2. Dynamic RNG topology update with refresh-triggered
recomputation and cached influence reuse.}

\subsubsection{3.3 Multi-GPU + Larger-Than-Memory Methodology and
Implementation}\label{multi-gpu-larger-than-memory-methodology-and-implementation}

FloatSOM is implemented for distributed GPU execution, with support for
RAM (Random Access Memory, CPU memory) spilling and disk-backed
execution when datasets exceed available VRAM and RAM, respectively.

\begin{figure}
\centering
\includesvg[width=0.5\linewidth,height=\textheight,keepaspectratio]{assets_manual/figures/fig_2.svg}
\caption{Figure 2}
\end{figure}

\emph{Figure 2. Multi-GPU data loading and NCCL synchronization
schematic. In disk-backed mode, data are sharded to worker-local
storage, loaded chunk-wise into pinned host memory, transferred to GPU
with transfer overlapped with compute, processed locally, and
synchronized by NCCL all-reduce before one weight update per iteration.
In RAM mode, data are instead pre-sharded into worker-local CPU RAM and
follow the same pinned-memory-to-GPU path without disk reads.}

\paragraph{3.3.1 General Multi-GPU Logic}\label{general-multi-gpu-logic}

Distributed execution uses Ray actors with one GPU per worker and NCCL
collectives for synchronous aggregation (Moritz et al. 2018). At each
iteration, every worker processes its assigned shard locally in
\texttt{n\_chunks}, accumulates worker-local statistics across those
chunks, and synchronizes once per iteration.

For worker \(g \in \{1,\dots,G\}\), let \(X_{t,g}^{(s)}\) denote the
worker-local shard of the selected iteration subset \(X_t^{(s)}\). Let
\(j\) index SOM nodes, let \(w_j^{(t)}\) denote the weight vector of
node \(j\) at iteration \(t\), let \(b(x)\) denote the best-matching
unit (BMU) of sample \(x\) under the current weights, let
\(h_{j,b(x)}^{(t)}\) denote the iteration-\(t\) neighborhood influence
between node \(j\) and the BMU (Best Matching Unit)of \(x\), and let
\(\eta_t\) denote the learning rate at iteration \(t\). The local
accumulators are:

\[
\begin{aligned}
 U_j^{(g)} &= \sum_{x \in X_{t,g}^{(s)}} \eta_t\,h_{j,b(x)}^{(t)}(x-w_j^{(t)}), \\
 H_j^{(g)} &= \sum_{x \in X_{t,g}^{(s)}} h_{j,b(x)}^{(t)}.
\end{aligned}
\tag{3}
\]

\(U_j^{(g)}\) is worker \(g\)'s summed learning-rate-scaled displacement
for node \(j\), and \(H_j^{(g)}\) is the corresponding summed
neighborhood influence used to normalize that displacement.

Global synchronized accumulators are:

\[
\begin{aligned}
U_j &= \sum_{g=1}^{G} U_j^{(g)}, \\
H_j &= \sum_{g=1}^{G} H_j^{(g)}.
\end{aligned}
\tag{4}
\]

In implementation, \(U_j^{(g)}\) and \(H_j^{(g)}\) are accumulated
across the worker's \texttt{n\_chunks} and synchronized once per
iteration. After synchronization, \(U_j\) and \(H_j\) define the global
update numerator and normalization denominator for node \(j\). Under
weighted normalization, the normalized update is given by \(U_j/H_j\) up
to numerical safeguards; optional momentum is then added to this
normalized update before applying the resulting change to the node
weight \(w_j^{(t)}\).

\DIFaddbegin \DIFadd{For \(P\) nodes and \(d\) features in float32, the two update
accumulators contain \(P d+P\) values, so their logical all-reduce
payload is \(4P(d+1)\) bytes per worker per iteration, plus one 4-byte
sample count. A ring all-reduce sends and receives approximately
\(2(G-1)/G\) times that payload per worker for \(G\) workers. This
communication is independent of the number of observations \(N\);
increasing \(N\) instead increases worker-local computation and data
staging. With observation chunks of at most \(C\) rows, the dominant
bounded worker arrays scale as \(O(Cd+Pd)\) rather than requiring the
full \(O(Nd)\) dataset in GPU memory. Graph-topology distance and
influence structures add terms that scale with \(P^2\) and can be tiled
or spilled when necessary, which explains why increasing node count is a
less favorable scaling regime than increasing sample count.
}

\DIFaddend \paragraph{3.3.2 Multi-GPU Implementation
Details}\label{multi-gpu-implementation-details}

As shown in Fig. 2, each worker uses a chunked loading path from CPU
memory to GPU memory. In streaming mode, data are distributed to
worker-local disk shards and then read chunk-by-chunk into pinned host
memory before transfer to GPU; in RAM mode, data are pre-sharded into
worker-local CPU RAM and follow the same pinned-memory path without disk
reads. FloatSOM overlaps data transfer and compute across CUDA streams,
uses JIT-compiled kernels (Just-In-Time) for the core BMU and update
steps, and synchronizes worker-local accumulators once per iteration
with NCCL all-reduce. Weights therefore remain resident on worker GPUs
across iterations. For large topologies that would otherwise exceed
VRAM, FloatSOM also supports spilling the associated graph-distance or
influence structures from VRAM to system RAM.

\paragraph{3.3.3 Larger-Than-Memory Capable Topology
Updates}\label{larger-than-memory-capable-topology-updates}

Additional larger-than-memory support for topology updates is provided
through topological chunking. Topological chunking applies the same idea
to topology-side computations within each worker. When graph-distance or
influence structures would otherwise exceed a worker's VRAM, those
computations are tiled and evaluated in bounded pieces rather than
materialized at once. This topology-side chunking is not depicted in the
Fig. 2 data path, but it follows the same per-worker bounded-memory
execution rule.

\subsubsection{3.4 Multi-Objective Hyperparameter
Optimization}\label{multi-objective-hyperparameter-optimization}

We use Optuna as an automated multi-objective hyperparameter
optimization tool to derive near-optimal performance and corresponding
hyperparameters for each FloatSOM configuration under a given sampling
and topology combination (Akiba et al. 2019).

\subsection{4. Experimental Setup}\label{experimental-setup}

We use two benchmark protocols: an Optuna quality benchmark and a speed
scaling benchmark. The first evaluates algorithmic quality and tuned
attainable performance, and the second evaluates runtime and distributed
scaling behavior. All production Optuna and speed benchmarks reported in
this manuscript were executed on Gadi at the National Computational
Infrastructure (NCI), Australia, on gpuvolta nodes (\emph{{HPC Systems}
\textbar{} {NCI}}, n.d.), using a consistent multi-GPU environment
across runs.

\subsubsection{4.1 Optuna benchmark
protocol}\label{optuna-benchmark-protocol}

We use Optuna-based multi-objective optimization to determine the best
attainable performance and corresponding hyperparameters for each
sampling (full and random) and topology (hexagonal, MST, and RNG)
combination, optimizing jointly for train and holdout quantization error
(\(QE_T\) and \(QE_H\)). In this protocol, each
dataset-topology-sampling configuration is optimized for 200 trials
across 10 seeds under variant-appropriate search constraints. The
reported runs used the standard in-memory batch path rather than the
Ray-distributed execution stack. This corresponds to:

\[
14_{\text{Datasets}} \times 10_{\text{Seeds}} \times 3_{\text{Topologies}} \times 2_{\text{SamplingMethods}} \times 200_{\text{Trials}}= 168{,}000 \text{ Runs}
\]

 \DIFaddbegin \DIFadd{The neighborhood-radius search space was shared across topology
families. In particular, \texttt{initial\_radius} was an
Optuna-optimized parameter for hexagonal, MST, and RNG runs, with the
same search interval of 0.5 to 10.0 in each case. Thus, the hexagonal
topology comparisons below use a tuned hexagonal comparator rather than
a default-radius hexagonal baseline.
}

\DIFadd{The main sampling comparison }\DIFaddend in Section 5.2  \DIFaddbegin \DIFadd{compares }\DIFaddend full versus random
 \DIFaddbegin \DIFadd{sampling within matched hexagonal runs}\DIFaddend . Topology comparisons in Sections
5.3-5.4 use  \DIFaddbegin \DIFadd{full-sampling }\DIFaddend runs across hexagonal, MST, and RNG.
Otherwise, the remaining \DIFaddbegin \DIFadd{quality }\DIFaddend analyses use full sampling. A separate
focused HDSSSOM pilot is reported in Section 5.2 and summarized in
Supplementary Table S2. This pilot additionally included the HDSSSOM
sampling methodology alongside the full and random sampling
methodologies.

\paragraph{4.1.1 Optuna benchmark datasets and
preprocessing}\label{optuna-benchmark-datasets-and-preprocessing}

The Optuna benchmark uses 14 synthetic and real scikit-learn datasets
spanning low- and high-dimensional regimes (Pedregosa et al. 2011);
dataset metadata are listed in Supplementary Table S1. Synthetic
datasets are generated according to the random seed, whereas real
datasets are loaded directly from sklearn. Across the Optuna protocol,
inputs are standardized feature-wise to zero mean and unit variance
using \texttt{StandardScaler}, then deterministically permuted and split
into fixed 70/30 train-holdout partitions. We retain both partitions to
evaluate both training-set representation and transfer to unseen
samples.

\paragraph{4.1.2 Optuna quality metrics}\label{optuna-quality-metrics}

Our primary quality metric is Quantization Error (\(QE\)) (Kohonen
2001). We report \(QE\) on both the training partition (\(QE_T\)) and
the holdout partition (\(QE_H\)). The training value reflects how well
the fitted SOM represents the samples it was trained on, whereas the
holdout value evaluates the same trained map after projecting previously
unseen samples onto it. The holdout partition is therefore not used
during fitting; it is used only to evaluate the trained map, analogous
to a test set in conventional model-training workflows.

Balanced \(QE\), denoted \(QE_B\), is defined as the mean of \(QE_T\)
and \(QE_H\). \(QE_B\) is therefore a composite endpoint that weights
data representation fidelity (train) and  \DIFaddbegin \DIFadd{generalizability }\DIFaddend (holdout)
equally. Note that in the Optuna runs, \(QE_T\) and \(QE_H\) are
optimized jointly as a two-objective vector, with \(QE_B\) only
calculated \emph{post hoc}.

\DIFaddbegin \DIFadd{We quantify local BMU-neighborhood ordering using Mean Tied Rank (MTR)
(Ramos et al. 2018). For each sample \(x_i\), non-winning units are
ranked by graph distance from the first BMU, \(b_i^{(1)}\), with tied
distance shells assigned their average ordinal rank. If the second BMU,
\(b_i^{(2)}\), lies in shell \(S_d\) and \(L_d\) units lie in closer
shells, its tied rank is \(\tau_i=L_d+(|S_d|+1)/2\), and
\(MTR=N^{-1}\sum_i \tau_i\). Lower MTR indicates closer local ordering
between the first and second BMUs.
}

\DIFadd{MTR is used instead of topographic error because one-hop adjacency is
not equivalent across hexagonal, MST, and RNG graphs; topographic error
also depends on lattice properties and map design (Ramos et al. 2018;
Neme and Miramontes 2005; Machón-González and López-García 2010). We
additionally report node utilization, defined as the fraction of nodes
selected as a BMU, and its complement, dead-node fraction. These post
hoc diagnostics are computed on training and holdout splits and balanced
as for \(QE\).
}

\paragraph{\DIFadd{4.1.3 XPySOM calibration
protocol}}\label{xpysom-calibration-protocol}

\DIFadd{We calibrated the FloatSOM batch pathway against XPySOM using the 14
benchmark datasets, 10 shared seeds, full sampling, a \(32\times32\)
map, and one run per dataset--seed condition without top-\(k\)
selection. Both implementations used the same deterministic 70/30
train--holdout split, number of batch epochs, and XPySOM-like untuned
hyperparameters. XPySOM's initialized node weights were transferred into
FloatSOM after reconciling the two implementations' node-coordinate
ordering, so paired runs began from corresponding prototypes. The
primary implementation-concordance comparison used the hexagonal
topology in both implementations; the associated MST and RNG FloatSOM
paths retained the same XPySOM hexagonal reference to quantify the
additional effect of changing the FloatSOM training topology.
}

\DIFadd{We recorded train, holdout, and balanced \(QE\) together with training
time. Within each dataset, paired percentage effects across the 10 seeds
were summarized using the Wilcoxon signed-rank location estimate, its
exact distribution-free 95\% interval, and a two-sided paired Wilcoxon
test; global rows pool the 140 dataset--seed pairs. Supplementary Figs.
S1--S3 and Supplementary Tables S4--S6 report these results. These are
difference tests rather than formal equivalence tests: a non-significant
hexagonal comparison is interpreted as no implementation-associated
\(QE\) difference detected, because no equivalence margin was
prespecified.
}

\DIFaddend \subsubsection{4.2 Speed scaling benchmark
protocol}\label{speed-scaling-benchmark-protocol}

The speed scaling benchmark evaluates runtime and distributed scaling
behavior in different compute and algorithm configurations. Speed
scaling is evaluated with runs across \(G\in\{1,2,4,8\}\) GPUs under a
series of fixed scaling protocols. Runtime summaries are computed from
repeated executions per configuration and reported as both absolute
training time and efficiency relative to the corresponding 1-GPU
comparison. We compute scaling efficiency as
\(E_G=(T_1/T_G)/G \times 100\%\), where \(T_1\) is the 1-GPU runtime and
\(T_G\) is the runtime on \(G\) GPUs. A value of 100\% corresponds to
ideal linear scaling; lower values indicate that communication,
orchestration, or I/O overhead has reduced the realized speedup, while
values above 100\% can occur when additional GPUs also change the memory
or data-staging regime. These scaling runs use the Ray-orchestrated
distributed execution layer built on top of the standard FloatSOM
training path (Moritz et al. 2018). Accordingly, the scaling figures in
Sections 6.1-6.2 and the runtime/scaling comparison reported later
against XPySOM should be interpreted as distributed-execution results
rather than the in-memory Optuna path.

For scaling-efficiency calculations, when a single-GPU baseline was
missing at a given axis value due to timeout, we estimated that baseline
by local linear extrapolation from the last available 1-GPU point on the
same curve. That is, runtime was assumed to scale proportionally with
the axis variable for the extrapolation step (for example, doubling
sample count or doubling dimensionality doubles the estimated 1-GPU
runtime).

Topology speed comparisons include hexagonal, MST, and RNG, with
harmonized workload settings so ratios isolate topology-associated
runtime effects. A dedicated random versus full comparison, run on
\(G\in\{1,2,4\}\) GPUs, is additionally included to isolate
sampling-specific runtime effects independently of the multi-GPU scaling
runs.

\DIFaddbegin \DIFadd{XPySOM was selected as the executable external baseline because it most
closely matches FloatSOM's Python/GPU batch-training regime. We also
attempted aweSOM on the standard speed workload, but all five runs
reached the 1800-s timeout before completing 60\% of \(N\) online
updates. Matching FloatSOM's 10 batch iterations would require \(10N\)
pointwise updates, so aweSOM was not included in the timed comparison.
}

\DIFaddend \paragraph{4.2.1 Scaling benchmark
datasets}\label{scaling-benchmark-datasets}

The speed benchmark uses synthetic random matrices with uniform values
in \([0,1]\). No train-holdout split is used here because the objective
is runtime rather than generalization, so all generated data are used
for training. We use a shared default configuration of \(10^7\) samples,
50 dimensions, a \(32 \times 32\) SOM grid, and 10 training iterations,
and then vary one workload axis at a time: sample count
(\(10^6\)--\(10^9\)), feature dimension (50--5000), or grid side length
(8--64). The exact axis values used for each scaling panel are provided
in Supplementary Table S3.

Additional CPU and RAM resources attached to each GPU are scaled
linearly with GPU count while keeping the software environment and
benchmark procedure consistent across runs. For scaling benchmarks, each
run is timed out if it exceeds 30 minutes of wall time. Per-GPU resource
allocation is fixed at one NVIDIA V100 (32 GB VRAM), 12 CPU cores, 90 GB
system RAM, with one full node comprising 4 GPUs and their associated
CPUs, and with 400 GB of associated local disk storage.

\subsubsection{4.3 Hyperparameter Tuning and
Stability}\label{hyperparameter-tuning-and-stability}

To quantify  \DIFaddbegin \DIFadd{tuning }\DIFaddend benefit, we  \DIFaddbegin \DIFadd{extracted the best Optuna settings for
each sampling-topology combination and formed deployable configurations
from the mean numeric and modal categorical values. These configurations
were derived from winning groups among datasets with at least 1000
samples, matching the sample-size regime expected for most applications.
}\DIFaddend

The stability analysis focuses on four tuned hyperparameters that govern
SOM training dynamics. These settings determine how the map is
initialized and how updates evolve over training: the initial radius
sets the early neighborhood scale around each BMU, the initialization
method sets the starting node weights, the radius decay type controls
the shift from broad global organization toward local refinement, and
the momentum-use parameter determines whether each update retains part
of the previous update direction.

\paragraph{4.3.1 Tuned Configuration versus Untuned Reference
Analysis}\label{tuned-configuration-versus-untuned-reference-analysis}

 \DIFaddbegin \DIFadd{Fixed configurations were rerun under matched dataset, seed, topology,
and split keys. The final topology diagnostic comprised all 14 benchmark
datasets, 20 shared random seeds, full sampling, and each of the
hexagonal, MST, and RNG topologies, giving 280 dataset--seed units per
paired contrast. Preprocessing and the deterministic 70/30
train--holdout split followed Section 4.1.1. The untuned profile used
the same XPySOM-like settings for all topologies: initial radius 5,
random initialization, exponential radius decay, initial momentum 0.5,
and momentum disabled. The tuned profile used the fixed
topology-specific full-sampling configurations derived above: initial
radii of 1.0266, 1.4639, and 1.4148 and initial momenta of 0.6069,
0.6146, and 0.6127 for hexagonal, MST, and RNG, respectively, with
random initialization, asymptotic radius decay, and momentum enabled.
Deployment analyses compare these tuned FloatSOM configurations with
untuned hexagonal XPySOM, while topology diagnostics compare tuned and
untuned FloatSOM within and between topologies. These reruns estimate
deployable fixed-configuration performance; the top-\(k\) Optuna
summaries instead estimate attainable performance within the search
budget.
}\DIFaddend

 \DIFaddbegin \DIFadd{For Supplementary Fig. S5, tuned and untuned values were paired within
dataset, seed, topology, and split. Supplementary Table S12 pools the
280 matched dataset--seed differences for each prespecified
profile--topology--metric contrast and reports the mean paired effect,
two-sided one-sample paired \(t\)-test confidence interval, Cohen's
\(d_z\), raw p-value, and win/loss/tie counts. These raw p-values are
descriptive diagnostics rather than a multiplicity-adjusted family;
effect sizes and confidence intervals are the primary summaries.
Supplementary Table S13 instead reports the corresponding mean observed
metric within each profile, dataset, and topology across the 20 seeds.
Dataset-level }\DIFaddend forest-plot  \DIFaddbegin \DIFadd{results use the convention described above,
and overall rows pool all matched pairs with }\DIFaddend the same paired \(t\)-test
and confidence-interval  \DIFaddbegin \DIFadd{calculation}\DIFaddend .

\paragraph{4.3.2 Hyperparameter stability and dataset-type
stratification}\label{hyperparameter-stability-and-dataset-type-stratification}

Hyperparameter stability is important because it determines whether a
method can be deployed reliably without bespoke retuning. We therefore
extracted the top-ranked tuned Optuna trial separately for each seed and
compared the recovered hyperparameter values across seeds. We define
hyperparameter stability as the variation in these recovered settings.
For numeric parameters, stability is measured by the relative difference
between runs, \(|a-b| / \max(|a|,|b|,\varepsilon)\), where \(a\) and
\(b\) are the values of a given parameter for the two compared seeds and
\(\varepsilon=10^{-12}\). These relative differences are then averaged
across numeric parameters, with lower values indicating higher
stability. For categorical parameters, stability is defined as the
mismatch rate across the same seed pairs. We report both mean
per-parameter stability and an equal-weight overall stability summary
across datasets.

\DIFaddbegin \paragraph{\DIFadd{4.3.3 Matched initial-radius sensitivity
analysis}}\label{matched-initial-radius-sensitivity-analysis}

\DIFadd{To determine whether topology differences could be explained by
topology-specific treatment of neighborhood radius, we performed a
controlled initial-radius sweep. Hexagonal, MST, and RNG maps were
evaluated at \(r\in\{0.5,0.75,1.0266,1.5,2,3,5\}\) using the same 20
random seeds on each of the 14 benchmark datasets. The value
\(r=1.0266\) was the selected hexagonal full-sampling radius and is
denoted \texttt{Hex\ optimal} in Fig. 9. Apart from topology and the
deliberately varied initial radius, the training configuration was
fixed: full sampling, random initialization, asymptotic radius decay,
momentum enabled with initial momentum 0.6069, and XPySOM-compatible
normalization. Runs were therefore matched by dataset and seed, ensuring
that topology families received the same radius values rather than
different search ranges or topology-specific radius schedules.
}

\DIFadd{We evaluated balanced \(QE\), balanced MTR, and balanced node
utilization. Because absolute \(QE\) scales differ substantially among
datasets, each dataset--seed \(QE_B\) value was divided by the matched
hexagonal \(QE_B\) at \(r=1.0266\). Ratios were averaged on the log
scale across seeds within each dataset and then across datasets with
equal dataset weight; exponentiating this mean yields the geometric-mean
ratio shown in Fig. 9A. Balanced MTR and balanced node utilization are
retained in their observed units in Fig. 9B and Fig. 9C, respectively.
Along each topology's response curve, direction-aware percentage changes
from \(r=1.0266\) were calculated for all 280 matched dataset--seed
units and tested against zero using two-sided one-sample \(t\)-tests;
Benjamini--Hochberg adjustment was applied across the 54 non-anchor
radius--topology--metric tests. These tests assess within-topology
radius response.
}

\DIFadd{We additionally conducted post hoc cross-topology \(QE_B\) contrasts at
each radius. For every topology pair, dataset, seed, and radius, we
calculated the log ratio of the comparator \(QE_B\) to the reference
\(QE_B\), averaged log ratios across the 20 seeds within each dataset,
and applied a two-sided one-sample \(t\)-test to the 14 equally weighted
dataset effects. Exponentiated means and confidence limits are
geometric-mean \(QE_B\) ratios, with values below 1 favouring the
comparator. Benjamini--Hochberg correction was applied across all 21
topology-pair--radius contrasts. Supplementary Table S15 reports the
estimates, intervals, raw p-values, and adjusted q-values.
}

\paragraph{\DIFadd{4.3.4 Execution-path concordance
analysis}}\label{execution-path-concordance-analysis}

\DIFadd{To determine whether the quality diagnostics depended on execution
pathway, we compared local CuPy and Ray streaming under matched inputs
and fixed tuned configurations. The complete protocol is provided in
Supplementary Methods S1, and the paired results are reported in
Supplementary Table S14. Because no equivalence margin was prespecified,
non-significant results are interpreted as no difference detected rather
than proof of formal equivalence.
}

\DIFaddend \subsubsection{4.4 Statistical analysis}\label{statistical-analysis}

All Optuna comparisons use matched pairs within dataset, seed, and split
units to control for substantial between-run heterogeneity. Within each
matched unit, trials were ranked by the target metric, the top five were
retained, and each condition was summarized by the median of those
retained trials, yielding a top-\(k\) summary with \(k=5\) intended to
estimate near-optimal attainable performance under a fixed number of
tuning trials, where each trial is one candidate hyperparameter
configuration evaluated by Optuna. Paired effects were then computed as
simple condition differences, with negative values favoring the first
condition for lower-is-better metrics. Dataset-level and global
summaries are shown as forest plots with 95\% confidence intervals from
two-sided paired one-sample \(t\)-tests \DIFaddbegin \DIFadd{. Supplementary Fig. S4 repeats
the topology comparisons at \(k\in\{1,3,5,10\}\), always using the
median of the retained trials within each matched unit, to assess
sensitivity to the number of near-optimal trials retained}\DIFaddend .

\DIFaddbegin \DIFadd{We report Benjamini-Hochberg q-values alongside raw p-values for related
dataset-level tests, and use q-values for figure markers and
significance counts. Each topology contrast in Figs. 6-7 forms a
separate family of 42 tests (14 datasets across \(QE_B\), \(QE_H\), and
\(QE_T\)). Pooled overall tests are reported separately with raw
p-values.
}

\DIFaddend \subsection{5. Results}\label{results}

\subsubsection{5.1 XPySOM calibration } \DIFaddbegin \label{xpysom-calibration}
\DIFaddend

Under matched-configuration XPySOM versus FloatSOM calibration on
hexagonal \(QE\) (Fig. S3),  \DIFaddbegin \DIFadd{no implementation-associated }\DIFaddend \(QE\)
 \DIFaddbegin \DIFadd{difference was detected}\DIFaddend . Accordingly, we  \DIFaddbegin \DIFadd{use }\DIFaddend hexagonal FloatSOM batch as
 \DIFaddbegin \DIFadd{the operational XPySOM reference }\DIFaddend in the benchmarks that follow\DIFaddbegin \DIFadd{, while
not claiming formal statistical equivalence}\DIFaddend . The runtime differences are
attributable to FloatSOM's JIT kernels, which incur a small startup
cost. This overhead is progressively amortized as workload size
increases, after which FloatSOM runs faster than XPySOM on larger
datasets.

\subsubsection{5.2 Comparison of Different Sampling
Methods}\label{comparison-of-different-sampling-methods}

We first report a focused HDSSSOM pilot as an elimination comparison
rather than as part of the broader sampling benchmark. This pilot used a
smaller, more limited configuration than the later full versus random
analysis and is summarized in Supplementary Table S2.

Under that pilot configuration, HDSSSOM was materially worse than full
sampling. In the hexagonal view shown in Fig. 3, full outperformed
HDSSSOM in all 50 paired comparisons (10 datasets \(\times\) 5 seeds; 50
wins, 0 losses, 0 ties), with dataset-level median balanced-\(QE\)
improvements ranging from 2.5\% to 209.7\% and a global median
improvement of 38.7\%. \DIFaddbegin \DIFadd{These HDSSSOM results should therefore be
interpreted as a pilot screen under the stated configuration rather than
as a comprehensive evaluation of all possible HDSSSOM schedules. }\DIFaddend Given
these large differences, we focused on the full and random sampling
methodologies for the remainder of the manuscript.

\pandocbounded{\includesvg[keepaspectratio]{assets_manual/figures/fig_3.svg}}
\emph{Figure 3. HDSSSOM screening pilot on \(QE_B\) (hexagonal
topology): full vs HDSSSOM, using the smaller pilot configuration
summarized in Supplementary Table S2 (10 datasets, 5 seeds, with all
other pilot settings held fixed). Panels report dataset-matched paired
top-\(k\) within-unit medians plus dataset-level paired-effect summaries
(Section  \DIFaddbegin \DIFadd{4.4}\DIFaddend ). Forest whiskers denote 95\% paired \(t\)-test confidence
intervals around the mean paired effect.}

The full versus random analysis in Fig. 4 was generated from matched
hexagonal Optuna runs in which the sampling selector was switched from
full to random. Effects were computed within matched seed-specific units
and then summarized across seeds. Fig. 4A-C summarize the matched full
versus random paired effects for balanced, holdout, and train \(QE\),
respectively, showing that full sampling generally provides equal or
better \(QE\) than random sampling. Notably, this full versus random
separation is much smaller than the full versus HDSSSOM pilot effect
shown in Fig. 3; in most dataset-level comparisons, the full versus
HDSSSOM improvement is at least twice as large - as evidenced by the
difference in axis scales between Fig. 3 and Fig. 4.

We find the effectiveness of random sampling relative to full sampling
is scale-dependent. Above \(10{,}000\) samples, paired \(QE\)
differences are not meaningfully detected, whereas in smaller datasets
the random arm shows higher variability and less stable outcomes (Fig.
4D-F). This is consistent with reduced per-iteration sample support
under random subsampling (Fig. 4D-F). In this benchmark, the
\(>10{,}000\) regime is therefore a useful practical proxy for more
stable random sampling behavior.

\pandocbounded{\includesvg[keepaspectratio]{assets_manual/figures/fig_4.svg}}
\emph{Figure 4. Full versus random sampling under the hexagonal
topology. Panels A-C report paired full-versus-random effects for
\(QE_B\), \(QE_H\), and \(QE_T\). Panels D-F show the corresponding
dataset-size regressions for the same three metrics, with numbered
datasets keyed in Supplementary Table S1. \DIFaddbegin \DIFadd{Figure significance markers
use Benjamini-Hochberg adjusted q-values. }\DIFaddend In the larger-dataset regime
observed here (\textgreater10,000 samples), little paired \(QE\)
separation is detected; at smaller dataset scales, random is more
variable and less stable.}

Nevertheless, full sampling remains the best sampling strategy for
optimal \(QE\) results. Accordingly, all remaining analyses reported
below rely on full sampling unless specified.

\subsubsection{5.3 Topology Results}\label{topology-results}

Topology comparisons  \DIFaddbegin \DIFadd{use }\DIFaddend \(QE\)  \DIFaddbegin \DIFadd{as the optimized endpoint and hexagonal
batch SOM as the }\DIFaddend regular-topology baseline; Fig. 5  \DIFaddbegin \DIFadd{illustrates the three
neighborhood structures. \texttt{initial\_radius} was optimized over the
same interval for every topology. The selected full-sampling values were
1.03 for hexagonal, 1.46 for }\DIFaddend MST, and  \DIFaddbegin \DIFadd{1.41 for RNG; under random
sampling they were 1.17, 1.82, and 1.77. Thus, the graph-topology gains
are not comparisons against a broader default-radius hexagonal map,
although radius was not isolated from the other tuned parameters.
}\DIFaddend

\pandocbounded{\includesvg[keepaspectratio]{assets_manual/figures/fig_5.svg}}
\emph{Figure 5. Representative  node and connection overlays for  \DIFaddbegin \DIFadd{matched
100-node }\DIFaddend hexagonal, MST, and RNG  \DIFaddbegin \DIFadd{SOMs. A--C: circles in native 2D. D--F:
KDD Cup 99, trained }\DIFaddend on \DIFaddbegin \DIFadd{standardized 41-dimensional observations and
displayed using one shared 2D PCA projection fitted to the observations
and applied to all three sets of prototypes. Training used seed 42, full
sampling, random initialization, 50 iterations, and the untuned
XPySOM-like settings used at this point in the manuscript: initial
learning rate 0.5, initial radius 5 with exponential decay, momentum
disabled, and XPySOM-compatible normalization. Grey observation clouds
show }\DIFaddend a \DIFaddbegin \DIFadd{deterministic maximum of }\DIFaddend 30,000  \DIFaddbegin \DIFadd{points; black lines are topology
connections and red markers are SOM nodes}\DIFaddend . \DIFaddbegin \DIFadd{PCA can distort graph
geometry in the original 41-dimensional space.}\DIFaddend }

\DIFaddbegin \DIFadd{In the KDD Cup 99 projection (Fig. 5D--F), the hexagonal topology
contains several apparently nonlocal connections spanning separated
regions of the projected prototype distribution, whereas the MST and RNG
connections more closely follow its local geometry. This pattern
illustrates how predetermined lattice neighbors can couple the updates
of prototypes that are not locally adjacent in the displayed data
structure, restricting how independently the nodes can redistribute.
Because the display is a two-dimensional PCA projection of the
41-dimensional training space, this visual pattern is illustrative
rather than direct evidence of nonlocality in the original space.
}

\DIFaddend \paragraph{5.3.1 MST}\label{mst}

We report dataset-wise paired improvement summaries (hexagonal over MST)
 \DIFaddbegin \DIFadd{using the statistical procedure in Section 4.4 }\DIFaddend in Fig. 6.

Overall, MST outperforms matched hexagonal on  \DIFaddbegin \DIFadd{Balanced }\DIFaddend QE (Fig. 6A),
indicating a net advantage across train and holdout performance. This
aggregate gain is driven more clearly by  \DIFaddbegin \DIFadd{Train }\DIFaddend QE (Fig. 6C), while
 \DIFaddbegin \DIFadd{Holdout }\DIFaddend QE is more mixed across datasets (Fig. 6B) and shows no clear
overall holdout advantage. The overall paired t-test p-values are
 \DIFaddbegin \DIFadd{Balanced }\DIFaddend QE (p=1.12e-05) \DIFaddbegin \DIFadd{; Holdout }\DIFaddend QE (p=0.15) \DIFaddbegin \DIFadd{; and Train }\DIFaddend QE (p=0.0064).
\DIFaddbegin \DIFadd{Supplementary Table S7 lists the corresponding per-dataset and overall
hexagonal-comparison effect estimates, 95\% confidence intervals, raw
p-values, and Benjamini-Hochberg q-values for MST and RNG.
}\DIFaddend

\pandocbounded{\includesvg[keepaspectratio]{assets_manual/figures/fig_6.svg}}
\emph{Figure 6. Hexagonal versus MST topology on \(QE\) metrics under
full sampling only. Panels A-C report paired full sampling only \(QE\)
effects for \(QE_B\), \(QE_H\), and \(QE_T\) across the available full
sampling datasets. Forest whiskers denote 95\% paired \(t\)-test
confidence intervals around the mean paired effect. \DIFaddbegin \DIFadd{Supplementary Table
S7 reports the per-dataset effect estimates, 95\% confidence intervals,
raw p-values, and dataset-level Benjamini-Hochberg adjusted q-values.}\DIFaddend }

\paragraph{5.3.2 RNG}\label{rng}

To evaluate RNG topology performance, we reuse the paired reporting
logic on hexagonal versus RNG, again centered on \(QE_B\) with \(QE_H\)
and \(QE_T\) in Fig. 7.

RNG has lower QE than matched hexagonal on the reported QE  \DIFaddbegin \DIFadd{endpoints
}\DIFaddend (Fig.  \DIFaddbegin \DIFadd{7A-7C}\DIFaddend ), with overall paired t-test p-values of  \DIFaddbegin \DIFadd{Balanced }\DIFaddend QE
(p=7.4e-10) \DIFaddbegin \DIFadd{; Holdout }\DIFaddend QE (p=0.0232) \DIFaddbegin \DIFadd{; and Train }\DIFaddend QE (p=4.69e-06).
\DIFaddbegin \DIFadd{Supplementary Table S7 lists the corresponding per-dataset and overall
hexagonal-comparison effect estimates, 95\% confidence intervals, raw
p-values, and Benjamini-Hochberg q-values for MST and RNG.
}\DIFaddend

The main trend in Fig. 7 is that RNG improves on hexagonal most clearly
in balanced QE and especially in train QE, with the separation most
apparent in the real and larger datasets where the added flexibility of
the graph neighborhood appears more useful than the fixed regular
lattice.

\pandocbounded{\includesvg[keepaspectratio]{assets_manual/figures/fig_7.svg}}
\emph{Figure 7. Hexagonal versus RNG topology on \(QE\) metrics under
full sampling only. Panels A-C report paired full sampling only \(QE\)
effects for \(QE_B\), \(QE_H\), and \(QE_T\) across the available full
sampling datasets. Forest whiskers denote 95\% paired \(t\)-test
confidence intervals around the mean paired effect. \DIFaddbegin \DIFadd{Supplementary Table
S7 reports the per-dataset effect estimates, 95\% confidence intervals,
raw p-values, and dataset-level Benjamini-Hochberg adjusted q-values.}\DIFaddend }

\DIFaddbegin \paragraph{\DIFadd{5.3.3 Direct MST--RNG
comparison}}\label{direct-mstrng-comparison}

\DIFadd{Fig. 8 directly compares MST and RNG under full sampling. The two graph
topologies are close on \(QE\), and neither dominates across all three
endpoints.
}

\pandocbounded{\includesvg[keepaspectratio]{assets_manual/figures/fig_8.svg}}
\DIFadd{\emph{Figure 8. MST versus RNG topology on \(QE\) metrics under full
sampling only. Panels A-C report paired full sampling only \(QE\)
effects for \(QE_B\), \(QE_H\), and \(QE_T\) across the available full
sampling datasets. Forest whiskers denote 95\% paired \(t\)-test
confidence intervals around the mean paired effect.}
}

\paragraph{\DIFadd{5.3.4 Initial-radius
sensitivity}}\label{initial-radius-sensitivity}

\DIFadd{Because initial radius was selected during tuning and differed across
topology families, we isolated its effect in a fixed-configuration
sweep. All non-radius hyperparameters were held at the tuned hexagonal
full-sampling configuration, and each of seven radii was evaluated for
hexagonal, MST, and RNG maps using the same 20 seeds on all 14 datasets.
Fig. 9A reports \(QE_B\) after normalizing every dataset--seed condition
to its matched hexagonal value at the selected radius (\(r=1.027\)). We
summarize these ratios by averaging log ratios across seeds within each
dataset, weighting the 14 datasets equally, and exponentiating the
cross-dataset mean. The resulting geometric-mean ratio is dimensionless:
1 denotes matched hexagonal performance at the selected radius, values
below 1 indicate lower \(QE_B\), and values above 1 indicate higher
\(QE_B\).
}

\DIFadd{The observed \(QE_B\) minimum occurs at \(r=0.75\) for hexagonal and at
the larger radius \(r=1.5\) for both MST and RNG (Fig. 9A). For MST and
RNG, however, the response is flat between \(r=1.027\) and \(r=1.5\):
the within-topology differences from the anchor are not significant
after adjustment (MST q=0.35; RNG q=0.21). Thus, \(r=1.5\) is the lowest
observed point in this grid rather than evidence of a sharply identified
optimum. The post hoc matched topology contrasts support the separation
in the useful radius region: at \(r=1.027\), MST and RNG have 1.51\% and
1.85\% lower geometric-mean \(QE_B\) than hexagonal (q=0.036 and
q=0.035), increasing to 3.30\% and 3.67\% at \(r=1.5\) (q=0.016 and
q=0.014; Supplementary Table S15). No RNG--MST contrast is significant
at any tested radius. The graph-topology advantage over hexagonal
therefore persists under identical radius treatment, while the two graph
topologies remain close competitors on \(QE_B\).
}

\DIFadd{The sweep also reveals a pronounced topology-specific
quantization--ordering trade-off. In Fig. 9B, increasing radius sharply
lowers hexagonal MTR but simultaneously and progressively worsens its
\(QE_B\) in Fig. 9A. Around the useful \(r=1.027\)--1.5 region, MST and
RNG do not show the same penalty: their \(QE_B\) remains essentially
flat, while RNG MTR improves and MST MTR remains stable. At still larger
radii, MST and RNG eventually also exchange higher \(QE_B\) for lower
MTR, but this trade-off is delayed and substantially weaker than for
hexagonal maps. The adaptive graph topologies therefore preserve a more
favourable joint \(QE\)/MTR operating region rather than merely shifting
the same hexagonal response to a different radius.
}

\DIFadd{Fig. 9C provides the complementary map-capacity check. Hexagonal node
utilization falls from 0.874 at the selected radius to 0.859 at
\(r=1.5\) and approximately 0.846 at \(r\geq2\). By contrast, MST and
RNG remain near 0.89 throughout the \(r=1.027\)--2 region in which their
\(QE_B\) response is also comparatively stable. These utilization
results do not support poorer use of the available SOM nodes as an
explanation for the graph topologies' lower \(QE_B\).
}

\DIFadd{RNG also has lower observed MTR than MST throughout the sweep (Fig. 9B);
for example, at the hexagonal selected radius, mean balanced MTR is 5.10
for RNG and 8.06 for MST. This agrees with the matched topology
diagnostics, in which RNG has the stronger MTR result. The significance
symbols in Fig. 9 should not be read as MST--RNG comparisons: they
compare each non-anchor radius with \(r=1.027\) within the same
topology. The panel therefore separates two findings: RNG maintains
better local ordering than MST in absolute MTR, while changing radius
away from the anchor produces topology-specific within-curve changes.
}

\pandocbounded{\includesvg[keepaspectratio]{assets_manual/figures/fig_9.svg}}
\DIFadd{\emph{Figure 9. Initial-radius sensitivity under matched fixed
configurations. A: dataset-balanced geometric-mean \(QE_B\) ratio
relative to the matched hexagonal configuration at its selected radius
(\(r=1.027\)); lower values are better and the dashed horizontal line
marks a ratio of 1. B: observed balanced Mean Tied Rank (MTR; lower is
better). C: observed balanced node utilization (higher is better). All
panels summarize 14 datasets with 20 matched seeds per topology and
radius. The dotted vertical line and \texttt{Hex\ optimal} tick mark
\(r=1.027\). Significance markers use Benjamini--Hochberg adjusted tests
comparing each non-anchor radius with \(r=1.027\) within topology across
the 54 radius--topology--metric tests. Confidence intervals are omitted
from the plotted panels for readability; Supplementary Table S15 reports
the post hoc matched cross-topology \(QE_B\) estimates, intervals, and
adjusted tests at every radius.}
}

\DIFadd{The fixed-configuration diagnostics provide the corresponding MTR and
node-use comparison (Table
\ref{tab:matched_topology_diagnostics_summary}). MST and RNG both
improved balanced \(QE\) relative to hexagonal maps with and without
tuning, but only RNG clearly lowered untuned balanced MTR. Under tuning,
no balanced-\(QE\) difference was detected between MST and RNG, while
RNG lowered balanced MTR by 2.96 tied-rank positions. RNG therefore gave
the strongest joint \(QE\)/MTR result, while MST remained a close
competitor on \(QE\).
}

\begin{table}[t]
\centering
\caption{\DIFaddFL{Matched topology diagnostics for balanced quantization error and Mean Tied Rank. Positive paired effects favor the second topology in each contrast after applying metric directionality; lower raw values are better for both $QE_B$ and $MTR_B$. Full confidence intervals, raw paired-test p-values, split-specific metrics, node utilization, and dead-node fraction are reported in Supplementary Table S12.}}
\label{tab:matched_topology_diagnostics_summary}
\begin{tabular}{lrrrrrr}
\toprule
& \multicolumn{3}{c}{\DIFaddFL{$QE_B$}} & \multicolumn{3}{c}{\DIFaddFL{$MTR_B$}} \\
\cmidrule\DIFaddFL{(lr)}{\DIFaddFL{2-4}}\cmidrule\DIFaddFL{(lr)}{\DIFaddFL{5-7}}
\DIFaddFL{Profile }& \DIFaddFL{MST vs hex }& \DIFaddFL{RNG vs hex }& \DIFaddFL{RNG vs MST }& \DIFaddFL{MST vs hex }& \DIFaddFL{RNG vs hex }& \DIFaddFL{RNG vs MST }\\
\midrule
\DIFaddFL{Untuned }& \DIFaddFL{+0.269 }& \DIFaddFL{+0.249 }& \DIFaddFL{-0.020 }& \DIFaddFL{-0.20 }& \DIFaddFL{+2.44 }& \DIFaddFL{+2.64 }\\
\DIFaddFL{Tuned }& \DIFaddFL{+0.064 }& \DIFaddFL{+0.052 }& \DIFaddFL{-0.012 }& \DIFaddFL{+22.71 }& \DIFaddFL{+25.67 }& \DIFaddFL{+2.96 }\\
\bottomrule
\end{tabular}
\end{table}

\DIFaddend \subsubsection{5.4 Hyperparameter Tuning and
Stability}\label{hyperparameter-tuning-and-stability-1}

\paragraph{5.4.1 Performance Gains from Hyperparameter
Tuning}\label{performance-gains-from-hyperparameter-tuning}

 \DIFaddbegin \DIFadd{Having shown that the graph-topology effects persist when radius is
varied under a common tuned non-radius configuration, we next quantify
the broader effect of complete topology-specific tuning relative to the
untuned reference.
}

\DIFadd{Fig. 10 compares the fixed tuned and untuned configurations using pooled
}\DIFaddend topology-level \(QE\)  \DIFaddbegin \DIFadd{summaries; topology-stratified results are
provided in Supplementary Fig. S5}\DIFaddend .

 \DIFaddbegin \pandocbounded{\includesvg[keepaspectratio]{assets_manual/figures/fig_10.svg}}
\DIFadd{\emph{Figure 10. Tuned configuration versus untuned reference \(QE\)
comparison across \(QE_B\), \(QE_H\), and \(QE_T\), pooled across all
topology runs under the matched pairing keys. Positive values indicate
the tuned configuration outperforms the untuned reference; the global
overall row pools all matched tuned configuration/untuned reference
pairs across datasets. Forest whiskers denote 95\% paired \(t\)-test
confidence intervals around the mean paired effect.}
}\DIFaddend

 \DIFaddbegin \DIFadd{The fixed tuned configurations had lower values }\DIFaddend for all three  \DIFaddbegin \DIFadd{pooled
\(QE\) metrics than the untuned references, but also had higher balanced
MTR: by 24.34 tied-rank positions for hexagonal maps compared with 1.43
for MST and 1.11 for RNG. This profile-associated local-ordering
trade-off was therefore much stronger for the fixed lattice and was not
accompanied by lower node utilization. The \(QE\) improvement was
observed in every topology family}\DIFaddend .

\paragraph{5.4.2 Hyperparameter Stability Across Topology and
Sampling}\label{hyperparameter-stability-across-topology-and-sampling}

To assess whether the derived hyperparameters are robust across runs and
datasets, we compare the within-topology seed-to-seed tuned-parameter
drift (Section 4.3.2). Overall, MST and RNG  \DIFaddbegin \DIFadd{exhibit less variability
than hexagonal maps under both full and random sampling (Fig. 11A-B}\DIFaddend ).
Across these panels, full sampling is generally the more stable setting.
Mirroring the sample-size-dependent \(QE\) performance, Fig.  \DIFaddbegin \DIFadd{11C shows
that random-sampling }\DIFaddend hyperparameter stability also improves with
increasing sample size.

 \DIFaddbegin \pandocbounded{\includesvg[keepaspectratio]{assets_manual/figures/fig_11.svg}}
\DIFaddend \emph{Figure  \DIFaddbegin \DIFadd{11. }\DIFaddend Hyperparameter stability by sampling mode. A: selected
parameter stability under full sampling for hexagonal, MST, and RNG
topologies (lower stability score is better). B: selected parameter
stability under random sampling for the same topologies. C: dataset-size
stability regression under random sampling, using the selected parameter
stability score against sample size (log10) across the included topology
families.}

\DIFaddbegin \subsubsection{\DIFadd{5.5 Execution-Path
Diagnostic}}\label{execution-path-diagnostic}

\DIFadd{The quality analyses above used the in-memory pathway. Repeating the
tuned diagnostics through Ray streaming with matched data and
configurations produced small differences, none of which reached
significance before or after Benjamini-Hochberg correction
(Supplementary Table S14). The remaining differences may reflect
floating-point accumulation order. The following speed benchmarks use
the Ray pathway.
}

\DIFaddend \subsection{6. Speed Scaling}\label{speed-scaling}

We next examine three aspects of FloatSOM runtime performance: the cost
of full relative to random sampling, the scaling behavior obtained with
additional GPUs, and the runtime implications of MST and RNG relative to
a hexagonal topology.

\subsubsection{6.1 Random versus Full Sampling
Runtime}\label{random-versus-full-sampling-runtime}

We report sample scaling runtime comparisons for full versus random
sampling using the synthetic datasets as outlined in Section 4.2.1. We
stratify by hexagonal, MST, and RNG, for 1, 2, and 4 GPUs (Fig.  \DIFaddbegin \DIFadd{12}\DIFaddend ).

\begin{figure}
\centering
 \DIFaddbeginFL \pandocbounded{\includesvg[keepaspectratio]{assets_manual/figures/fig_12.svg}}
\DIFaddendFL \caption{Figure  \DIFaddbeginFL \DIFaddFL{12}\DIFaddendFL }
\end{figure}

\emph{Figure  \DIFaddbegin \DIFadd{12. }\DIFaddend Sample scaling runtime comparison of full versus random
sampling across \(G\in\{1,2,4\}\) GPUs (A,B,C). Curves report mean
wall-clock training time (s) under harmonized settings; error bars
denote \(\pm 1\) standard deviation across \(n=3\) repeated runs per
configuration. Color encodes topology (hexagonal, MST, RNG), and line
style encodes sampling mode (full vs.~random). Shaded x-axis regions
indicate sample-size ranges that could not be run in that panel relative
to the shared axis maximum due to timeouts. Lower values indicate faster
execution.}

Across topologies, random sampling is faster than full sampling across
all 1-, 2-, and 4-GPU comparisons, with similar proportional reductions
at a given dataset size. With more GPUs, larger datasets can also be
processed before timing out, although the runtime advantage of random
over full narrows at the largest sample counts.

For the 1-GPU random sampling runs, the last successful 1-GPU run was at
100M samples, the same number of samples as run on full samples, despite
random being much faster. We reason that the failed subsequent runs on
larger samples were due to the transition from RAM operation to
disk-backed operation. Even in random sampling mode, the implementation
still stages the full dataset initially and loads full chunks before the
subsampler selects the training samples. Therefore, the relevant cost is
no longer only the reduced number of selected samples. We therefore
treat the 1-GPU random failure at 500M samples as a disk-mode systems
limitation rather than as evidence against the general random versus
full runtime ordering.

\subsubsection{6.2 Multi-GPU topology scaling and Larger-Than-Memory
context}\label{multi-gpu-topology-scaling-and-larger-than-memory-context}

To examine parallel scaling, we benchmarked FloatSOM under full sampling
across \(G\in\{1,2,4,8\}\) GPUs using workload-scaling benchmarks that
extend from standard in-memory settings to larger workloads that exceed
available memory.  \DIFaddbegin \DIFadd{Section 6.2 primarily concerns sample scaling. In that
setting, the topologies exhibit similar scaling characteristics: as
sample count increases, the GPU-count response and efficiency curves
have similar qualitative shapes for RNG (Fig. 13B,E) and for the
corresponding hexagonal and MST outputs (Fig. S6). Conversely, when the
number of SOM nodes is increased (grid-size scaling), the topologies
differ substantially. That grid-size regime is analyzed in Section 6.3,
where MST and RNG take 8.19x and 27.07x the hexagonal runtime,
respectively, at the largest tested grid size}\DIFaddend .

\begin{figure}
\centering
 \DIFaddbeginFL \pandocbounded{\includesvg[keepaspectratio]{assets_manual/figures/fig_13.svg}}
\DIFaddendFL \caption{Figure  \DIFaddbeginFL \DIFaddFL{13}\DIFaddendFL }
\end{figure}

\emph{Figure  \DIFaddbegin \DIFadd{13. }\DIFaddend Multi-GPU scaling across \(G\in\{1,2,4,8\}\) GPUs.
Panels A-C show runtime (s) for dimension, sample, and grid size scaling
workloads, respectively. Panels D-F show scaling efficiency for the same
workloads, computed from the single-GPU baseline and the corresponding
\(G\)-GPU runtime. Runtime error bars denote \(\pm 1\) standard
deviation across \(n=3\) repeated runs per configuration; the 100\%
efficiency reference line indicates ideal linear scaling.}

\paragraph{6.2.1 GPU Scaling and Larger-Than-Memory Runtime
Performance}\label{gpu-scaling-and-larger-than-memory-runtime-performance}

Fig.  \DIFaddbegin \DIFadd{13 }\DIFaddend shows that increasing GPU count improves performance in the
sample scaling regime by increasing parallel throughput and delaying the
transition to disk-backed execution. In the sample scaling benchmark,
the 500,000,000 sample dataset requires disk backing under the 2-GPU
configuration, whereas the 8-GPU configuration remains in RAM mode until
the 1,000,000,000 sample dataset; when staging is still required, the
disk-to-GPU path is distributed across more nodes. The 8-GPU RNG
configuration processes 1,000,000,000 samples in 369.41 s (6.16 min).
Note that this speed includes the time required to transfer data from
shared storage to node-local shards, alongside the overhead associated
with operating across multiple HPC nodes.

 \DIFaddbegin \DIFadd{Grid-size scaling is }\DIFaddend the main exception: at the largest tested grid size
(64), runtime  \DIFaddbegin \DIFadd{decreases }\DIFaddend from 934.01 s (15.57 min) on 1 GPU to  880.83 s
(14.68 min) on 8 GPUs, a 5.69\% reduction \DIFaddbegin \DIFadd{. Once }\DIFaddend map-size  \DIFaddbegin \DIFadd{and
}\DIFaddend topology-refresh costs dominate, additional GPUs  \DIFaddbegin \DIFadd{therefore contribute
little further }\DIFaddend speedup.

\paragraph{6.2.2 GPU Scaling Efficiency}\label{gpu-scaling-efficiency}

We next consider GPU efficiency under strong scaling, relative to ideal
linear scaling. At smaller dataset sizes, efficiency is lower. As
workload size increases, efficiency rises sharply and in some regions
exceeds 100\%. When a direct 1-GPU baseline was unavailable at a given
axis value, the efficiency denominator was constructed by local linear
extrapolation from the last available 1-GPU point on that curve (Section
4.2), so some values should be interpreted with care if the underlying
1-GPU runtime is nonlinear over that range. \DIFaddbegin \DIFadd{Efficiencies above 100\%
should also be interpreted as a combined consequence of parallelism and
a changed memory/data-staging regime, not as evidence of superlinear
compute scaling. Together, Fig. 13 and }\DIFaddend Fig.  \DIFaddbegin \DIFadd{S6 support the same
qualitative sample-scaling efficiency trend across topologies. This does
not extend to grid-size scaling; Section 6.3 shows that when node count
increases, topology-dependent absolute runtimes diverge.
}\DIFaddend

\subsubsection{6.3 Topology Runtime
Comparisons}\label{topology-runtime-comparisons}

With that systems context in place, we next compare topology runtimes
across hexagonal, MST, and RNG configurations on 8 GPUs (Fig.  \DIFaddbegin \DIFadd{14}\DIFaddend ).

\begin{figure}
\centering
 \DIFaddbeginFL \pandocbounded{\includesvg[keepaspectratio]{assets_manual/figures/fig_14.svg}}
\DIFaddendFL \caption{Figure  \DIFaddbeginFL \DIFaddFL{14}\DIFaddendFL }
\end{figure}

\emph{Figure  \DIFaddbegin \DIFadd{14. }\DIFaddend Topology runtime comparison at fixed \(G=8\) GPUs under
full sampling. Panels A-C report mean wall-clock runtime (s) for
dimension, sample, and grid size scaling workloads, respectively, with
topology traces for hexagonal, MST, and RNG. Error bars denote \(\pm 1\)
standard deviation across \(n=3\) repeated runs per configuration. The
largest-axis 8-GPU topology runtime summaries are listed in
Supplementary Table S11.}

In Fig.  \DIFaddbegin \DIFadd{14A-B}\DIFaddend , the topologies scale similarly as input complexity and
data volume increase: even at the largest tested axis values, the
maximum pairwise runtime spread remains modest at dimension scaling
(4.70\% at 5,000 dimensions) and sample scaling (3.26\% at 1,000,000,000
samples).

However, when the grid itself is enlarged in Fig.  \DIFaddbegin \DIFadd{14C}\DIFaddend ,
topology-dependent runtime differences become readily evident. At the
largest tested grid size (grid size 64), the 8-GPU mean runtimes are
32.54 s (0.54 min) for hexagonal, 266.45 s (4.44 min) for MST, and
880.83 s (14.68 min) for RNG, corresponding to 8-GPU MST and RNG
runtimes that are 8.19x and 27.07x the hexagonal runtime, respectively.

\subsection{7. Final FloatSOM RNG Comparison with
XPySOM}\label{final-floatsom-rng-comparison-with-xpysom}

Fig.  \DIFaddbegin \DIFadd{15 compares untuned hexagonal XPySOM with tuned FloatSOM RNG and
therefore combines implementation, tuning, and topology effects }\DIFaddend (Mancini
et al. 2020). \DIFaddbegin \DIFadd{Sections 5.1, 5.3, and 5.4 separate these components;
Supplementary Figs. S7-S8 provide the corresponding tuned hexagonal and
MST comparisons.
}\DIFaddend

\begin{figure}
\centering
 \DIFaddbeginFL \pandocbounded{\includesvg[keepaspectratio]{assets_manual/figures/fig_15.svg}}
\DIFaddendFL \caption{Figure  \DIFaddbeginFL \DIFaddFL{15}\DIFaddendFL }
\end{figure}

\emph{Figure  \DIFaddbegin \DIFadd{15. }\DIFaddend Integrated deployment comparison of  \DIFaddbegin \DIFadd{untuned }\DIFaddend hexagonal
XPySOM versus tuned FloatSOM RNG\DIFaddbegin \DIFadd{, combining implementation, tuning, and
topology effects}\DIFaddend . Panels A-C compare \(QE_B\), \(QE_H\), and \(QE_T\)
 \DIFaddbegin \DIFadd{under }\DIFaddend full sampling \DIFaddbegin \DIFadd{; panel }\DIFaddend D provides the scaling/runtime context .} At
the overall level, Fig.  \DIFaddbegin \DIFadd{15 }\DIFaddend shows median percentage improvements of
\(QE_B\) (14.5\%); \(QE_H\) (9.1\%); and \(QE_T\) (22.5\%) for tuned
FloatSOM RNG relative to  \DIFaddbegin \DIFadd{untuned }\DIFaddend hexagonal XPySOM, capturing the
combined deployment effect of  \DIFaddbegin \DIFadd{implementation, topology choice, }\DIFaddend and
tuning on \(QE\).

For the  \DIFaddbegin \DIFadd{untuned }\DIFaddend hexagonal XPySOM reference in Fig.  \DIFaddbegin \DIFadd{15}\DIFaddend , workloads beyond
the \(10^8\)-sample case were not processed because they exceeded
available VRAM and XPySOM requires the full dataset to be loaded into
memory. Taken together, Fig.  \DIFaddbegin \DIFadd{15 }\DIFaddend shows the point at which the workload
size is large enough to warrant the extra startup overhead incurred by
tuned FloatSOM RNG: tuned FloatSOM RNG delivers better \(QE\) than the
 \DIFaddbegin \DIFadd{untuned }\DIFaddend hexagonal XPySOM baseline, while also running faster and scaling
to larger workloads.

\subsection{8. Discussion}\label{discussion}

 \DIFaddbegin \DIFadd{FloatSOM combines graph-based SOM topologies with distributed}\DIFaddend ,
out-of-memory  \DIFaddbegin \DIFadd{GPU training. The benchmarks show how sampling, topology,
}\DIFaddend tuning, and workload  \DIFaddbegin \DIFadd{size affect both \(QE\) and runtime}\DIFaddend .

 \DIFaddbegin \DIFadd{The }\DIFaddend trade-off  \DIFaddbegin \DIFadd{between full and random sampling is }\DIFaddend scale dependent.
 \DIFaddbegin \DIFadd{Random sampling is less stable on small datasets, while }\DIFaddend above
\(10{,}000\) samples \DIFaddbegin \DIFadd{we detect little }\DIFaddend paired \(QE\)  \DIFaddbegin \DIFadd{difference}\DIFaddend . Full
sampling is therefore  \DIFaddbegin \DIFadd{preferable }\DIFaddend when stability is  \DIFaddbegin \DIFadd{critical; random
sampling offers higher throughput on larger datasets, although
disk-backed operation limits its }\DIFaddend I/O \DIFaddbegin \DIFadd{advantage}\DIFaddend .

 MST and RNG  \DIFaddbegin \DIFadd{both improve \(QE\) over }\DIFaddend the fixed hexagonal  \DIFaddbegin \DIFadd{lattice and are
close competitors on \(QE\), while RNG gives the stronger MTR result.
This training-time use of MST must be distinguished from the more common
post hoc use of an MST to connect or visualize prototypes learned by a
conventional SOM, as in deployed workflows such as FlowSOM (Van Gassen
et al. 2015). Earlier SOM literature considered MST-defined learning
neighborhoods (}\DIFaddend Kangas et al. 1990\DIFaddbegin \DIFadd{), but FloatSOM recomputes MST or RNG
connectivity from the evolving node weights and uses graph hop distance
directly in the neighborhood update throughout training. A post hoc
graph can change only the displayed connectivity between already learned
prototypes; it cannot change their locations or the resulting \(QE\).
The observed \(QE\) differentials therefore show that MST- and
RNG-mediated training produces different, and on average more
favourable, prototype locations in data space rather than merely a
different set of connections drawn between the same nodes.
}

\DIFadd{A geometric interpretation consistent with these results is that a fixed
lattice imposes uniform, predetermined coupling: when a prototype moves,
it can unnecessarily influence lattice neighbors that are unrelated in
the learned data space, potentially displacing useful prototypes and
increasing dead nodes. The apparently nonlocal hexagonal connections in
the KDD Cup 99 projection (Fig. 5D) provide a visual illustration of
this coupling and the associated restriction on independent prototype
redistribution, subject to the distortions inherent in the
two-dimensional PCA display. MST minimizes such coupling, allowing
prototypes to redistribute more freely along irregular or elongated data
structures. This freedom is also its limitation, because a tree cannot
retain multiple locally appropriate connections where a concentrated
region is better represented by a mesh. RNG occupies an intermediate
position. Its nodes can influence several neighbors, but those
connections arise from the evolving prototype geometry rather than being
imposed before training; RNG can remain sparse where appropriate and
form multiple connections in dense regions. The node-utilization and
dead-node results in Fig. 9, together with \(QE\) and MTR, are
consistent with this account, but they do not identify a unique causal
mechanism (Kohonen 2013; Kangas et al. 1990}\DIFaddend ; Toussaint  \DIFaddbegin \DIFadd{1980)}\DIFaddend .

 \DIFaddbegin \DIFadd{\(QE\) and MTR capture different properties: \(QE\) measures
vector-quantization fidelity, while MTR measures local ordering between
the first and second BMUs. The fixed QE-tuned profiles are associated
with higher MTR in every topology, but the topological cost is far
larger for hexagonal maps: balanced MTR is 24.34 tied-rank positions
higher than in the untuned profile, compared with 1.43 for MST and 1.11
for RNG. The matched radius sweep reinforces this distinction. Around
\(r=1.027\)--1.5, MST and RNG retain nearly flat \(QE\) while MTR is
stable or improves and node utilization remains high; hexagonal maps
instead show a pronounced exchange of lower MTR for worse \(QE\) and
lower node utilization as radius increases (Fig. 9). Thus, }\DIFaddend the  \DIFaddbegin \DIFadd{QE-tuned
}\DIFaddend MST and RNG  \DIFaddbegin \DIFadd{profiles are associated with nowhere near the ordering
penalty observed for the fixed hexagonal lattice. Topology and
hyperparameter choice should therefore be considered jointly rather than
treating topology as a post-training visualization choice. The lower
seed-to-seed variability of MST and RNG may likewise reflect their
ability to rebuild connectivity from evolving node locations, unlike the
fixed hexagonal lattice }\DIFaddend (Kohonen 2013; Kangas et al. 1990).

 \DIFaddbegin \DIFadd{The quality experiments use the in-memory pathway, whereas the scaling
experiments use Ray. Within the distributed benchmark, scaling depends
on workload size: overhead dominates small problems, while larger
workloads benefit from }\DIFaddend parallel throughput and  \DIFaddbegin \DIFadd{may }\DIFaddend remain in RAM  \DIFaddbegin \DIFadd{when
}\DIFaddend the 1-GPU  \DIFaddbegin \DIFadd{case requires disk backing. Efficiencies }\DIFaddend above 100\%  \DIFaddbegin \DIFadd{reflect
this change in memory regime rather than }\DIFaddend super-linear  \DIFaddbegin \DIFadd{computation}\DIFaddend .

 \DIFaddbegin \DIFadd{RNG is the preferred joint }\DIFaddend \(QE\) \DIFaddbegin \DIFadd{/MTR topology when topology overhead is
acceptable; for \(QE\) alone, MST and RNG are close competitors and MST
offers the lower-cost graph option. This recommendation depends on grid
size. At grid size 64, hexagonal, MST, and RNG required 32.54, 266.45,
and 880.83 s on 8 GPUs, respectively. Hexagonal is therefore the
throughput-oriented choice for very large grids, while MST provides a
compromise when a graph topology is desired. More generally, full
sampling favors stability, random sampling favors throughput on large
datasets, and additional GPUs are most useful when they keep the
workload in RAM.
}\DIFaddend

 \DIFaddbegin \DIFadd{FloatSOM brings these topology and systems choices together in a single
}\DIFaddend large-scale  \DIFaddbegin \DIFadd{SOM framework}\DIFaddend .

\subsection{9. Acknowledgements}\label{acknowledgements}

This work was supported by computational resources provided by the
Australian Government through the National Computational Infrastructure
(NCI) under the ANU Merit Allocation Scheme.

We also acknowledge the computational services provided by the
University of Bern, the University of Sydney, and the Walter and Eliza
Hall Institute.

We thank Prof.~Hanna Suominen for her input and advice.

\subsection{10. References}\label{references}

\protect\phantomsection\label{refs}
\begin{CSLReferences}{1}{1}
\bibitem[\citeproctext]{ref-akibaOptunaNextgenerationHyperparameter2019}
Akiba, Takuya, Shotaro Sano, Toshihiko Yanase, Takeru Ohta, and Masanori
Koyama. 2019. \emph{Optuna: {A Next-generation Hyperparameter
Optimization Framework}}. arXiv:1907.10902. arXiv.
\url{https://doi.org/10.48550/arXiv.1907.10902}.

\bibitem[\citeproctext]{ref-akindukoSOMStochasticInitialization2016}
Akinduko, Ayodeji A., Evgeny M. Mirkes, and Alexander N. Gorban. 2016.
{``{SOM}: {Stochastic} Initialization Versus Principal Components.''}
\emph{Information Sciences} 364--365 (October): 213--21.
\url{https://doi.org/10.1016/j.ins.2015.10.013}.

\bibitem[\citeproctext]{ref-alahakoonDynamicSelforganizingMaps2000}
Alahakoon, Damminda, Saman Halgamuge, and Srinivasan Bala. 2000.
{``Dynamic Self-Organizing Maps with Controlled Growth for Knowledge
Discovery. {IEEE Transactions} on {Neural Networks}, 11(3), 601-614.''}
\emph{Neural Networks, IEEE Transactions on} 11 (June): 601--14.
\url{https://doi.org/10.1109/72.846732}.

\bibitem[\citeproctext]{ref-forestSurveyImplementationPerformance2020}
Forest, Florent, Mustapha Lebbah, Hanane Azzag, and Jérôme Lacaille.
2020. {``A {Survey} and {Implementation} of {Performance Metrics} for
{Self-Organized Maps}.''} In \emph{arXiv.org}.
Https://arxiv.org/abs/2011.05847v1.

\DIFaddbegin \bibitem[\citeproctext]{ref-haAweSOMCPUGPUaccelerated2025}
\DIFadd{Ha, Trung, Joonas Nättilä, and Jordy Davelaar. 2025. }{\DIFadd{``}{\DIFadd{aweSOM}}\DIFadd{: A
}{\DIFadd{CPU}}\DIFadd{/}{\DIFadd{GPU-accelerated Self-organizing Map}} \DIFadd{and }{\DIFadd{Statistically Combined
Ensemble Framework}} \DIFadd{for }{\DIFadd{Machine-learning Clustering Analysis}}\DIFadd{.''}}
\DIFadd{\emph{Journal of Open Source Software} 10 (108): 7613.
}\url{https://doi.org/10.21105/joss.07613}\DIFadd{.
}

\DIFaddend \bibitem[\citeproctext]{ref-HPCSystemsNCI}
\emph{{HPC Systems} \textbar{} {NCI}}. n.d.
Https://nci.org.au/our-systems/hpc-systems.

\bibitem[\citeproctext]{ref-itoDesignSpaceExploration2016}
{Ito, Keiichi, Ivo Couckuyt, Roberto d'Ippolito, and Tom Dhaene}. 2016.
{``Design Space Exploration Using {Self-Organizing Map} Based Adaptive
Sampling.''} \emph{Applied Soft Computing} 43 (June): 337--46.
\url{https://doi.org/10.1016/j.asoc.2016.02.036}.

\bibitem[\citeproctext]{ref-jangUseMinimalSpanning2009}
Jang, Yoo-Jin, Myung-Hoe Huh, and Mi-Ra Park. 2009. {``Use of {Minimal
Spanning Trees} on {Self-Organizing Maps}.''} \emph{Korean Journal of
Applied Statistics} 22 (January).
\url{https://doi.org/10.5351/KJAS.2009.22.2.415}.

\bibitem[\citeproctext]{ref-kangasVariantsSelforganizingMaps1990}
Kangas, J. A., T. K. Kohonen, and J. T. Laaksonen. 1990. {``Variants of
Self-Organizing Maps.''} \emph{IEEE Transactions on Neural Networks} 1
(1): 93--99. \url{https://doi.org/10.1109/72.80208}.

\bibitem[\citeproctext]{ref-kohonenSelforganizingMap1990}
Kohonen, T. 1990. {``The Self-Organizing Map.''} \emph{Proceedings of
the IEEE} 78 (9): 1464--80. \url{https://doi.org/10.1109/5.58325}.

\bibitem[\citeproctext]{ref-kohonenSelfOrganizingMaps2001}
Kohonen, Teuvo. 2001. \emph{Self-{Organizing Maps}}. Edited by Thomas S.
Huang, Teuvo Kohonen, and Manfred R. Schroeder. Vol. 30. Springer
{Series} in {Information Sciences}. Springer.
\url{https://doi.org/10.1007/978-3-642-56927-2}.

\bibitem[\citeproctext]{ref-kohonenEssentialsSelforganizingMap2013}
Kohonen, Teuvo. 2013. {``Essentials of the Self-Organizing Map.''}
\emph{Neural Networks: The Official Journal of the International Neural
Network Society} 37 (January): 52--65.
\url{https://doi.org/10.1016/j.neunet.2012.09.018}.

\bibitem[\citeproctext]{ref-kratochvilGigaSOMjlHighperformanceClustering2020}
Kratochvíl, Miroslav, Oliver Hunewald, Laurent Heirendt, et al. 2020.
{``{GigaSOM}.jl: {High-performance} Clustering and Visualization of Huge
Cytometry Datasets.''} \emph{GigaScience} 9 (11): giaa127.
\url{https://doi.org/10.1093/gigascience/giaa127}.

\bibitem[\citeproctext]{ref-kruskalShortestSpanningSubtree1956}
Kruskal, Joseph B. 1956. {``On the Shortest Spanning Subtree of a Graph
and the Traveling Salesman Problem.''} \emph{Proceedings of the American
Mathematical Society} 7 (1): 48--50.
\url{https://doi.org/10.1090/S0002-9939-1956-0078686-7}.

\bibitem[\citeproctext]{ref-liuRobustPhenotypingHighly2023}
Liu, Candace C., Noah F. Greenwald, Alex Kong, et al. 2023. {``Robust
Phenotyping of Highly Multiplexed Tissue Imaging Data Using Pixel-Level
Clustering.''} \emph{Nature Communications} 14 (1): 4618.
\url{https://doi.org/10.1038/s41467-023-40068-5}.

\DIFaddbegin \bibitem[\citeproctext]{ref-machon-gonzalezFLSOMIndividualKernel2010}
\DIFadd{Machón-González, Iván, and Hilario López-García. 2010. }{\DIFadd{``}{\DIFadd{FLSOM}} \DIFadd{with
Individual Kernel Radii Formation and Application to Optimization of a
Pickling Line.''}} \DIFadd{\emph{Engineering Applications of Artificial
Intelligence} 23 (3): 411--19.
}\url{https://doi.org/10.1016/j.engappai.2010.01.016}\DIFadd{.
}

\DIFaddend \bibitem[\citeproctext]{ref-manciniXPySomHighPerformanceSelfOrganizing2020}
Mancini, Riccardo, Antonio Ritacco, Giacomo Lanciano, and Tommaso
Cucinotta. 2020. {``{XPySom}: {High-Performance Self-Organizing
Maps}.''} \emph{2020 {IEEE} 32nd {International Symposium} on {Computer
Architecture} and {High Performance Computing} ({SBAC-PAD})}, September,
209--16. \url{https://doi.org/10.1109/SBAC-PAD49847.2020.00037}.

\bibitem[\citeproctext]{ref-moritzRayDistributedFramework2018}
Moritz, Philipp, Robert Nishihara, Stephanie Wang, et al. 2018. {``Ray:
A Distributed Framework for Emerging {AI} Applications.''}
\emph{Proceedings of the 13th {USENIX} Conference on {Operating Systems
Design} and {Implementation}} (USA), {OSDI}'18, October, 561--77.

\DIFaddbegin \bibitem[\citeproctext]{ref-nemeStatisticalPropertiesLattices2005}
\DIFadd{Neme, Antonio, and Pedro Miramontes. 2005. }{\DIFadd{``Statistical }{\DIFadd{Properties}}
\DIFadd{of }{\DIFadd{Lattices Affect Topographic Error}} \DIFadd{in }{\DIFadd{Self-organizing Maps}}\DIFadd{.''}} \DIFadd{In
\emph{Artificial {Neural Networks}: {Biological Inspirations} -- {ICANN}
2005}, edited by Włodzisław Duch, Janusz Kacprzyk, Erkki Oja, and
Sławomir Zadrożny. Springer. }\url{https://doi.org/10.1007/11550822_67}\DIFadd{.
}

\DIFaddend \bibitem[\citeproctext]{ref-pedregosaScikitlearnMachineLearning2011}
Pedregosa, Fabian, Gaël Varoquaux, Alexandre Gramfort, et al. 2011.
{``Scikit-Learn: {Machine Learning} in {Python}.''} \emph{Journal of
Machine Learning Research} 12 (85): 2825--30.

\DIFaddbegin \bibitem[\citeproctext]{ref-ramosROLELATTICEDIMENSIONALITY2018}
\DIFadd{Ramos, A. Dıaz, E. López-Rubioy, and E. J. Palomo. 2018. }{\DIFadd{``}{\DIFadd{THE ROLE OF
THE LATTICE DIMENSIONALITY IN THE SELF-ORGANIZING MAP}}\DIFadd{.''}} \DIFadd{\emph{Neural
Network World}, no. 1.
}

\DIFaddend \bibitem[\citeproctext]{ref-spanakisAMSOMAdaptiveMoving2016}
Spanakis, Gerasimos, and Gerhard Weiss. 2016. {``{AMSOM}: {Adaptive
Moving Self-organizing Map} for {Clustering} and {Visualization}.''} In
\emph{arXiv.org}. Https://arxiv.org/abs/1605.06047v1.
\url{https://doi.org/10.5220/0005704801290140}.

\bibitem[\citeproctext]{ref-toussaintRelativeNeighbourhoodGraph1980}
Toussaint, Godfried T. 1980. {``The Relative Neighbourhood Graph of a
Finite Planar Set.''} \emph{Pattern Recognition} 12 (4): 261--68.
\url{https://doi.org/10.1016/0031-3203(80)90066-7}.

\DIFaddbegin \bibitem[\citeproctext]{ref-vangassenFlowSOMUsingSelforganizing2015}
\DIFadd{Van Gassen, Sofie, Britt Callebaut, Mary J. Van Helden, et al. 2015.
}{\DIFadd{``}{\DIFadd{FlowSOM}}\DIFadd{: }{\DIFadd{Using}} \DIFadd{Self-Organizing Maps for Visualization and
Interpretation of Cytometry Data.''}} \DIFadd{\emph{Cytometry Part A} 87 (7):
636--45. }\url{https://doi.org/10.1002/cyto.a.22625}\DIFadd{.
}

\DIFaddend \bibitem[\citeproctext]{ref-vasighiDirectedBatchGrowing2017}
Vasighi, Mahdi, and Homa Amini. 2017. {``A Directed Batch Growing
Approach to Enhance the Topology Preservation of Self-Organizing Map.''}
\emph{Applied Soft Computing} 55 (June): 424--35.
\url{https://doi.org/10.1016/j.asoc.2017.02.015}.

\bibitem[\citeproctext]{ref-vettigliJustGlowingMinisom2018}
Vettigli, Giuseppe. 2018. \emph{{JustGlowing}/Minisom}.

\bibitem[\citeproctext]{ref-wetmoreSpeedingSelfOrganizingFeature2005}
Wetmore, Leigh, Malcolm I. Heywood, and A. Nur Zincir-Heywood. 2005.
{``Speeding up the {Self-Organizing Feature Map Using Dynamic Subset
Selection}.''} \emph{Neural Processing Letters} 22 (1): 17--32.
\url{https://doi.org/10.1007/s11063-004-7775-6}.

\bibitem[\citeproctext]{ref-whiteTopologyMattersNetwork2008}
White, Denis, and A. Ross Kiester. 2008. {``Topology Matters: {Network}
Topology Affects Outcomes from Community Ecology Neutral Models.''}
\emph{Computers, Environment and Urban Systems} 32 (2): 165--71.
\url{https://doi.org/10.1016/j.compenvurbsys.2007.11.002}.

\bibitem[\citeproctext]{ref-wittekSomocluEfficientParallel2017}
Wittek, Peter, Shi Chao Gao, Ik Soo Lim, and Li Zhao. 2017. {``Somoclu:
{An Efficient Parallel Library} for {Self-Organizing Maps}.''}
\emph{Journal of Statistical Software} 78 (June): 1--21.
\url{https://doi.org/10.18637/jss.v078.i09}.

\end{CSLReferences}

\subsection{11. Supplementary \DIFaddbegin \DIFadd{Methods and }\DIFaddend Tables (End
Matter)} \DIFaddbegin \label{supplementary-methods-and-tables-end-matter}
\DIFaddend

\DIFaddbegin \subsubsection{\DIFadd{Supplementary Methods S1. Execution-path concordance
protocol}}\label{supplementary-methods-s1.-execution-path-concordance-protocol}

\DIFadd{The execution-path analysis used all 14 datasets, the 10 shared seeds
42--51, full sampling, and the hexagonal, MST, and RNG topologies on one
NVIDIA V100 GPU. For each dataset--seed--topology unit, local CuPy and
Ray streaming received the same standardized data, deterministic
train--holdout split, topology-specific tuned parameters from Section
4.3.1, and evaluation metrics; execution profile was the only
deliberately changed factor. Differences were defined as Ray streaming
minus local CuPy.
}

\DIFadd{For each topology, Supplementary Table S14 reports paired differences
for balanced \(QE\), MTR, node utilization, and dead-node fraction
across 140 matched dataset--seed units. Mean differences and 95\%
confidence intervals were calculated from two-sided one-sample paired
\(t\)-tests against zero, with Benjamini--Hochberg correction applied
jointly across the 12 topology--metric tests.
}

\subsubsection{\DIFadd{Supplementary Tables}}\label{supplementary-tables}

\DIFaddend \textbf{Supplementary Table S1. Dataset metadata and numbered point key
for the Figure 4 sampling mode analysis.}

{\def\LTcaptype{none} % do not increment counter
\begin{longtable}[]{@{}
  >{\raggedright\arraybackslash}p{(\linewidth - 8\tabcolsep) * \real{0.2000}}
  >{\raggedright\arraybackslash}p{(\linewidth - 8\tabcolsep) * \real{0.2000}}
  >{\raggedright\arraybackslash}p{(\linewidth - 8\tabcolsep) * \real{0.2000}}
  >{\raggedright\arraybackslash}p{(\linewidth - 8\tabcolsep) * \real{0.2000}}
  >{\raggedright\arraybackslash}p{(\linewidth - 8\tabcolsep) * \real{0.2000}}@{}}
\toprule\noalign{}
\begin{minipage}[b]{\linewidth}\raggedright
dataset\_index
\end{minipage} & \begin{minipage}[b]{\linewidth}\raggedright
dataset
\end{minipage} & \begin{minipage}[b]{\linewidth}\raggedright
dimension\_count
\end{minipage} & \begin{minipage}[b]{\linewidth}\raggedright
sample\_size
\end{minipage} & \begin{minipage}[b]{\linewidth}\raggedright
dataset\_type
\end{minipage} \\
\midrule\noalign{}
\endhead
\bottomrule\noalign{}
\endlastfoot
1 & iris & 4 & 150 & real \\
2 & wine & 13 & 178 & real \\
3 & olivetti\_faces & 4096 & 400 & real \\
4 & diabetes & 10 & 442 & real \\
5 & breast\_cancer & 30 & 569 & real \\
6 & digits & 64 & 1797 & real \\
7 & california\_housing & 8 & 20640 & real \\
8 & blobs & 2 & 30000 & synthetic \\
9 & circles & 2 & 30000 & synthetic \\
10 & moons & 2 & 30000 & synthetic \\
11 & s\_curve & 3 & 30000 & synthetic \\
12 & swiss\_roll & 3 & 30000 & synthetic \\
13 & kddcup99 & 41 & 494021 & real \\
14 & covertype & 54 & 581012 & real \\
\end{longtable}
}

\textbf{Supplementary Table S2. HDSSSOM pilot configuration summary for
Figure 3.}

{\def\LTcaptype{none} % do not increment counter
\begin{longtable}[]{@{}
  >{\raggedright\arraybackslash}p{(\linewidth - 2\tabcolsep) * \real{0.5000}}
  >{\raggedright\arraybackslash}p{(\linewidth - 2\tabcolsep) * \real{0.5000}}@{}}
\toprule\noalign{}
\begin{minipage}[b]{\linewidth}\raggedright
field
\end{minipage} & \begin{minipage}[b]{\linewidth}\raggedright
value
\end{minipage} \\
\midrule\noalign{}
\endhead
\bottomrule\noalign{}
\endlastfoot
pilot purpose & Initial HDSSSOM screening before the broader sampling
comparison \\
configuration scope & Smaller hexagonal Optuna pilot \\
datasets & \texttt{swiss\_roll}, \texttt{moons}, \texttt{circles},
\texttt{blobs}, \texttt{s\_curve}, \texttt{breast\_cancer},
\texttt{wine}, \texttt{iris}, \texttt{digits},
\texttt{olivetti\_faces} \\
dataset count & 10 \\
seed count & 5 \\
sampling methods present & full, random, HDSSSOM \\
topology & \texttt{hexagonal} \\
algorithm family in campaign & \texttt{ \DIFaddbegin \DIFadd{full\_batch}\DIFaddend } \\
optimization split setup & \texttt{evaluation-split=both}, yielding
\texttt{QE\_H} and \texttt{QE\_T}; Figure 3 reports paired
\texttt{QE\_B} \\
trials per scenario & 200 \\
main text figure slice & hexagonal / full vs HDSSSOM \\
\end{longtable}
}

\textbf{Supplementary Table S3. Exact axis values used for the sample
count, feature dimension, and grid size scaling benchmarks in Section
4.2.1.}

{\def\LTcaptype{none} % do not increment counter
\begin{longtable}[]{@{}
  >{\raggedright\arraybackslash}p{(\linewidth - 2\tabcolsep) * \real{0.5000}}
  >{\raggedright\arraybackslash}p{(\linewidth - 2\tabcolsep) * \real{0.5000}}@{}}
\toprule\noalign{}
\begin{minipage}[b]{\linewidth}\raggedright
scaling axis
\end{minipage} & \begin{minipage}[b]{\linewidth}\raggedright
exact values
\end{minipage} \\
\midrule\noalign{}
\endhead
\bottomrule\noalign{}
\endlastfoot
sample count & \(10^6\), \(5 \times 10^6\), \(10^7\), \(5 \times 10^7\),
\(10^8\), \(5 \times 10^8\), \(10^9\) \\
feature dimension & 50, 100, 200, 500, 1000, 2000, 5000 \\
grid side length & 8, 16, 24, 32, 48, 64 \\
\end{longtable}
}

\textbf{Supplementary Table S4. \DIFaddbegin \DIFadd{Untuned }\DIFaddend FloatSOM \DIFaddbegin \DIFadd{MST }\DIFaddend versus \DIFaddbegin \DIFadd{hexagonal
}\DIFaddend XPySOM .} The \texttt{dataset\_index} column matches the numbered points
in Supplementary Figure S1 panel D.  \DIFaddbegin \DIFadd{Positive percentage values and
positive signed effects favor FloatSOM; the compact embedded table lists
wins as FloatSOM/XPySOM/ties}\DIFaddend .

\DIFaddbegin {\def\LTcaptype{none} %DIF >  do not increment counter
\begin{longtable}[]{@{}
  >{\raggedright\arraybackslash}p{(\linewidth - 20\tabcolsep) * \real{0.0909}}
  >{\raggedright\arraybackslash}p{(\linewidth - 20\tabcolsep) * \real{0.0909}}
  >{\raggedright\arraybackslash}p{(\linewidth - 20\tabcolsep) * \real{0.0909}}
  >{\raggedright\arraybackslash}p{(\linewidth - 20\tabcolsep) * \real{0.0909}}
  >{\raggedright\arraybackslash}p{(\linewidth - 20\tabcolsep) * \real{0.0909}}
  >{\raggedright\arraybackslash}p{(\linewidth - 20\tabcolsep) * \real{0.0909}}
  >{\raggedright\arraybackslash}p{(\linewidth - 20\tabcolsep) * \real{0.0909}}
  >{\raggedright\arraybackslash}p{(\linewidth - 20\tabcolsep) * \real{0.0909}}
  >{\raggedright\arraybackslash}p{(\linewidth - 20\tabcolsep) * \real{0.0909}}
  >{\raggedright\arraybackslash}p{(\linewidth - 20\tabcolsep) * \real{0.0909}}
  >{\raggedright\arraybackslash}p{(\linewidth - 20\tabcolsep) * \real{0.0909}}@{}}
\toprule\noalign{}
\begin{minipage}[b]{\linewidth}\raggedright
\DIFadd{idx
}\end{minipage} & \begin{minipage}[b]{\linewidth}\raggedright
\DIFadd{dataset
}\end{minipage} & \begin{minipage}[b]{\linewidth}\raggedright
\DIFadd{metric
}\end{minipage} & \begin{minipage}[b]{\linewidth}\raggedright
\DIFadd{split
}\end{minipage} & \begin{minipage}[b]{\linewidth}\raggedright
\DIFadd{wins F/X/tie
}\end{minipage} & \begin{minipage}[b]{\linewidth}\raggedright
\DIFadd{median \%
}\end{minipage} & \begin{minipage}[b]{\linewidth}\raggedright
\DIFadd{mean \%
}\end{minipage} & \begin{minipage}[b]{\linewidth}\raggedright
\DIFadd{95\% CI
}\end{minipage} & \begin{minipage}[b]{\linewidth}\raggedright
\DIFadd{p
}\end{minipage} & \begin{minipage}[b]{\linewidth}\raggedright
\DIFadd{n
}\end{minipage} & \begin{minipage}[b]{\linewidth}\raggedright
\DIFadd{signed effect
}\end{minipage} \\
\midrule\noalign{}
\endhead
\bottomrule\noalign{}
\endlastfoot
\DIFadd{8 }& \DIFadd{blobs }& \DIFadd{QE }& \DIFadd{train }& \DIFadd{0/10/0 }& \DIFadd{-3.6521 }& \DIFadd{-3.5612 }& {[}\DIFadd{-4.6291,
-2.687}{]} & \DIFadd{0.002 }& \DIFadd{10 }& \DIFadd{-1 }\\
\DIFadd{5 }& \DIFadd{breast\_cancer }& \DIFadd{QE }& \DIFadd{train }& \DIFadd{10/0/0 }& \DIFadd{6.0461 }& \DIFadd{5.7839 }& {[}\DIFadd{5.0244,
6.3682}{]} & \DIFadd{0.002 }& \DIFadd{10 }& \DIFadd{1 }\\
\DIFadd{7 }& \DIFadd{california\_housing }& \DIFadd{QE }& \DIFadd{train }& \DIFadd{10/0/0 }& \DIFadd{4.2199 }& \DIFadd{4.3194 }&
{[}\DIFadd{3.7617, 4.9195}{]} & \DIFadd{0.002 }& \DIFadd{10 }& \DIFadd{1 }\\
\DIFadd{9 }& \DIFadd{circles }& \DIFadd{QE }& \DIFadd{train }& \DIFadd{0/10/0 }& \DIFadd{-1.7629 }& \DIFadd{-2.1879 }& {[}\DIFadd{-3.655,
-1.0873}{]} & \DIFadd{0.002 }& \DIFadd{10 }& \DIFadd{-1 }\\
\DIFadd{14 }& \DIFadd{covertype }& \DIFadd{QE }& \DIFadd{train }& \DIFadd{10/0/0 }& \DIFadd{15.3837 }& \DIFadd{15.3479 }& {[}\DIFadd{13.9347,
16.925}{]} & \DIFadd{0.002 }& \DIFadd{10 }& \DIFadd{1 }\\
\DIFadd{4 }& \DIFadd{diabetes }& \DIFadd{QE }& \DIFadd{train }& \DIFadd{10/0/0 }& \DIFadd{7.0091 }& \DIFadd{7.0306 }& {[}\DIFadd{6.1295,
7.9882}{]} & \DIFadd{0.002 }& \DIFadd{10 }& \DIFadd{1 }\\
\DIFadd{6 }& \DIFadd{digits }& \DIFadd{QE }& \DIFadd{train }& \DIFadd{10/0/0 }& \DIFadd{5.3038 }& \DIFadd{5.3987 }& {[}\DIFadd{5.0889,
5.8056}{]} & \DIFadd{0.002 }& \DIFadd{10 }& \DIFadd{1 }\\
\DIFadd{1 }& \DIFadd{iris }& \DIFadd{QE }& \DIFadd{train }& \DIFadd{10/0/0 }& \DIFadd{28.4174 }& \DIFadd{28.3671 }& {[}\DIFadd{24.2551,
32.3913}{]} & \DIFadd{0.002 }& \DIFadd{10 }& \DIFadd{1 }\\
\DIFadd{13 }& \DIFadd{kddcup99 }& \DIFadd{QE }& \DIFadd{train }& \DIFadd{10/0/0 }& \DIFadd{19.4133 }& \DIFadd{19.693 }& {[}\DIFadd{16.1432,
23.2133}{]} & \DIFadd{0.002 }& \DIFadd{10 }& \DIFadd{1 }\\
\DIFadd{10 }& \DIFadd{moons }& \DIFadd{QE }& \DIFadd{train }& \DIFadd{1/9/0 }& \DIFadd{-3.7058 }& \DIFadd{-3.3568 }& {[}\DIFadd{-5.5241,
-1.575}{]} & \DIFadd{0.0098 }& \DIFadd{10 }& \DIFadd{-0.8 }\\
\DIFadd{3 }& \DIFadd{olivetti\_faces }& \DIFadd{QE }& \DIFadd{train }& \DIFadd{10/0/0 }& \DIFadd{8.4513 }& \DIFadd{8.5474 }& {[}\DIFadd{7.1779,
9.6246}{]} & \DIFadd{0.002 }& \DIFadd{10 }& \DIFadd{1 }\\
\DIFadd{11 }& \DIFadd{s\_curve }& \DIFadd{QE }& \DIFadd{train }& \DIFadd{10/0/0 }& \DIFadd{3.3827 }& \DIFadd{3.4166 }& {[}\DIFadd{1.906,
5.1934}{]} & \DIFadd{0.002 }& \DIFadd{10 }& \DIFadd{1 }\\
\DIFadd{12 }& \DIFadd{swiss\_roll }& \DIFadd{QE }& \DIFadd{train }& \DIFadd{10/0/0 }& \DIFadd{5.271 }& \DIFadd{5.0983 }& {[}\DIFadd{3.8948,
6.3949}{]} & \DIFadd{0.002 }& \DIFadd{10 }& \DIFadd{1 }\\
\DIFadd{2 }& \DIFadd{wine }& \DIFadd{QE }& \DIFadd{train }& \DIFadd{10/0/0 }& \DIFadd{20.0992 }& \DIFadd{20.1816 }& {[}\DIFadd{18.2671,
22.1871}{]} & \DIFadd{0.002 }& \DIFadd{10 }& \DIFadd{1 }\\
\DIFadd{- }& \DIFadd{GLOBAL }& \DIFadd{QE }& \DIFadd{train }& \DIFadd{111/29/0 }& \DIFadd{6.8961 }& \DIFadd{8.1485 }& {[}\DIFadd{5.6474,
9.0003}{]} & \DIFadd{1.49e-17 }& \DIFadd{140 }& \DIFadd{0.5857 }\\
\DIFadd{8 }& \DIFadd{blobs }& \DIFadd{QE }& \DIFadd{holdout }& \DIFadd{0/10/0 }& \DIFadd{-4.1422 }& \DIFadd{-3.9793 }& {[}\DIFadd{-5.2037,
-2.9583}{]} & \DIFadd{0.002 }& \DIFadd{10 }& \DIFadd{-1 }\\
\DIFadd{5 }& \DIFadd{breast\_cancer }& \DIFadd{QE }& \DIFadd{holdout }& \DIFadd{2/8/0 }& \DIFadd{-0.6423 }& \DIFadd{-0.4735 }&
{[}\DIFadd{-1.2282, 0.3513}{]} & \DIFadd{0.1309 }& \DIFadd{10 }& \DIFadd{-0.6 }\\
\DIFadd{7 }& \DIFadd{california\_housing }& \DIFadd{QE }& \DIFadd{holdout }& \DIFadd{10/0/0 }& \DIFadd{3.8135 }& \DIFadd{3.8367 }&
{[}\DIFadd{3.1391, 4.5531}{]} & \DIFadd{0.002 }& \DIFadd{10 }& \DIFadd{1 }\\
\DIFadd{9 }& \DIFadd{circles }& \DIFadd{QE }& \DIFadd{holdout }& \DIFadd{0/10/0 }& \DIFadd{-2.5762 }& \DIFadd{-3.1314 }& {[}\DIFadd{-4.9792,
-1.4919}{]} & \DIFadd{0.002 }& \DIFadd{10 }& \DIFadd{-1 }\\
\DIFadd{14 }& \DIFadd{covertype }& \DIFadd{QE }& \DIFadd{holdout }& \DIFadd{10/0/0 }& \DIFadd{15.4626 }& \DIFadd{15.3225 }& {[}\DIFadd{13.8381,
16.8307}{]} & \DIFadd{0.002 }& \DIFadd{10 }& \DIFadd{1 }\\
\DIFadd{4 }& \DIFadd{diabetes }& \DIFadd{QE }& \DIFadd{holdout }& \DIFadd{4/6/0 }& \DIFadd{-0.1376 }& \DIFadd{0.0531 }& {[}\DIFadd{-1.2671,
1.4709}{]} & \DIFadd{1 }& \DIFadd{10 }& \DIFadd{-0.2 }\\
\DIFadd{6 }& \DIFadd{digits }& \DIFadd{QE }& \DIFadd{holdout }& \DIFadd{10/0/0 }& \DIFadd{2.5919 }& \DIFadd{2.7034 }& {[}\DIFadd{2.2909,
3.2567}{]} & \DIFadd{0.002 }& \DIFadd{10 }& \DIFadd{1 }\\
\DIFadd{1 }& \DIFadd{iris }& \DIFadd{QE }& \DIFadd{holdout }& \DIFadd{4/6/0 }& \DIFadd{0.1297 }& \DIFadd{-0.4991 }& {[}\DIFadd{-4.9164,
4.9562}{]} & \DIFadd{1 }& \DIFadd{10 }& \DIFadd{-0.2 }\\
\DIFadd{13 }& \DIFadd{kddcup99 }& \DIFadd{QE }& \DIFadd{holdout }& \DIFadd{10/0/0 }& \DIFadd{19.2029 }& \DIFadd{19.4801 }& {[}\DIFadd{15.8968,
23.1041}{]} & \DIFadd{0.002 }& \DIFadd{10 }& \DIFadd{1 }\\
\DIFadd{10 }& \DIFadd{moons }& \DIFadd{QE }& \DIFadd{holdout }& \DIFadd{1/9/0 }& \DIFadd{-3.6371 }& \DIFadd{-3.4648 }& {[}\DIFadd{-5.2819,
-1.8559}{]} & \DIFadd{0.0059 }& \DIFadd{10 }& \DIFadd{-0.8 }\\
\DIFadd{3 }& \DIFadd{olivetti\_faces }& \DIFadd{QE }& \DIFadd{holdout }& \DIFadd{10/0/0 }& \DIFadd{1.915 }& \DIFadd{1.8944 }& {[}\DIFadd{1.55,
2.2547}{]} & \DIFadd{0.002 }& \DIFadd{10 }& \DIFadd{1 }\\
\DIFadd{11 }& \DIFadd{s\_curve }& \DIFadd{QE }& \DIFadd{holdout }& \DIFadd{10/0/0 }& \DIFadd{3.1294 }& \DIFadd{3.2429 }& {[}\DIFadd{1.9458,
4.4824}{]} & \DIFadd{0.002 }& \DIFadd{10 }& \DIFadd{1 }\\
\DIFadd{12 }& \DIFadd{swiss\_roll }& \DIFadd{QE }& \DIFadd{holdout }& \DIFadd{10/0/0 }& \DIFadd{5.1413 }& \DIFadd{5.1294 }& {[}\DIFadd{3.9965,
6.4501}{]} & \DIFadd{0.002 }& \DIFadd{10 }& \DIFadd{1 }\\
\DIFadd{2 }& \DIFadd{wine }& \DIFadd{QE }& \DIFadd{holdout }& \DIFadd{7/3/0 }& \DIFadd{0.1337 }& \DIFadd{-0.0378 }& {[}\DIFadd{-1.2204,
0.9745}{]} & \DIFadd{0.8457 }& \DIFadd{10 }& \DIFadd{0.4 }\\
\DIFadd{- }& \DIFadd{GLOBAL }& \DIFadd{QE }& \DIFadd{holdout }& \DIFadd{88/52/0 }& \DIFadd{1.7869 }& \DIFadd{2.8626 }& {[}\DIFadd{0.8691,
2.7426}{]} & \DIFadd{6.33e-05 }& \DIFadd{140 }& \DIFadd{0.2571 }\\
\DIFadd{8 }& \DIFadd{blobs }& \DIFadd{QE\_B }& \DIFadd{both }& \DIFadd{0/10/0 }& \DIFadd{-3.91 }& \DIFadd{-3.7703 }& {[}\DIFadd{-4.93,
-2.7417}{]} & \DIFadd{0.002 }& \DIFadd{10 }& \DIFadd{-1 }\\
\DIFadd{5 }& \DIFadd{breast\_cancer }& \DIFadd{QE\_B }& \DIFadd{both }& \DIFadd{10/0/0 }& \DIFadd{2.4588 }& \DIFadd{2.4478 }&
{[}\DIFadd{1.8439, 2.987}{]} & \DIFadd{0.002 }& \DIFadd{10 }& \DIFadd{1 }\\
\DIFadd{7 }& \DIFadd{california\_housing }& \DIFadd{QE\_B }& \DIFadd{both }& \DIFadd{10/0/0 }& \DIFadd{3.9644 }& \DIFadd{4.0771 }&
{[}\DIFadd{3.6418, 4.7324}{]} & \DIFadd{0.002 }& \DIFadd{10 }& \DIFadd{1 }\\
\DIFadd{9 }& \DIFadd{circles }& \DIFadd{QE\_B }& \DIFadd{both }& \DIFadd{0/10/0 }& \DIFadd{-2.0639 }& \DIFadd{-2.6605 }& {[}\DIFadd{-4.3931,
-1.3869}{]} & \DIFadd{0.002 }& \DIFadd{10 }& \DIFadd{-1 }\\
\DIFadd{14 }& \DIFadd{covertype }& \DIFadd{QE\_B }& \DIFadd{both }& \DIFadd{10/0/0 }& \DIFadd{15.3975 }& \DIFadd{15.3352 }& {[}\DIFadd{13.8864,
16.8854}{]} & \DIFadd{0.002 }& \DIFadd{10 }& \DIFadd{1 }\\
\DIFadd{4 }& \DIFadd{diabetes }& \DIFadd{QE\_B }& \DIFadd{both }& \DIFadd{10/0/0 }& \DIFadd{3.1657 }& \DIFadd{3.1796 }& {[}\DIFadd{2.2103,
4.1719}{]} & \DIFadd{0.002 }& \DIFadd{10 }& \DIFadd{1 }\\
\DIFadd{6 }& \DIFadd{digits }& \DIFadd{QE\_B }& \DIFadd{both }& \DIFadd{10/0/0 }& \DIFadd{3.997 }& \DIFadd{4.0139 }& {[}\DIFadd{3.6762,
4.3228}{]} & \DIFadd{0.002 }& \DIFadd{10 }& \DIFadd{1 }\\
\DIFadd{1 }& \DIFadd{iris }& \DIFadd{QE\_B }& \DIFadd{both }& \DIFadd{10/0/0 }& \DIFadd{9.1407 }& \DIFadd{9.5616 }& {[}\DIFadd{5.9156,
13.6397}{]} & \DIFadd{0.002 }& \DIFadd{10 }& \DIFadd{1 }\\
\DIFadd{13 }& \DIFadd{kddcup99 }& \DIFadd{QE\_B }& \DIFadd{both }& \DIFadd{10/0/0 }& \DIFadd{19.3064 }& \DIFadd{19.5854 }& {[}\DIFadd{16.0187,
23.1588}{]} & \DIFadd{0.002 }& \DIFadd{10 }& \DIFadd{1 }\\
\DIFadd{10 }& \DIFadd{moons }& \DIFadd{QE\_B }& \DIFadd{both }& \DIFadd{1/9/0 }& \DIFadd{-3.4326 }& \DIFadd{-3.4102 }& {[}\DIFadd{-5.2683,
-1.7163}{]} & \DIFadd{0.0059 }& \DIFadd{10 }& \DIFadd{-0.8 }\\
\DIFadd{3 }& \DIFadd{olivetti\_faces }& \DIFadd{QE\_B }& \DIFadd{both }& \DIFadd{10/0/0 }& \DIFadd{4.8444 }& \DIFadd{4.8685 }&
{[}\DIFadd{4.365, 5.4342}{]} & \DIFadd{0.002 }& \DIFadd{10 }& \DIFadd{1 }\\
\DIFadd{11 }& \DIFadd{s\_curve }& \DIFadd{QE\_B }& \DIFadd{both }& \DIFadd{10/0/0 }& \DIFadd{3.292 }& \DIFadd{3.3299 }& {[}\DIFadd{1.8567,
4.777}{]} & \DIFadd{0.002 }& \DIFadd{10 }& \DIFadd{1 }\\
\DIFadd{12 }& \DIFadd{swiss\_roll }& \DIFadd{QE\_B }& \DIFadd{both }& \DIFadd{10/0/0 }& \DIFadd{5.2033 }& \DIFadd{5.1138 }& {[}\DIFadd{4.0306,
6.4161}{]} & \DIFadd{0.002 }& \DIFadd{10 }& \DIFadd{1 }\\
\DIFadd{2 }& \DIFadd{wine }& \DIFadd{QE\_B }& \DIFadd{both }& \DIFadd{10/0/0 }& \DIFadd{7.8522 }& \DIFadd{7.7725 }& {[}\DIFadd{6.9044, 8.686}{]}
& \DIFadd{0.002 }& \DIFadd{10 }& \DIFadd{1 }\\
\DIFadd{- }& \DIFadd{GLOBAL }& \DIFadd{QE\_B }& \DIFadd{both }& \DIFadd{111/29/0 }& \DIFadd{4.2807 }& \DIFadd{4.9603 }& {[}\DIFadd{3.4978,
5.2317}{]} & \DIFadd{1.35e-13 }& \DIFadd{140 }& \DIFadd{0.5857 }\\
\DIFadd{8 }& \DIFadd{blobs }& \DIFadd{train time }& \DIFadd{train }& \DIFadd{0/10/0 }& \DIFadd{-730.1165 }& \DIFadd{-757.4247 }&
{[}\DIFadd{-864.8227, -727.7904}{]} & \DIFadd{0.002 }& \DIFadd{10 }& \DIFadd{-1 }\\
\DIFadd{5 }& \DIFadd{breast\_cancer }& \DIFadd{train time }& \DIFadd{train }& \DIFadd{0/10/0 }& \DIFadd{-1011.2684 }&
\DIFadd{-1004.5725 }& {[}\DIFadd{-1016.3642, -976.4259}{]} & \DIFadd{0.002 }& \DIFadd{10 }& \DIFadd{-1 }\\
\DIFadd{7 }& \DIFadd{california\_housing }& \DIFadd{train time }& \DIFadd{train }& \DIFadd{0/10/0 }& \DIFadd{-972.0213 }&
\DIFadd{-950.4554 }& {[}\DIFadd{-1011.7249, -888.6339}{]} & \DIFadd{0.002 }& \DIFadd{10 }& \DIFadd{-1 }\\
\DIFadd{9 }& \DIFadd{circles }& \DIFadd{train time }& \DIFadd{train }& \DIFadd{0/10/0 }& \DIFadd{-728.984 }& \DIFadd{-728.8823 }&
{[}\DIFadd{-733.0182, -724.2827}{]} & \DIFadd{0.002 }& \DIFadd{10 }& \DIFadd{-1 }\\
\DIFadd{14 }& \DIFadd{covertype }& \DIFadd{train time }& \DIFadd{train }& \DIFadd{0/10/0 }& \DIFadd{-37.4988 }& \DIFadd{-37.8968 }&
{[}\DIFadd{-42.3539, -32.87}{]} & \DIFadd{0.002 }& \DIFadd{10 }& \DIFadd{-1 }\\
\DIFadd{4 }& \DIFadd{diabetes }& \DIFadd{train time }& \DIFadd{train }& \DIFadd{0/10/0 }& \DIFadd{-1010.0585 }& \DIFadd{-1043.324 }&
{[}\DIFadd{-1183.0245, -1002.8581}{]} & \DIFadd{0.002 }& \DIFadd{10 }& \DIFadd{-1 }\\
\DIFadd{6 }& \DIFadd{digits }& \DIFadd{train time }& \DIFadd{train }& \DIFadd{0/10/0 }& \DIFadd{-1004.8808 }& \DIFadd{-1005.1876 }&
{[}\DIFadd{-1011.5931, -998.5865}{]} & \DIFadd{0.002 }& \DIFadd{10 }& \DIFadd{-1 }\\
\DIFadd{1 }& \DIFadd{iris }& \DIFadd{train time }& \DIFadd{train }& \DIFadd{0/10/0 }& \DIFadd{-1037.1907 }& \DIFadd{-1036.0797 }&
{[}\DIFadd{-1042.8531, -1029.1162}{]} & \DIFadd{0.002 }& \DIFadd{10 }& \DIFadd{-1 }\\
\DIFadd{13 }& \DIFadd{kddcup99 }& \DIFadd{train time }& \DIFadd{train }& \DIFadd{0/10/0 }& \DIFadd{-36.9051 }& \DIFadd{-35.1756 }&
{[}\DIFadd{-38.6464, -32.0512}{]} & \DIFadd{0.002 }& \DIFadd{10 }& \DIFadd{-1 }\\
\DIFadd{10 }& \DIFadd{moons }& \DIFadd{train time }& \DIFadd{train }& \DIFadd{0/10/0 }& \DIFadd{-734.37 }& \DIFadd{-734.8929 }&
{[}\DIFadd{-740.0781, -731.5506}{]} & \DIFadd{0.002 }& \DIFadd{10 }& \DIFadd{-1 }\\
\DIFadd{3 }& \DIFadd{olivetti\_faces }& \DIFadd{train time }& \DIFadd{train }& \DIFadd{0/10/0 }& \DIFadd{-1019.8923 }&
\DIFadd{-1016.4567 }& {[}\DIFadd{-1030.2974, -1000.3188}{]} & \DIFadd{0.002 }& \DIFadd{10 }& \DIFadd{-1 }\\
\DIFadd{11 }& \DIFadd{s\_curve }& \DIFadd{train time }& \DIFadd{train }& \DIFadd{0/10/0 }& \DIFadd{-736.215 }& \DIFadd{-736.5435 }&
{[}\DIFadd{-741.9456, -732.008}{]} & \DIFadd{0.002 }& \DIFadd{10 }& \DIFadd{-1 }\\
\DIFadd{12 }& \DIFadd{swiss\_roll }& \DIFadd{train time }& \DIFadd{train }& \DIFadd{1/9/0 }& \DIFadd{-726.5489 }& \DIFadd{-649.8592 }&
{[}\DIFadd{-732.2277, -336.7629}{]} & \DIFadd{0.0039 }& \DIFadd{10 }& \DIFadd{-0.8 }\\
\DIFadd{2 }& \DIFadd{wine }& \DIFadd{train time }& \DIFadd{train }& \DIFadd{0/10/0 }& \DIFadd{-1030.4764 }& \DIFadd{-1064.4839 }&
{[}\DIFadd{-1204.4846, -1025.0949}{]} & \DIFadd{0.002 }& \DIFadd{10 }& \DIFadd{-1 }\\
\DIFadd{- }& \DIFadd{GLOBAL }& \DIFadd{train time }& \DIFadd{train }& \DIFadd{1/139/0 }& \DIFadd{-868.8976 }& \DIFadd{-771.5168 }&
{[}\DIFadd{-875.2748, -760.3798}{]} & \DIFadd{1.59e-24 }& \DIFadd{140 }& \DIFadd{-0.9857 }\\
\end{longtable}
}

\DIFaddend \textbf{Supplementary Table S5. \DIFaddbegin \DIFadd{Untuned }\DIFaddend FloatSOM \DIFaddbegin \DIFadd{RNG }\DIFaddend versus \DIFaddbegin \DIFadd{hexagonal
}\DIFaddend XPySOM .} The \texttt{dataset\_index} column matches the numbered points
in Supplementary Figure S2 panel D.  \DIFaddbegin \DIFadd{Positive percentage values and
positive signed effects favor FloatSOM; the compact embedded table lists
wins as FloatSOM/XPySOM/ties}\DIFaddend .

\DIFaddbegin {\def\LTcaptype{none} %DIF >  do not increment counter
\begin{longtable}[]{@{}
  >{\raggedright\arraybackslash}p{(\linewidth - 20\tabcolsep) * \real{0.0909}}
  >{\raggedright\arraybackslash}p{(\linewidth - 20\tabcolsep) * \real{0.0909}}
  >{\raggedright\arraybackslash}p{(\linewidth - 20\tabcolsep) * \real{0.0909}}
  >{\raggedright\arraybackslash}p{(\linewidth - 20\tabcolsep) * \real{0.0909}}
  >{\raggedright\arraybackslash}p{(\linewidth - 20\tabcolsep) * \real{0.0909}}
  >{\raggedright\arraybackslash}p{(\linewidth - 20\tabcolsep) * \real{0.0909}}
  >{\raggedright\arraybackslash}p{(\linewidth - 20\tabcolsep) * \real{0.0909}}
  >{\raggedright\arraybackslash}p{(\linewidth - 20\tabcolsep) * \real{0.0909}}
  >{\raggedright\arraybackslash}p{(\linewidth - 20\tabcolsep) * \real{0.0909}}
  >{\raggedright\arraybackslash}p{(\linewidth - 20\tabcolsep) * \real{0.0909}}
  >{\raggedright\arraybackslash}p{(\linewidth - 20\tabcolsep) * \real{0.0909}}@{}}
\toprule\noalign{}
\begin{minipage}[b]{\linewidth}\raggedright
\DIFadd{idx
}\end{minipage} & \begin{minipage}[b]{\linewidth}\raggedright
\DIFadd{dataset
}\end{minipage} & \begin{minipage}[b]{\linewidth}\raggedright
\DIFadd{metric
}\end{minipage} & \begin{minipage}[b]{\linewidth}\raggedright
\DIFadd{split
}\end{minipage} & \begin{minipage}[b]{\linewidth}\raggedright
\DIFadd{wins F/X/tie
}\end{minipage} & \begin{minipage}[b]{\linewidth}\raggedright
\DIFadd{median \%
}\end{minipage} & \begin{minipage}[b]{\linewidth}\raggedright
\DIFadd{mean \%
}\end{minipage} & \begin{minipage}[b]{\linewidth}\raggedright
\DIFadd{95\% CI
}\end{minipage} & \begin{minipage}[b]{\linewidth}\raggedright
\DIFadd{p
}\end{minipage} & \begin{minipage}[b]{\linewidth}\raggedright
\DIFadd{n
}\end{minipage} & \begin{minipage}[b]{\linewidth}\raggedright
\DIFadd{signed effect
}\end{minipage} \\
\midrule\noalign{}
\endhead
\bottomrule\noalign{}
\endlastfoot
\DIFadd{8 }& \DIFadd{blobs }& \DIFadd{QE }& \DIFadd{train }& \DIFadd{0/10/0 }& \DIFadd{-2.7298 }& \DIFadd{-2.6927 }& {[}\DIFadd{-3.7585,
-1.6597}{]} & \DIFadd{0.002 }& \DIFadd{10 }& \DIFadd{-1 }\\
\DIFadd{5 }& \DIFadd{breast\_cancer }& \DIFadd{QE }& \DIFadd{train }& \DIFadd{10/0/0 }& \DIFadd{5.4998 }& \DIFadd{5.465 }& {[}\DIFadd{5.0315,
5.9377}{]} & \DIFadd{0.002 }& \DIFadd{10 }& \DIFadd{1 }\\
\DIFadd{7 }& \DIFadd{california\_housing }& \DIFadd{QE }& \DIFadd{train }& \DIFadd{10/0/0 }& \DIFadd{3.6187 }& \DIFadd{3.641 }&
{[}\DIFadd{3.0746, 4.2359}{]} & \DIFadd{0.002 }& \DIFadd{10 }& \DIFadd{1 }\\
\DIFadd{9 }& \DIFadd{circles }& \DIFadd{QE }& \DIFadd{train }& \DIFadd{1/9/0 }& \DIFadd{-1.7013 }& \DIFadd{-2.1091 }& {[}\DIFadd{-3.9805,
-0.3778}{]} & \DIFadd{0.0137 }& \DIFadd{10 }& \DIFadd{-0.8 }\\
\DIFadd{14 }& \DIFadd{covertype }& \DIFadd{QE }& \DIFadd{train }& \DIFadd{10/0/0 }& \DIFadd{13.6239 }& \DIFadd{13.5016 }& {[}\DIFadd{10.9438,
15.887}{]} & \DIFadd{0.002 }& \DIFadd{10 }& \DIFadd{1 }\\
\DIFadd{4 }& \DIFadd{diabetes }& \DIFadd{QE }& \DIFadd{train }& \DIFadd{10/0/0 }& \DIFadd{6.2882 }& \DIFadd{6.3332 }& {[}\DIFadd{5.2073,
7.5187}{]} & \DIFadd{0.002 }& \DIFadd{10 }& \DIFadd{1 }\\
\DIFadd{6 }& \DIFadd{digits }& \DIFadd{QE }& \DIFadd{train }& \DIFadd{10/0/0 }& \DIFadd{4.4749 }& \DIFadd{4.5437 }& {[}\DIFadd{4.073, 5.0017}{]}
& \DIFadd{0.002 }& \DIFadd{10 }& \DIFadd{1 }\\
\DIFadd{1 }& \DIFadd{iris }& \DIFadd{QE }& \DIFadd{train }& \DIFadd{10/0/0 }& \DIFadd{25.6985 }& \DIFadd{25.9471 }& {[}\DIFadd{20.7469,
31.2163}{]} & \DIFadd{0.002 }& \DIFadd{10 }& \DIFadd{1 }\\
\DIFadd{13 }& \DIFadd{kddcup99 }& \DIFadd{QE }& \DIFadd{train }& \DIFadd{10/0/0 }& \DIFadd{18.8172 }& \DIFadd{18.2227 }& {[}\DIFadd{14.5027,
22.0324}{]} & \DIFadd{0.002 }& \DIFadd{10 }& \DIFadd{1 }\\
\DIFadd{10 }& \DIFadd{moons }& \DIFadd{QE }& \DIFadd{train }& \DIFadd{2/8/0 }& \DIFadd{-2.6397 }& \DIFadd{-2.3844 }& {[}\DIFadd{-4.0907,
-0.5701}{]} & \DIFadd{0.0195 }& \DIFadd{10 }& \DIFadd{-0.6 }\\
\DIFadd{3 }& \DIFadd{olivetti\_faces }& \DIFadd{QE }& \DIFadd{train }& \DIFadd{10/0/0 }& \DIFadd{6.9757 }& \DIFadd{7.4422 }& {[}\DIFadd{5.9654,
9.3791}{]} & \DIFadd{0.002 }& \DIFadd{10 }& \DIFadd{1 }\\
\DIFadd{11 }& \DIFadd{s\_curve }& \DIFadd{QE }& \DIFadd{train }& \DIFadd{10/0/0 }& \DIFadd{3.0675 }& \DIFadd{2.9986 }& {[}\DIFadd{1.5732,
4.5318}{]} & \DIFadd{0.002 }& \DIFadd{10 }& \DIFadd{1 }\\
\DIFadd{12 }& \DIFadd{swiss\_roll }& \DIFadd{QE }& \DIFadd{train }& \DIFadd{10/0/0 }& \DIFadd{5.7819 }& \DIFadd{5.3025 }& {[}\DIFadd{4.1229,
6.2356}{]} & \DIFadd{0.002 }& \DIFadd{10 }& \DIFadd{1 }\\
\DIFadd{2 }& \DIFadd{wine }& \DIFadd{QE }& \DIFadd{train }& \DIFadd{10/0/0 }& \DIFadd{18.7804 }& \DIFadd{18.6217 }& {[}\DIFadd{15.9669,
21.416}{]} & \DIFadd{0.002 }& \DIFadd{10 }& \DIFadd{1 }\\
\DIFadd{- }& \DIFadd{GLOBAL }& \DIFadd{QE }& \DIFadd{train }& \DIFadd{113/27/0 }& \DIFadd{6.1574 }& \DIFadd{7.4881 }& {[}\DIFadd{5.0617,
8.1204}{]} & \DIFadd{5.8e-18 }& \DIFadd{140 }& \DIFadd{0.6143 }\\
\DIFadd{8 }& \DIFadd{blobs }& \DIFadd{QE }& \DIFadd{holdout }& \DIFadd{0/10/0 }& \DIFadd{-3.0403 }& \DIFadd{-2.9146 }& {[}\DIFadd{-3.8854,
-1.8988}{]} & \DIFadd{0.002 }& \DIFadd{10 }& \DIFadd{-1 }\\
\DIFadd{5 }& \DIFadd{breast\_cancer }& \DIFadd{QE }& \DIFadd{holdout }& \DIFadd{3/7/0 }& \DIFadd{-0.1679 }& \DIFadd{-0.0809 }&
{[}\DIFadd{-0.6344, 0.5026}{]} & \DIFadd{0.4316 }& \DIFadd{10 }& \DIFadd{-0.4 }\\
\DIFadd{7 }& \DIFadd{california\_housing }& \DIFadd{QE }& \DIFadd{holdout }& \DIFadd{10/0/0 }& \DIFadd{3.5473 }& \DIFadd{3.4616 }&
{[}\DIFadd{2.8161, 3.9656}{]} & \DIFadd{0.002 }& \DIFadd{10 }& \DIFadd{1 }\\
\DIFadd{9 }& \DIFadd{circles }& \DIFadd{QE }& \DIFadd{holdout }& \DIFadd{2/8/0 }& \DIFadd{-2.435 }& \DIFadd{-2.5817 }& {[}\DIFadd{-4.2591,
-0.7108}{]} & \DIFadd{0.0098 }& \DIFadd{10 }& \DIFadd{-0.6 }\\
\DIFadd{14 }& \DIFadd{covertype }& \DIFadd{QE }& \DIFadd{holdout }& \DIFadd{10/0/0 }& \DIFadd{13.5339 }& \DIFadd{13.4773 }& {[}\DIFadd{11.0872,
15.767}{]} & \DIFadd{0.002 }& \DIFadd{10 }& \DIFadd{1 }\\
\DIFadd{4 }& \DIFadd{diabetes }& \DIFadd{QE }& \DIFadd{holdout }& \DIFadd{8/2/0 }& \DIFadd{0.993 }& \DIFadd{0.9845 }& {[}\DIFadd{0.0524,
1.9348}{]} & \DIFadd{0.0488 }& \DIFadd{10 }& \DIFadd{0.6 }\\
\DIFadd{6 }& \DIFadd{digits }& \DIFadd{QE }& \DIFadd{holdout }& \DIFadd{10/0/0 }& \DIFadd{2.6213 }& \DIFadd{2.6133 }& {[}\DIFadd{2.0362,
3.1169}{]} & \DIFadd{0.002 }& \DIFadd{10 }& \DIFadd{1 }\\
\DIFadd{1 }& \DIFadd{iris }& \DIFadd{QE }& \DIFadd{holdout }& \DIFadd{3/7/0 }& \DIFadd{-1.3579 }& \DIFadd{-1.2948 }& {[}\DIFadd{-4.6315,
2.0699}{]} & \DIFadd{0.375 }& \DIFadd{10 }& \DIFadd{-0.4 }\\
\DIFadd{13 }& \DIFadd{kddcup99 }& \DIFadd{QE }& \DIFadd{holdout }& \DIFadd{10/0/0 }& \DIFadd{18.8346 }& \DIFadd{18.178 }& {[}\DIFadd{14.6321,
21.9247}{]} & \DIFadd{0.002 }& \DIFadd{10 }& \DIFadd{1 }\\
\DIFadd{10 }& \DIFadd{moons }& \DIFadd{QE }& \DIFadd{holdout }& \DIFadd{1/9/0 }& \DIFadd{-2.9592 }& \DIFadd{-2.5631 }& {[}\DIFadd{-4.2128,
-0.9289}{]} & \DIFadd{0.0137 }& \DIFadd{10 }& \DIFadd{-0.8 }\\
\DIFadd{3 }& \DIFadd{olivetti\_faces }& \DIFadd{QE }& \DIFadd{holdout }& \DIFadd{8/2/0 }& \DIFadd{1.5548 }& \DIFadd{1.5684 }&
{[}\DIFadd{0.0193, 3.2179}{]} & \DIFadd{0.0371 }& \DIFadd{10 }& \DIFadd{0.6 }\\
\DIFadd{11 }& \DIFadd{s\_curve }& \DIFadd{QE }& \DIFadd{holdout }& \DIFadd{10/0/0 }& \DIFadd{2.8467 }& \DIFadd{3.0331 }& {[}\DIFadd{1.8244,
4.3726}{]} & \DIFadd{0.002 }& \DIFadd{10 }& \DIFadd{1 }\\
\DIFadd{12 }& \DIFadd{swiss\_roll }& \DIFadd{QE }& \DIFadd{holdout }& \DIFadd{10/0/0 }& \DIFadd{5.5918 }& \DIFadd{5.3626 }& {[}\DIFadd{3.9067,
6.5351}{]} & \DIFadd{0.002 }& \DIFadd{10 }& \DIFadd{1 }\\
\DIFadd{2 }& \DIFadd{wine }& \DIFadd{QE }& \DIFadd{holdout }& \DIFadd{5/5/0 }& \DIFadd{-0.0877 }& \DIFadd{-0.2419 }& {[}\DIFadd{-1.2255,
0.6779}{]} & \DIFadd{0.6953 }& \DIFadd{10 }& \DIFadd{0 }\\
\DIFadd{- }& \DIFadd{GLOBAL }& \DIFadd{QE }& \DIFadd{holdout }& \DIFadd{90/50/0 }& \DIFadd{1.7513 }& \DIFadd{2.7858 }& {[}\DIFadd{0.9907,
2.6251}{]} & \DIFadd{1.02e-05 }& \DIFadd{140 }& \DIFadd{0.2857 }\\
\DIFadd{8 }& \DIFadd{blobs }& \DIFadd{QE\_B }& \DIFadd{both }& \DIFadd{0/10/0 }& \DIFadd{-2.8719 }& \DIFadd{-2.8032 }& {[}\DIFadd{-3.8029,
-1.9134}{]} & \DIFadd{0.002 }& \DIFadd{10 }& \DIFadd{-1 }\\
\DIFadd{5 }& \DIFadd{breast\_cancer }& \DIFadd{QE\_B }& \DIFadd{both }& \DIFadd{10/0/0 }& \DIFadd{2.4644 }& \DIFadd{2.4959 }& {[}\DIFadd{2.075,
2.9683}{]} & \DIFadd{0.002 }& \DIFadd{10 }& \DIFadd{1 }\\
\DIFadd{7 }& \DIFadd{california\_housing }& \DIFadd{QE\_B }& \DIFadd{both }& \DIFadd{10/0/0 }& \DIFadd{3.6308 }& \DIFadd{3.5512 }&
{[}\DIFadd{2.9851, 4.109}{]} & \DIFadd{0.002 }& \DIFadd{10 }& \DIFadd{1 }\\
\DIFadd{9 }& \DIFadd{circles }& \DIFadd{QE\_B }& \DIFadd{both }& \DIFadd{2/8/0 }& \DIFadd{-2.1386 }& \DIFadd{-2.3461 }& {[}\DIFadd{-3.9949,
-0.5829}{]} & \DIFadd{0.0098 }& \DIFadd{10 }& \DIFadd{-0.6 }\\
\DIFadd{14 }& \DIFadd{covertype }& \DIFadd{QE\_B }& \DIFadd{both }& \DIFadd{10/0/0 }& \DIFadd{13.6109 }& \DIFadd{13.4894 }& {[}\DIFadd{11.0155,
15.826}{]} & \DIFadd{0.002 }& \DIFadd{10 }& \DIFadd{1 }\\
\DIFadd{4 }& \DIFadd{diabetes }& \DIFadd{QE\_B }& \DIFadd{both }& \DIFadd{10/0/0 }& \DIFadd{3.3183 }& \DIFadd{3.3817 }& {[}\DIFadd{2.369,
4.1991}{]} & \DIFadd{0.002 }& \DIFadd{10 }& \DIFadd{1 }\\
\DIFadd{6 }& \DIFadd{digits }& \DIFadd{QE\_B }& \DIFadd{both }& \DIFadd{10/0/0 }& \DIFadd{3.5551 }& \DIFadd{3.5512 }& {[}\DIFadd{3.034,
3.9851}{]} & \DIFadd{0.002 }& \DIFadd{10 }& \DIFadd{1 }\\
\DIFadd{1 }& \DIFadd{iris }& \DIFadd{QE\_B }& \DIFadd{both }& \DIFadd{10/0/0 }& \DIFadd{7.914 }& \DIFadd{8.1514 }& {[}\DIFadd{4.7671,
10.8707}{]} & \DIFadd{0.002 }& \DIFadd{10 }& \DIFadd{1 }\\
\DIFadd{13 }& \DIFadd{kddcup99 }& \DIFadd{QE\_B }& \DIFadd{both }& \DIFadd{10/0/0 }& \DIFadd{18.8261 }& \DIFadd{18.1997 }& {[}\DIFadd{14.5584,
21.9789}{]} & \DIFadd{0.002 }& \DIFadd{10 }& \DIFadd{1 }\\
\DIFadd{10 }& \DIFadd{moons }& \DIFadd{QE\_B }& \DIFadd{both }& \DIFadd{2/8/0 }& \DIFadd{-2.8707 }& \DIFadd{-2.4735 }& {[}\DIFadd{-4.1471,
-0.7217}{]} & \DIFadd{0.0195 }& \DIFadd{10 }& \DIFadd{-0.6 }\\
\DIFadd{3 }& \DIFadd{olivetti\_faces }& \DIFadd{QE\_B }& \DIFadd{both }& \DIFadd{10/0/0 }& \DIFadd{4.1052 }& \DIFadd{4.1971 }&
{[}\DIFadd{2.6366, 5.7848}{]} & \DIFadd{0.002 }& \DIFadd{10 }& \DIFadd{1 }\\
\DIFadd{11 }& \DIFadd{s\_curve }& \DIFadd{QE\_B }& \DIFadd{both }& \DIFadd{10/0/0 }& \DIFadd{2.9385 }& \DIFadd{3.0161 }& {[}\DIFadd{1.6726,
4.3754}{]} & \DIFadd{0.002 }& \DIFadd{10 }& \DIFadd{1 }\\
\DIFadd{12 }& \DIFadd{swiss\_roll }& \DIFadd{QE\_B }& \DIFadd{both }& \DIFadd{10/0/0 }& \DIFadd{5.7209 }& \DIFadd{5.3327 }& {[}\DIFadd{4.0146,
6.334}{]} & \DIFadd{0.002 }& \DIFadd{10 }& \DIFadd{1 }\\
\DIFadd{2 }& \DIFadd{wine }& \DIFadd{QE\_B }& \DIFadd{both }& \DIFadd{10/0/0 }& \DIFadd{7.1548 }& \DIFadd{7.0442 }& {[}\DIFadd{5.7434,
8.3772}{]} & \DIFadd{0.002 }& \DIFadd{10 }& \DIFadd{1 }\\
\DIFadd{- }& \DIFadd{GLOBAL }& \DIFadd{QE\_B }& \DIFadd{both }& \DIFadd{114/26/0 }& \DIFadd{3.9446 }& \DIFadd{4.6277 }& {[}\DIFadd{3.2069,
4.8515}{]} & \DIFadd{1.27e-14 }& \DIFadd{140 }& \DIFadd{0.6286 }\\
\DIFadd{8 }& \DIFadd{blobs }& \DIFadd{train time }& \DIFadd{train }& \DIFadd{0/10/0 }& \DIFadd{-725.754 }& \DIFadd{-726.3513 }&
{[}\DIFadd{-729.2884, -724.0566}{]} & \DIFadd{0.002 }& \DIFadd{10 }& \DIFadd{-1 }\\
\DIFadd{5 }& \DIFadd{breast\_cancer }& \DIFadd{train time }& \DIFadd{train }& \DIFadd{0/10/0 }& \DIFadd{-1028.7409 }&
\DIFadd{-1021.2214 }& {[}\DIFadd{-1033.9516, -990.4514}{]} & \DIFadd{0.002 }& \DIFadd{10 }& \DIFadd{-1 }\\
\DIFadd{7 }& \DIFadd{california\_housing }& \DIFadd{train time }& \DIFadd{train }& \DIFadd{0/10/0 }& \DIFadd{-959.1137 }&
\DIFadd{-935.4774 }& {[}\DIFadd{-997.3973, -863.1072}{]} & \DIFadd{0.002 }& \DIFadd{10 }& \DIFadd{-1 }\\
\DIFadd{9 }& \DIFadd{circles }& \DIFadd{train time }& \DIFadd{train }& \DIFadd{0/10/0 }& \DIFadd{-728.9236 }& \DIFadd{-729.3583 }&
{[}\DIFadd{-734.228, -723.7248}{]} & \DIFadd{0.002 }& \DIFadd{10 }& \DIFadd{-1 }\\
\DIFadd{14 }& \DIFadd{covertype }& \DIFadd{train time }& \DIFadd{train }& \DIFadd{0/10/0 }& \DIFadd{-38.0038 }& \DIFadd{-37.1495 }&
{[}\DIFadd{-41.0157, -32.7859}{]} & \DIFadd{0.002 }& \DIFadd{10 }& \DIFadd{-1 }\\
\DIFadd{4 }& \DIFadd{diabetes }& \DIFadd{train time }& \DIFadd{train }& \DIFadd{0/10/0 }& \DIFadd{-1029.3356 }& \DIFadd{-1028.7558 }&
{[}\DIFadd{-1034.5884, -1022.6478}{]} & \DIFadd{0.002 }& \DIFadd{10 }& \DIFadd{-1 }\\
\DIFadd{6 }& \DIFadd{digits }& \DIFadd{train time }& \DIFadd{train }& \DIFadd{0/10/0 }& \DIFadd{-1029.1819 }& \DIFadd{-1028.6713 }&
{[}\DIFadd{-1034.196, -1022.846}{]} & \DIFadd{0.002 }& \DIFadd{10 }& \DIFadd{-1 }\\
\DIFadd{1 }& \DIFadd{iris }& \DIFadd{train time }& \DIFadd{train }& \DIFadd{0/10/0 }& \DIFadd{-1030.5871 }& \DIFadd{-1029.7524 }&
{[}\DIFadd{-1035.7818, -1023.3311}{]} & \DIFadd{0.002 }& \DIFadd{10 }& \DIFadd{-1 }\\
\DIFadd{13 }& \DIFadd{kddcup99 }& \DIFadd{train time }& \DIFadd{train }& \DIFadd{0/10/0 }& \DIFadd{-36.6803 }& \DIFadd{-35.3473 }&
{[}\DIFadd{-39.0245, -32.0055}{]} & \DIFadd{0.002 }& \DIFadd{10 }& \DIFadd{-1 }\\
\DIFadd{10 }& \DIFadd{moons }& \DIFadd{train time }& \DIFadd{train }& \DIFadd{0/10/0 }& \DIFadd{-731.1551 }& \DIFadd{-730.6761 }&
{[}\DIFadd{-733.0006, -727.7656}{]} & \DIFadd{0.002 }& \DIFadd{10 }& \DIFadd{-1 }\\
\DIFadd{3 }& \DIFadd{olivetti\_faces }& \DIFadd{train time }& \DIFadd{train }& \DIFadd{0/10/0 }& \DIFadd{-1054.7323 }&
\DIFadd{-1050.5264 }& {[}\DIFadd{-1064.173, -1034.3248}{]} & \DIFadd{0.002 }& \DIFadd{10 }& \DIFadd{-1 }\\
\DIFadd{11 }& \DIFadd{s\_curve }& \DIFadd{train time }& \DIFadd{train }& \DIFadd{0/10/0 }& \DIFadd{-749.2016 }& \DIFadd{-800.7466 }&
{[}\DIFadd{-1011.2814, -744.996}{]} & \DIFadd{0.002 }& \DIFadd{10 }& \DIFadd{-1 }\\
\DIFadd{12 }& \DIFadd{swiss\_roll }& \DIFadd{train time }& \DIFadd{train }& \DIFadd{1/9/0 }& \DIFadd{-741.1606 }& \DIFadd{-661.2161 }&
{[}\DIFadd{-745.7479, -339.5738}{]} & \DIFadd{0.0039 }& \DIFadd{10 }& \DIFadd{-0.8 }\\
\DIFadd{2 }& \DIFadd{wine }& \DIFadd{train time }& \DIFadd{train }& \DIFadd{0/10/0 }& \DIFadd{-1040.5015 }& \DIFadd{-1039.5076 }&
{[}\DIFadd{-1043.8579, -1034.1717}{]} & \DIFadd{0.002 }& \DIFadd{10 }& \DIFadd{-1 }\\
\DIFadd{- }& \DIFadd{GLOBAL }& \DIFadd{train time }& \DIFadd{train }& \DIFadd{1/139/0 }& \DIFadd{-877.7086 }& \DIFadd{-775.3398 }&
{[}\DIFadd{-884.2881, -753.573}{]} & \DIFadd{1.59e-24 }& \DIFadd{140 }& \DIFadd{-0.9857 }\\
\end{longtable}
}

\DIFaddend \textbf{Supplementary Table S6.  \DIFaddbegin \DIFadd{Hexagonal FloatSOM--XPySOM
implementation }\DIFaddend calibration .} The \texttt{dataset\_index} column matches
the numbered points in Supplementary Figure S3 panel D.  \DIFaddbegin \DIFadd{Positive
percentage values and positive signed effects favor FloatSOM; the
compact embedded table lists wins as FloatSOM/XPySOM/ties}\DIFaddend .

\DIFaddbegin {\def\LTcaptype{none} %DIF >  do not increment counter
\begin{longtable}[]{@{}
  >{\raggedright\arraybackslash}p{(\linewidth - 20\tabcolsep) * \real{0.0909}}
  >{\raggedright\arraybackslash}p{(\linewidth - 20\tabcolsep) * \real{0.0909}}
  >{\raggedright\arraybackslash}p{(\linewidth - 20\tabcolsep) * \real{0.0909}}
  >{\raggedright\arraybackslash}p{(\linewidth - 20\tabcolsep) * \real{0.0909}}
  >{\raggedright\arraybackslash}p{(\linewidth - 20\tabcolsep) * \real{0.0909}}
  >{\raggedright\arraybackslash}p{(\linewidth - 20\tabcolsep) * \real{0.0909}}
  >{\raggedright\arraybackslash}p{(\linewidth - 20\tabcolsep) * \real{0.0909}}
  >{\raggedright\arraybackslash}p{(\linewidth - 20\tabcolsep) * \real{0.0909}}
  >{\raggedright\arraybackslash}p{(\linewidth - 20\tabcolsep) * \real{0.0909}}
  >{\raggedright\arraybackslash}p{(\linewidth - 20\tabcolsep) * \real{0.0909}}
  >{\raggedright\arraybackslash}p{(\linewidth - 20\tabcolsep) * \real{0.0909}}@{}}
\toprule\noalign{}
\begin{minipage}[b]{\linewidth}\raggedright
\DIFadd{idx
}\end{minipage} & \begin{minipage}[b]{\linewidth}\raggedright
\DIFadd{dataset
}\end{minipage} & \begin{minipage}[b]{\linewidth}\raggedright
\DIFadd{metric
}\end{minipage} & \begin{minipage}[b]{\linewidth}\raggedright
\DIFadd{split
}\end{minipage} & \begin{minipage}[b]{\linewidth}\raggedright
\DIFadd{wins F/X/tie
}\end{minipage} & \begin{minipage}[b]{\linewidth}\raggedright
\DIFadd{median \%
}\end{minipage} & \begin{minipage}[b]{\linewidth}\raggedright
\DIFadd{mean \%
}\end{minipage} & \begin{minipage}[b]{\linewidth}\raggedright
\DIFadd{95\% CI
}\end{minipage} & \begin{minipage}[b]{\linewidth}\raggedright
\DIFadd{p
}\end{minipage} & \begin{minipage}[b]{\linewidth}\raggedright
\DIFadd{n
}\end{minipage} & \begin{minipage}[b]{\linewidth}\raggedright
\DIFadd{signed effect
}\end{minipage} \\
\midrule\noalign{}
\endhead
\bottomrule\noalign{}
\endlastfoot
\DIFadd{8 }& \DIFadd{blobs }& \DIFadd{QE }& \DIFadd{train }& \DIFadd{7/3/0 }& \DIFadd{0.0042 }& \DIFadd{0.0034 }& {[}\DIFadd{-0.0023, 0.0089}{]}
& \DIFadd{0.1602 }& \DIFadd{10 }& \DIFadd{0.4 }\\
\DIFadd{5 }& \DIFadd{breast\_cancer }& \DIFadd{QE }& \DIFadd{train }& \DIFadd{1/1/8 }& \DIFadd{1.46e-08 }& \DIFadd{2.92e-09 }&
{[}\DIFadd{-1.04e-05, 1.04e-05}{]} & \DIFadd{1 }& \DIFadd{10 }& \DIFadd{0 }\\
\DIFadd{7 }& \DIFadd{california\_housing }& \DIFadd{QE }& \DIFadd{train }& \DIFadd{7/1/2 }& \DIFadd{1.13e-05 }& \DIFadd{0.0016 }&
{[}\DIFadd{-0.0046, 0.0095}{]} & \DIFadd{0.1484 }& \DIFadd{10 }& \DIFadd{0.75 }\\
\DIFadd{9 }& \DIFadd{circles }& \DIFadd{QE }& \DIFadd{train }& \DIFadd{6/4/0 }& \DIFadd{-0.0156 }& \DIFadd{-0.0142 }& {[}\DIFadd{-0.0446,
0.0107}{]} & \DIFadd{0.5566 }& \DIFadd{10 }& \DIFadd{0.2 }\\
\DIFadd{14 }& \DIFadd{covertype }& \DIFadd{QE }& \DIFadd{train }& \DIFadd{3/7/0 }& \DIFadd{-0.000333 }& \DIFadd{-0.0021 }& {[}\DIFadd{-0.0152,
0.0059}{]} & \DIFadd{0.3223 }& \DIFadd{10 }& \DIFadd{-0.4 }\\
\DIFadd{4 }& \DIFadd{diabetes }& \DIFadd{QE }& \DIFadd{train }& \DIFadd{2/2/6 }& \DIFadd{-8.76e-08 }& \DIFadd{-3.5e-08 }& {[}\DIFadd{-8.88e-06,
8.68e-06}{]} & \DIFadd{0.625 }& \DIFadd{10 }& \DIFadd{0 }\\
\DIFadd{6 }& \DIFadd{digits }& \DIFadd{QE }& \DIFadd{train }& \DIFadd{0/2/8 }& \DIFadd{-1.08e-05 }& \DIFadd{-2.16e-06 }& {[}\DIFadd{-1.08e-05,
-1.08e-05}{]} & \DIFadd{0.5 }& \DIFadd{10 }& \DIFadd{-1 }\\
\DIFadd{1 }& \DIFadd{iris }& \DIFadd{QE }& \DIFadd{train }& \DIFadd{3/1/6 }& \DIFadd{7.85e-06 }& \DIFadd{1.66e-06 }& {[}\DIFadd{-7.64e-06,
8.43e-06}{]} & \DIFadd{0.25 }& \DIFadd{10 }& \DIFadd{0.5 }\\
\DIFadd{13 }& \DIFadd{kddcup99 }& \DIFadd{QE }& \DIFadd{train }& \DIFadd{4/6/0 }& \DIFadd{-0.037 }& \DIFadd{-0.2112 }& {[}\DIFadd{-0.5893,
0.0103}{]} & \DIFadd{0.1309 }& \DIFadd{10 }& \DIFadd{-0.2 }\\
\DIFadd{10 }& \DIFadd{moons }& \DIFadd{QE }& \DIFadd{train }& \DIFadd{8/2/0 }& \DIFadd{0.0021 }& \DIFadd{0.0114 }& {[}\DIFadd{-0.000313,
0.0481}{]} & \DIFadd{0.084 }& \DIFadd{10 }& \DIFadd{0.6 }\\
\DIFadd{3 }& \DIFadd{olivetti\_faces }& \DIFadd{QE }& \DIFadd{train }& \DIFadd{2/4/4 }& \DIFadd{-0.6875 }& \DIFadd{-0.3985 }&
{[}\DIFadd{-1.5487, 0.3296}{]} & \DIFadd{0.2188 }& \DIFadd{10 }& \DIFadd{-0.3333 }\\
\DIFadd{11 }& \DIFadd{s\_curve }& \DIFadd{QE }& \DIFadd{train }& \DIFadd{9/1/0 }& \DIFadd{0.0024 }& \DIFadd{0.0041 }& {[}\DIFadd{1.78e-05,
0.0092}{]} & \DIFadd{0.0195 }& \DIFadd{10 }& \DIFadd{0.8 }\\
\DIFadd{12 }& \DIFadd{swiss\_roll }& \DIFadd{QE }& \DIFadd{train }& \DIFadd{6/4/0 }& \DIFadd{0.0019 }& \DIFadd{0.0019 }& {[}\DIFadd{-0.0046,
0.0087}{]} & \DIFadd{0.625 }& \DIFadd{10 }& \DIFadd{0.2 }\\
\DIFadd{2 }& \DIFadd{wine }& \DIFadd{QE }& \DIFadd{train }& \DIFadd{0/1/9 }& \DIFadd{-1.02e-05 }& \DIFadd{-1.02e-06 }& {[}\DIFadd{-1.02e-05,
-1.02e-05}{]} & \DIFadd{1 }& \DIFadd{10 }& \DIFadd{-1 }\\
\DIFadd{- }& \DIFadd{GLOBAL }& \DIFadd{QE }& \DIFadd{train }& \DIFadd{58/39/43 }& \DIFadd{1.51e-05 }& \DIFadd{-0.0431 }& {[}\DIFadd{-0.000899,
0.0016}{]} & \DIFadd{0.484 }& \DIFadd{140 }& \DIFadd{0.1959 }\\
\DIFadd{8 }& \DIFadd{blobs }& \DIFadd{QE }& \DIFadd{holdout }& \DIFadd{7/3/0 }& \DIFadd{0.0026 }& \DIFadd{0.0032 }& {[}\DIFadd{-0.001,
0.0074}{]} & \DIFadd{0.2324 }& \DIFadd{10 }& \DIFadd{0.4 }\\
\DIFadd{5 }& \DIFadd{breast\_cancer }& \DIFadd{QE }& \DIFadd{holdout }& \DIFadd{1/1/8 }& \DIFadd{-2.38e-07 }& \DIFadd{-4.76e-08 }&
{[}\DIFadd{-9.3e-06, 8.82e-06}{]} & \DIFadd{1 }& \DIFadd{10 }& \DIFadd{0 }\\
\DIFadd{7 }& \DIFadd{california\_housing }& \DIFadd{QE }& \DIFadd{holdout }& \DIFadd{7/2/1 }& \DIFadd{7.32e-06 }& \DIFadd{-0.0011 }&
{[}\DIFadd{-0.0113, 0.0069}{]} & \DIFadd{0.4961 }& \DIFadd{10 }& \DIFadd{0.5556 }\\
\DIFadd{9 }& \DIFadd{circles }& \DIFadd{QE }& \DIFadd{holdout }& \DIFadd{6/4/0 }& \DIFadd{0.0017 }& \DIFadd{0.0011 }& {[}\DIFadd{-0.0504,
0.0522}{]} & \DIFadd{0.9219 }& \DIFadd{10 }& \DIFadd{0.2 }\\
\DIFadd{14 }& \DIFadd{covertype }& \DIFadd{QE }& \DIFadd{holdout }& \DIFadd{3/7/0 }& \DIFadd{-0.000325 }& \DIFadd{-0.0022 }&
{[}\DIFadd{-0.0154, 0.0059}{]} & \DIFadd{0.3223 }& \DIFadd{10 }& \DIFadd{-0.4 }\\
\DIFadd{4 }& \DIFadd{diabetes }& \DIFadd{QE }& \DIFadd{holdout }& \DIFadd{0/2/8 }& \DIFadd{-7.11e-06 }& \DIFadd{-1.42e-06 }&
{[}\DIFadd{-7.17e-06, -7.05e-06}{]} & \DIFadd{0.5 }& \DIFadd{10 }& \DIFadd{-1 }\\
\DIFadd{6 }& \DIFadd{digits }& \DIFadd{QE }& \DIFadd{holdout }& \DIFadd{1/1/8 }& \DIFadd{6.39e-08 }& \DIFadd{1.28e-08 }& {[}\DIFadd{-1.05e-05,
1.06e-05}{]} & \DIFadd{1 }& \DIFadd{10 }& \DIFadd{0 }\\
\DIFadd{1 }& \DIFadd{iris }& \DIFadd{QE }& \DIFadd{holdout }& \DIFadd{4/2/4 }& \DIFadd{5.77e-06 }& \DIFadd{3.63e-06 }& {[}\DIFadd{-8.4e-06,
1.94e-05}{]} & \DIFadd{0.4375 }& \DIFadd{10 }& \DIFadd{0.3333 }\\
\DIFadd{13 }& \DIFadd{kddcup99 }& \DIFadd{QE }& \DIFadd{holdout }& \DIFadd{4/6/0 }& \DIFadd{-0.0381 }& \DIFadd{-0.2277 }& {[}\DIFadd{-0.6099,
0.0098}{]} & \DIFadd{0.1602 }& \DIFadd{10 }& \DIFadd{-0.2 }\\
\DIFadd{10 }& \DIFadd{moons }& \DIFadd{QE }& \DIFadd{holdout }& \DIFadd{6/3/1 }& \DIFadd{0.0026 }& \DIFadd{0.0049 }& {[}\DIFadd{-0.009,
0.0308}{]} & \DIFadd{0.4961 }& \DIFadd{10 }& \DIFadd{0.3333 }\\
\DIFadd{3 }& \DIFadd{olivetti\_faces }& \DIFadd{QE }& \DIFadd{holdout }& \DIFadd{4/2/4 }& \DIFadd{0.2291 }& \DIFadd{0.1813 }&
{[}\DIFadd{-0.3845, 1.4392}{]} & \DIFadd{0.3125 }& \DIFadd{10 }& \DIFadd{0.3333 }\\
\DIFadd{11 }& \DIFadd{s\_curve }& \DIFadd{QE }& \DIFadd{holdout }& \DIFadd{10/0/0 }& \DIFadd{0.003 }& \DIFadd{0.0034 }& {[}\DIFadd{0.000108,
0.0061}{]} & \DIFadd{0.002 }& \DIFadd{10 }& \DIFadd{1 }\\
\DIFadd{12 }& \DIFadd{swiss\_roll }& \DIFadd{QE }& \DIFadd{holdout }& \DIFadd{6/4/0 }& \DIFadd{0.000823 }& \DIFadd{0.0011 }&
{[}\DIFadd{-0.0056, 0.007}{]} & \DIFadd{0.625 }& \DIFadd{10 }& \DIFadd{0.2 }\\
\DIFadd{2 }& \DIFadd{wine }& \DIFadd{QE }& \DIFadd{holdout }& \DIFadd{4/1/5 }& \DIFadd{6.17e-06 }& \DIFadd{2.54e-06 }& {[}\DIFadd{-1.18e-05,
1.31e-05}{]} & \DIFadd{0.4375 }& \DIFadd{10 }& \DIFadd{0.6 }\\
\DIFadd{- }& \DIFadd{GLOBAL }& \DIFadd{QE }& \DIFadd{holdout }& \DIFadd{63/38/39 }& \DIFadd{0.000105 }& \DIFadd{-0.0026 }&
{[}\DIFadd{-0.000154, 0.0021}{]} & \DIFadd{0.1633 }& \DIFadd{140 }& \DIFadd{0.2475 }\\
\DIFadd{8 }& \DIFadd{blobs }& \DIFadd{QE\_B }& \DIFadd{both }& \DIFadd{8/2/0 }& \DIFadd{0.0025 }& \DIFadd{0.0033 }& {[}\DIFadd{0.000283,
0.0064}{]} & \DIFadd{0.0488 }& \DIFadd{10 }& \DIFadd{0.6 }\\
\DIFadd{5 }& \DIFadd{breast\_cancer }& \DIFadd{QE\_B }& \DIFadd{both }& \DIFadd{2/2/6 }& \DIFadd{-1.11e-08 }& \DIFadd{-4.43e-09 }&
{[}\DIFadd{-4.92e-06, 4.92e-06}{]} & \DIFadd{1 }& \DIFadd{10 }& \DIFadd{0 }\\
\DIFadd{7 }& \DIFadd{california\_housing }& \DIFadd{QE\_B }& \DIFadd{both }& \DIFadd{9/1/0 }& \DIFadd{7.42e-06 }& \DIFadd{0.000257 }&
{[}\DIFadd{-0.0069, 0.0082}{]} & \DIFadd{0.0645 }& \DIFadd{10 }& \DIFadd{0.8 }\\
\DIFadd{9 }& \DIFadd{circles }& \DIFadd{QE\_B }& \DIFadd{both }& \DIFadd{6/4/0 }& \DIFadd{-0.0012 }& \DIFadd{-0.0065 }& {[}\DIFadd{-0.0372,
0.0231}{]} & \DIFadd{0.9219 }& \DIFadd{10 }& \DIFadd{0.2 }\\
\DIFadd{14 }& \DIFadd{covertype }& \DIFadd{QE\_B }& \DIFadd{both }& \DIFadd{3/7/0 }& \DIFadd{-0.000329 }& \DIFadd{-0.0022 }&
{[}\DIFadd{-0.0153, 0.0059}{]} & \DIFadd{0.3223 }& \DIFadd{10 }& \DIFadd{-0.4 }\\
\DIFadd{4 }& \DIFadd{diabetes }& \DIFadd{QE\_B }& \DIFadd{both }& \DIFadd{1/3/6 }& \DIFadd{-3.92e-06 }& \DIFadd{-7.82e-07 }&
{[}\DIFadd{-3.93e-06, 3.95e-06}{]} & \DIFadd{0.875 }& \DIFadd{10 }& \DIFadd{-0.5 }\\
\DIFadd{6 }& \DIFadd{digits }& \DIFadd{QE\_B }& \DIFadd{both }& \DIFadd{1/2/7 }& \DIFadd{-3.98e-06 }& \DIFadd{-1.06e-06 }&
{[}\DIFadd{-1.06e-05, 5.35e-06}{]} & \DIFadd{0.75 }& \DIFadd{10 }& \DIFadd{-0.3333 }\\
\DIFadd{1 }& \DIFadd{iris }& \DIFadd{QE\_B }& \DIFadd{both }& \DIFadd{4/2/4 }& \DIFadd{4.93e-06 }& \DIFadd{2.86e-06 }& {[}\DIFadd{-7.76e-06,
1.4e-05}{]} & \DIFadd{0.2188 }& \DIFadd{10 }& \DIFadd{0.3333 }\\
\DIFadd{13 }& \DIFadd{kddcup99 }& \DIFadd{QE\_B }& \DIFadd{both }& \DIFadd{4/6/0 }& \DIFadd{-0.0376 }& \DIFadd{-0.2194 }& {[}\DIFadd{-0.5996,
0.0099}{]} & \DIFadd{0.1602 }& \DIFadd{10 }& \DIFadd{-0.2 }\\
\DIFadd{10 }& \DIFadd{moons }& \DIFadd{QE\_B }& \DIFadd{both }& \DIFadd{7/3/0 }& \DIFadd{0.0018 }& \DIFadd{0.0081 }& {[}\DIFadd{-0.003,
0.0377}{]} & \DIFadd{0.2754 }& \DIFadd{10 }& \DIFadd{0.4 }\\
\DIFadd{3 }& \DIFadd{olivetti\_faces }& \DIFadd{QE\_B }& \DIFadd{both }& \DIFadd{4/3/3 }& \DIFadd{-0.0802 }& \DIFadd{-0.0794 }&
{[}\DIFadd{-0.789, 0.4756}{]} & \DIFadd{0.9375 }& \DIFadd{10 }& \DIFadd{0.1429 }\\
\DIFadd{11 }& \DIFadd{s\_curve }& \DIFadd{QE\_B }& \DIFadd{both }& \DIFadd{10/0/0 }& \DIFadd{0.0037 }& \DIFadd{0.0038 }& {[}\DIFadd{8.5e-05,
0.0073}{]} & \DIFadd{0.002 }& \DIFadd{10 }& \DIFadd{1 }\\
\DIFadd{12 }& \DIFadd{swiss\_roll }& \DIFadd{QE\_B }& \DIFadd{both }& \DIFadd{6/4/0 }& \DIFadd{0.000814 }& \DIFadd{0.0015 }&
{[}\DIFadd{-0.0051, 0.0079}{]} & \DIFadd{0.6953 }& \DIFadd{10 }& \DIFadd{0.2 }\\
\DIFadd{2 }& \DIFadd{wine }& \DIFadd{QE\_B }& \DIFadd{both }& \DIFadd{4/2/4 }& \DIFadd{2.01e-06 }& \DIFadd{1.15e-06 }& {[}\DIFadd{-7.44e-06,
7.84e-06}{]} & \DIFadd{0.6875 }& \DIFadd{10 }& \DIFadd{0.3333 }\\
\DIFadd{- }& \DIFadd{GLOBAL }& \DIFadd{QE\_B }& \DIFadd{both }& \DIFadd{69/41/30 }& \DIFadd{1.58e-05 }& \DIFadd{-0.0208 }&
{[}\DIFadd{-3.92e-06, 0.0011}{]} & \DIFadd{0.0923 }& \DIFadd{140 }& \DIFadd{0.2545 }\\
\DIFadd{8 }& \DIFadd{blobs }& \DIFadd{train time }& \DIFadd{train }& \DIFadd{0/10/0 }& \DIFadd{-82.1972 }& \DIFadd{-82.2698 }&
{[}\DIFadd{-82.616, -81.8815}{]} & \DIFadd{0.002 }& \DIFadd{10 }& \DIFadd{-1 }\\
\DIFadd{5 }& \DIFadd{breast\_cancer }& \DIFadd{train time }& \DIFadd{train }& \DIFadd{0/10/0 }& \DIFadd{-124.4887 }& \DIFadd{-123.0875
}& {[}\DIFadd{-125.5044, -117.1834}{]} & \DIFadd{0.002 }& \DIFadd{10 }& \DIFadd{-1 }\\
\DIFadd{7 }& \DIFadd{california\_housing }& \DIFadd{train time }& \DIFadd{train }& \DIFadd{0/10/0 }& \DIFadd{-112.7705 }&
\DIFadd{-107.7269 }& {[}\DIFadd{-120.1093, -94.1388}{]} & \DIFadd{0.002 }& \DIFadd{10 }& \DIFadd{-1 }\\
\DIFadd{9 }& \DIFadd{circles }& \DIFadd{train time }& \DIFadd{train }& \DIFadd{0/10/0 }& \DIFadd{-82.011 }& \DIFadd{-82.0468 }&
{[}\DIFadd{-82.6543, -81.5143}{]} & \DIFadd{0.002 }& \DIFadd{10 }& \DIFadd{-1 }\\
\DIFadd{14 }& \DIFadd{covertype }& \DIFadd{train time }& \DIFadd{train }& \DIFadd{8/2/0 }& \DIFadd{4.5544 }& \DIFadd{3.1364 }&
{[}\DIFadd{-0.2416, 4.7034}{]} & \DIFadd{0.1055 }& \DIFadd{10 }& \DIFadd{0.6 }\\
\DIFadd{4 }& \DIFadd{diabetes }& \DIFadd{train time }& \DIFadd{train }& \DIFadd{0/10/0 }& \DIFadd{-124.7165 }& \DIFadd{-130.8947 }&
{[}\DIFadd{-156.4431, -123.3642}{]} & \DIFadd{0.002 }& \DIFadd{10 }& \DIFadd{-1 }\\
\DIFadd{6 }& \DIFadd{digits }& \DIFadd{train time }& \DIFadd{train }& \DIFadd{0/10/0 }& \DIFadd{-124.2555 }& \DIFadd{-124.4075 }&
{[}\DIFadd{-126.3659, -122.7873}{]} & \DIFadd{0.002 }& \DIFadd{10 }& \DIFadd{-1 }\\
\DIFadd{1 }& \DIFadd{iris }& \DIFadd{train time }& \DIFadd{train }& \DIFadd{0/10/0 }& \DIFadd{-126.2679 }& \DIFadd{-127.1989 }&
{[}\DIFadd{-132.6232, -124.6813}{]} & \DIFadd{0.002 }& \DIFadd{10 }& \DIFadd{-1 }\\
\DIFadd{13 }& \DIFadd{kddcup99 }& \DIFadd{train time }& \DIFadd{train }& \DIFadd{10/0/0 }& \DIFadd{12.295 }& \DIFadd{12.5548 }&
{[}\DIFadd{8.0787, 15.7426}{]} & \DIFadd{0.002 }& \DIFadd{10 }& \DIFadd{1 }\\
\DIFadd{10 }& \DIFadd{moons }& \DIFadd{train time }& \DIFadd{train }& \DIFadd{0/10/0 }& \DIFadd{-82.0957 }& \DIFadd{-88.2408 }&
{[}\DIFadd{-112.9703, -81.6979}{]} & \DIFadd{0.002 }& \DIFadd{10 }& \DIFadd{-1 }\\
\DIFadd{3 }& \DIFadd{olivetti\_faces }& \DIFadd{train time }& \DIFadd{train }& \DIFadd{0/10/0 }& \DIFadd{-200.8962 }&
\DIFadd{-199.5774 }& {[}\DIFadd{-203.2204, -194.7869}{]} & \DIFadd{0.002 }& \DIFadd{10 }& \DIFadd{-1 }\\
\DIFadd{11 }& \DIFadd{s\_curve }& \DIFadd{train time }& \DIFadd{train }& \DIFadd{0/10/0 }& \DIFadd{-83.1004 }& \DIFadd{-83.8438 }&
{[}\DIFadd{-87.0546, -82.4761}{]} & \DIFadd{0.002 }& \DIFadd{10 }& \DIFadd{-1 }\\
\DIFadd{12 }& \DIFadd{swiss\_roll }& \DIFadd{train time }& \DIFadd{train }& \DIFadd{1/9/0 }& \DIFadd{-81.642 }& \DIFadd{-65.9893 }&
{[}\DIFadd{-82.6048, -3.0855}{]} & \DIFadd{0.0039 }& \DIFadd{10 }& \DIFadd{-0.8 }\\
\DIFadd{2 }& \DIFadd{wine }& \DIFadd{train time }& \DIFadd{train }& \DIFadd{0/10/0 }& \DIFadd{-126.5213 }& \DIFadd{-126.5147 }&
{[}\DIFadd{-127.7591, -125.2633}{]} & \DIFadd{0.002 }& \DIFadd{10 }& \DIFadd{-1 }\\
\DIFadd{- }& \DIFadd{GLOBAL }& \DIFadd{train time }& \DIFadd{train }& \DIFadd{19/121/0 }& \DIFadd{-102.7688 }& \DIFadd{-94.7219 }&
{[}\DIFadd{-103.9162, -95.6444}{]} & \DIFadd{1.02e-22 }& \DIFadd{140 }& \DIFadd{-0.7286 }\\
\end{longtable}
}

\DIFaddend \textbf{Supplementary Table S7. Paired  \DIFaddbegin \DIFadd{topology-comparison effects }\DIFaddend for
hexagonal versus MST and hexagonal versus RNG across  \DIFaddbegin \DIFadd{Balanced }\DIFaddend QE,
 \DIFaddbegin \DIFadd{Holdout }\DIFaddend QE, and  \DIFaddbegin \DIFadd{Train }\DIFaddend QE.} Rows list metric/dataset entries, including
the OVERALL row.  \DIFaddbegin \DIFadd{Effect columns report the paired mean percent
improvement of hexagonal over the comparator topology, using hexagonal
QE as the reference denominator; positive values favor hexagonal and
negative values favor MST or RNG. The CI columns give the corresponding
95\% paired \(t\)-test confidence intervals. Raw }\DIFaddend p-values  \DIFaddbegin \DIFadd{are retained
for audit, and q-values report Benjamini-Hochberg adjustment across the
42 dataset-level tests in each topology-comparison family (14 datasets
by three QE endpoints, separately for MST and RNG); OVERALL rows are
pooled summaries and retain q=NA. The embedded table is reproduced from
}\DIFaddend \texttt{assets/tables/supp\_table\_topology\_hex\_vs\_mst\_rng\_pvalues.tsv}.

\DIFaddbegin {\def\LTcaptype{none} %DIF >  do not increment counter
\begin{longtable}[]{@{}
  >{\raggedright\arraybackslash}p{(\linewidth - 18\tabcolsep) * \real{0.1000}}
  >{\raggedright\arraybackslash}p{(\linewidth - 18\tabcolsep) * \real{0.1000}}
  >{\raggedright\arraybackslash}p{(\linewidth - 18\tabcolsep) * \real{0.1000}}
  >{\raggedright\arraybackslash}p{(\linewidth - 18\tabcolsep) * \real{0.1000}}
  >{\raggedright\arraybackslash}p{(\linewidth - 18\tabcolsep) * \real{0.1000}}
  >{\raggedright\arraybackslash}p{(\linewidth - 18\tabcolsep) * \real{0.1000}}
  >{\raggedright\arraybackslash}p{(\linewidth - 18\tabcolsep) * \real{0.1000}}
  >{\raggedright\arraybackslash}p{(\linewidth - 18\tabcolsep) * \real{0.1000}}
  >{\raggedright\arraybackslash}p{(\linewidth - 18\tabcolsep) * \real{0.1000}}
  >{\raggedright\arraybackslash}p{(\linewidth - 18\tabcolsep) * \real{0.1000}}@{}}
\toprule\noalign{}
\begin{minipage}[b]{\linewidth}\raggedright
\DIFadd{metric
}\end{minipage} & \begin{minipage}[b]{\linewidth}\raggedright
\DIFadd{dataset
}\end{minipage} & \begin{minipage}[b]{\linewidth}\raggedright
\DIFadd{MST\_effect\_pct
}\end{minipage} & \begin{minipage}[b]{\linewidth}\raggedright
\DIFadd{MST\_95pct\_CI
}\end{minipage} & \begin{minipage}[b]{\linewidth}\raggedright
\DIFadd{MST\_p
}\end{minipage} & \begin{minipage}[b]{\linewidth}\raggedright
\DIFadd{MST\_q
}\end{minipage} & \begin{minipage}[b]{\linewidth}\raggedright
\DIFadd{RNG\_effect\_pct
}\end{minipage} & \begin{minipage}[b]{\linewidth}\raggedright
\DIFadd{RNG\_95pct\_CI
}\end{minipage} & \begin{minipage}[b]{\linewidth}\raggedright
\DIFadd{RNG\_p
}\end{minipage} & \begin{minipage}[b]{\linewidth}\raggedright
\DIFadd{RNG\_q
}\end{minipage} \\
\midrule\noalign{}
\endhead
\bottomrule\noalign{}
\endlastfoot
\DIFadd{Balanced QE }& \DIFadd{blobs }& \DIFadd{-0.234 }& {[}\DIFadd{-0.4136, -0.0543}{]} & \DIFadd{p=0.0163 }&
\DIFadd{q=0.0403 }& \DIFadd{-0.4278 }& {[}\DIFadd{-0.6395, -0.2161}{]} & \DIFadd{p=0.00134 }& \DIFadd{q=0.00353 }\\
\DIFadd{Balanced QE }& \DIFadd{circles }& \DIFadd{-0.0117 }& {[}\DIFadd{-0.1666, 0.1432}{]} & \DIFadd{p=0.868 }&
\DIFadd{q=0.889 }& \DIFadd{-0.1507 }& {[}\DIFadd{-0.28, -0.02145}{]} & \DIFadd{p=0.027 }& \DIFadd{q=0.0494 }\\
\DIFadd{Balanced QE }& \DIFadd{moons }& \DIFadd{0.04526 }& {[}\DIFadd{-0.09841, 0.1889}{]} & \DIFadd{p=0.494 }&
\DIFadd{q=0.576 }& \DIFadd{-0.1834 }& {[}\DIFadd{-0.3183, -0.04856}{]} & \DIFadd{p=0.0132 }& \DIFadd{q=0.0277 }\\
\DIFadd{Balanced QE }& \DIFadd{s\_curve }& \DIFadd{-0.2258 }& {[}\DIFadd{-0.324, -0.1277}{]} & \DIFadd{p=0.000559 }&
\DIFadd{q=0.00261 }& \DIFadd{-0.2554 }& {[}\DIFadd{-0.3406, -0.1702}{]} & \DIFadd{p=8.08e-05 }&
\DIFadd{q=0.000565 }\\
\DIFadd{Balanced QE }& \DIFadd{swiss\_roll }& \DIFadd{-0.2335 }& {[}\DIFadd{-0.6601, 0.1931}{]} & \DIFadd{p=0.247 }&
\DIFadd{q=0.358 }& \DIFadd{-0.3258 }& {[}\DIFadd{-0.6783, 0.02675}{]} & \DIFadd{p=0.0661 }& \DIFadd{q=0.0992 }\\
\DIFadd{Balanced QE }& \DIFadd{breast\_cancer }& \DIFadd{0.2996 }& {[}\DIFadd{-0.1996, 0.7987}{]} & \DIFadd{p=0.208
}& \DIFadd{q=0.336 }& \DIFadd{-0.1832 }& {[}\DIFadd{-0.6168, 0.2505}{]} & \DIFadd{p=0.364 }& \DIFadd{q=0.413 }\\
\DIFadd{Balanced QE }& \DIFadd{wine }& \DIFadd{-3.999 }& {[}\DIFadd{-5.789, -2.209}{]} & \DIFadd{p=0.000686 }&
\DIFadd{q=0.00288 }& \DIFadd{-5.319 }& {[}\DIFadd{-7.359, -3.28}{]} & \DIFadd{p=0.000229 }& \DIFadd{q=0.0012 }\\
\DIFadd{Balanced QE }& \DIFadd{iris }& \DIFadd{-1.004 }& {[}\DIFadd{-4.964, 2.955}{]} & \DIFadd{p=0.58 }& \DIFadd{q=0.641 }&
\DIFadd{-4.83 }& {[}\DIFadd{-7.089, -2.572}{]} & \DIFadd{p=0.000924 }& \DIFadd{q=0.00299 }\\
\DIFadd{Balanced QE }& \DIFadd{digits }& \DIFadd{-0.1745 }& {[}\DIFadd{-0.4361, 0.08697}{]} & \DIFadd{p=0.165 }&
\DIFadd{q=0.298 }& \DIFadd{-0.1445 }& {[}\DIFadd{-0.4203, 0.1314}{]} & \DIFadd{p=0.266 }& \DIFadd{q=0.32 }\\
\DIFadd{Balanced QE }& \DIFadd{olivetti\_faces }& \DIFadd{-1.265 }& {[}\DIFadd{-2.634, 0.1046}{]} & \DIFadd{p=0.0663
}& \DIFadd{q=0.146 }& \DIFadd{-1.23 }& {[}\DIFadd{-2.492, 0.03187}{]} & \DIFadd{p=0.0549 }& \DIFadd{q=0.0887 }\\
\DIFadd{Balanced QE }& \DIFadd{diabetes }& \DIFadd{0.1492 }& {[}\DIFadd{-0.5331, 0.8315}{]} & \DIFadd{p=0.633 }&
\DIFadd{q=0.681 }& \DIFadd{-0.3543 }& {[}\DIFadd{-1.031, 0.3227}{]} & \DIFadd{p=0.267 }& \DIFadd{q=0.32 }\\
\DIFadd{Balanced QE }& \DIFadd{california\_housing }& \DIFadd{-0.28 }& {[}\DIFadd{-0.3942, -0.1659}{]} &
\DIFadd{p=0.000356 }& \DIFadd{q=0.00187 }& \DIFadd{-0.184 }& {[}\DIFadd{-0.2426, -0.1254}{]} & \DIFadd{p=5.65e-05 }&
\DIFadd{q=0.000565 }\\
\DIFadd{Balanced QE }& \DIFadd{covertype }& \DIFadd{-0.3486 }& {[}\DIFadd{-1.411, 0.7143}{]} & \DIFadd{p=0.477 }&
\DIFadd{q=0.574 }& \DIFadd{-0.1093 }& {[}\DIFadd{-0.857, 0.6383}{]} & \DIFadd{p=0.748 }& \DIFadd{q=0.786 }\\
\DIFadd{Balanced QE }& \DIFadd{kddcup99 }& \DIFadd{-5.859 }& {[}\DIFadd{-6.623, -5.096}{]} & \DIFadd{p=3.14e-08 }&
\DIFadd{q=6.6e-07 }& \DIFadd{-4.466 }& {[}\DIFadd{-5.768, -3.165}{]} & \DIFadd{p=2.81e-05 }& \DIFadd{q=0.000394 }\\
\DIFadd{Balanced QE }& \DIFadd{OVERALL }& \DIFadd{-0.9387 }& {[}\DIFadd{-1.346, -0.5316}{]} & \DIFadd{p=1.12e-05 }&
\DIFadd{q=NA }& \DIFadd{-1.297 }& {[}\DIFadd{-1.685, -0.9097}{]} & \DIFadd{p=7.4e-10 }& \DIFadd{q=NA }\\
\DIFadd{Holdout QE }& \DIFadd{blobs }& \DIFadd{-0.1188 }& {[}\DIFadd{-0.4239, 0.1864}{]} & \DIFadd{p=0.401 }& \DIFadd{q=0.544
}& \DIFadd{-0.4212 }& {[}\DIFadd{-0.8076, -0.03478}{]} & \DIFadd{p=0.0358 }& \DIFadd{q=0.0627 }\\
\DIFadd{Holdout QE }& \DIFadd{circles }& \DIFadd{0.2031 }& {[}\DIFadd{-0.05154, 0.4578}{]} & \DIFadd{p=0.105 }&
\DIFadd{q=0.209 }& \DIFadd{-0.02455 }& {[}\DIFadd{-0.2726, 0.2235}{]} & \DIFadd{p=0.828 }& \DIFadd{q=0.828 }\\
\DIFadd{Holdout QE }& \DIFadd{moons }& \DIFadd{0.1236 }& {[}\DIFadd{-0.08663, 0.3339}{]} & \DIFadd{p=0.216 }& \DIFadd{q=0.336
}& \DIFadd{-0.2158 }& {[}\DIFadd{-0.3211, -0.1106}{]} & \DIFadd{p=0.00122 }& \DIFadd{q=0.00342 }\\
\DIFadd{Holdout QE }& \DIFadd{s\_curve }& \DIFadd{-0.1582 }& {[}\DIFadd{-0.2432, -0.07327}{]} & \DIFadd{p=0.00226 }&
\DIFadd{q=0.00792 }& \DIFadd{-0.1842 }& {[}\DIFadd{-0.2673, -0.1011}{]} & \DIFadd{p=0.000724 }& \DIFadd{q=0.00299 }\\
\DIFadd{Holdout QE }& \DIFadd{swiss\_roll }& \DIFadd{-0.303 }& {[}\DIFadd{-0.6657, 0.05966}{]} & \DIFadd{p=0.0913 }&
\DIFadd{q=0.192 }& \DIFadd{-0.3654 }& {[}\DIFadd{-0.6648, -0.06609}{]} & \DIFadd{p=0.0221 }& \DIFadd{q=0.0421 }\\
\DIFadd{Holdout QE }& \DIFadd{breast\_cancer }& \DIFadd{0.4924 }& {[}\DIFadd{0.1735, 0.8113}{]} & \DIFadd{p=0.0068 }&
\DIFadd{q=0.0204 }& \DIFadd{0.2997 }& {[}\DIFadd{-0.0605, 0.6599}{]} & \DIFadd{p=0.0925 }& \DIFadd{q=0.129 }\\
\DIFadd{Holdout QE }& \DIFadd{wine }& \DIFadd{1.582 }& {[}\DIFadd{0.9599, 2.203}{]} & \DIFadd{p=0.000275 }& \DIFadd{q=0.00165
}& \DIFadd{1.148 }& {[}\DIFadd{0.524, 1.772}{]} & \DIFadd{p=0.00244 }& \DIFadd{q=0.0057 }\\
\DIFadd{Holdout QE }& \DIFadd{iris }& \DIFadd{2.725 }& {[}\DIFadd{0.7442, 4.705}{]} & \DIFadd{p=0.0125 }& \DIFadd{q=0.0349 }&
\DIFadd{2.032 }& {[}\DIFadd{0.06363, 4.001}{]} & \DIFadd{p=0.0444 }& \DIFadd{q=0.0745 }\\
\DIFadd{Holdout QE }& \DIFadd{digits }& \DIFadd{-0.2882 }& {[}\DIFadd{-0.5927, 0.01621}{]} & \DIFadd{p=0.0609 }&
\DIFadd{q=0.142 }& \DIFadd{-0.1447 }& {[}\DIFadd{-0.4001, 0.1107}{]} & \DIFadd{p=0.232 }& \DIFadd{q=0.305 }\\
\DIFadd{Holdout QE }& \DIFadd{olivetti\_faces }& \DIFadd{-1.322 }& {[}\DIFadd{-1.723, -0.9199}{]} &
\DIFadd{p=3.93e-05 }& \DIFadd{q=0.000412 }& \DIFadd{-1.713 }& {[}\DIFadd{-2.152, -1.275}{]} & \DIFadd{p=9.96e-06 }&
\DIFadd{q=0.000333 }\\
\DIFadd{Holdout QE }& \DIFadd{diabetes }& \DIFadd{0.003896 }& {[}\DIFadd{-0.4209, 0.4287}{]} & \DIFadd{p=0.984 }&
\DIFadd{q=0.984 }& \DIFadd{-0.273 }& {[}\DIFadd{-0.901, 0.355}{]} & \DIFadd{p=0.351 }& \DIFadd{q=0.41 }\\
\DIFadd{Holdout QE }& \DIFadd{california\_housing }& \DIFadd{-0.2211 }& {[}\DIFadd{-0.3483, -0.09394}{]} &
\DIFadd{p=0.00344 }& \DIFadd{q=0.0111 }& \DIFadd{-0.1404 }& {[}\DIFadd{-0.247, -0.03388}{]} & \DIFadd{p=0.0154 }&
\DIFadd{q=0.0308 }\\
\DIFadd{Holdout QE }& \DIFadd{covertype }& \DIFadd{-0.3532 }& {[}\DIFadd{-1.419, 0.7122}{]} & \DIFadd{p=0.472 }&
\DIFadd{q=0.574 }& \DIFadd{-0.08568 }& {[}\DIFadd{-0.8426, 0.6712}{]} & \DIFadd{p=0.804 }& \DIFadd{q=0.823 }\\
\DIFadd{Holdout QE }& \DIFadd{kddcup99 }& \DIFadd{-5.871 }& {[}\DIFadd{-6.673, -5.068}{]} & \DIFadd{p=4.8e-08 }&
\DIFadd{q=6.72e-07 }& \DIFadd{-4.643 }& {[}\DIFadd{-5.902, -3.383}{]} & \DIFadd{p=1.58e-05 }& \DIFadd{q=0.000333 }\\
\DIFadd{Holdout QE }& \DIFadd{OVERALL }& \DIFadd{-0.2504 }& {[}\DIFadd{-0.5924, 0.09159}{]} & \DIFadd{p=0.15 }& \DIFadd{q=NA
}& \DIFadd{-0.3379 }& {[}\DIFadd{-0.6291, -0.0468}{]} & \DIFadd{p=0.0232 }& \DIFadd{q=NA }\\
\DIFadd{Train QE }& \DIFadd{blobs }& \DIFadd{-0.2851 }& {[}\DIFadd{-0.5038, -0.06639}{]} & \DIFadd{p=0.0162 }&
\DIFadd{q=0.0403 }& \DIFadd{-0.4429 }& {[}\DIFadd{-0.6468, -0.2391}{]} & \DIFadd{p=0.00083 }& \DIFadd{q=0.00299 }\\
\DIFadd{Train QE }& \DIFadd{circles }& \DIFadd{-0.1236 }& {[}\DIFadd{-0.2862, 0.03895}{]} & \DIFadd{p=0.119 }&
\DIFadd{q=0.228 }& \DIFadd{-0.1656 }& {[}\DIFadd{-0.2318, -0.09945}{]} & \DIFadd{p=0.000309 }& \DIFadd{q=0.00144 }\\
\DIFadd{Train QE }& \DIFadd{moons }& \DIFadd{-0.04045 }& {[}\DIFadd{-0.1551, 0.07423}{]} & \DIFadd{p=0.445 }& \DIFadd{q=0.574
}& \DIFadd{-0.157 }& {[}\DIFadd{-0.2696, -0.04434}{]} & \DIFadd{p=0.0117 }& \DIFadd{q=0.0258 }\\
\DIFadd{Train QE }& \DIFadd{s\_curve }& \DIFadd{-0.2144 }& {[}\DIFadd{-0.2961, -0.1327}{]} & \DIFadd{p=0.000219 }&
\DIFadd{q=0.00153 }& \DIFadd{-0.2643 }& {[}\DIFadd{-0.3584, -0.1702}{]} & \DIFadd{p=0.000132 }&
\DIFadd{q=0.000794 }\\
\DIFadd{Train QE }& \DIFadd{swiss\_roll }& \DIFadd{-0.2022 }& {[}\DIFadd{-0.6787, 0.2742}{]} & \DIFadd{p=0.362 }&
\DIFadd{q=0.507 }& \DIFadd{-0.3216 }& {[}\DIFadd{-0.6867, 0.04341}{]} & \DIFadd{p=0.0774 }& \DIFadd{q=0.112 }\\
\DIFadd{Train QE }& \DIFadd{breast\_cancer }& \DIFadd{0.9588 }& {[}\DIFadd{-0.4954, 2.413}{]} & \DIFadd{p=0.17 }&
\DIFadd{q=0.298 }& \DIFadd{-0.263 }& {[}\DIFadd{-1.543, 1.016}{]} & \DIFadd{p=0.653 }& \DIFadd{q=0.722 }\\
\DIFadd{Train QE }& \DIFadd{wine }& \DIFadd{-17.32 }& {[}\DIFadd{-23.63, -11.02}{]} & \DIFadd{p=0.000156 }& \DIFadd{q=0.00131
}& \DIFadd{-22.4 }& {[}\DIFadd{-33.21, -11.6}{]} & \DIFadd{p=0.00114 }& \DIFadd{q=0.00341 }\\
\DIFadd{Train QE }& \DIFadd{iris }& \DIFadd{-5.888 }& {[}\DIFadd{-28.17, 16.4}{]} & \DIFadd{p=0.565 }& \DIFadd{q=0.641 }&
\DIFadd{-25.13 }& {[}\DIFadd{-36.77, -13.48}{]} & \DIFadd{p=0.000872 }& \DIFadd{q=0.00299 }\\
\DIFadd{Train QE }& \DIFadd{digits }& \DIFadd{-0.2587 }& {[}\DIFadd{-0.7176, 0.2002}{]} & \DIFadd{p=0.234 }& \DIFadd{q=0.351
}& \DIFadd{-0.3143 }& {[}\DIFadd{-0.6485, 0.01992}{]} & \DIFadd{p=0.0623 }& \DIFadd{q=0.0969 }\\
\DIFadd{Train QE }& \DIFadd{olivetti\_faces }& \DIFadd{-1.791 }& {[}\DIFadd{-4.796, 1.215}{]} & \DIFadd{p=0.211 }&
\DIFadd{q=0.336 }& \DIFadd{-1.553 }& {[}\DIFadd{-4.04, 0.9338}{]} & \DIFadd{p=0.191 }& \DIFadd{q=0.259 }\\
\DIFadd{Train QE }& \DIFadd{diabetes }& \DIFadd{0.2611 }& {[}\DIFadd{-1.533, 2.055}{]} & \DIFadd{p=0.75 }& \DIFadd{q=0.787 }&
\DIFadd{-0.8799 }& {[}\DIFadd{-2.463, 0.7033}{]} & \DIFadd{p=0.24 }& \DIFadd{q=0.306 }\\
\DIFadd{Train QE }& \DIFadd{california\_housing }& \DIFadd{-0.3185 }& {[}\DIFadd{-0.4766, -0.1604}{]} &
\DIFadd{p=0.00137 }& \DIFadd{q=0.00523 }& \DIFadd{-0.2556 }& {[}\DIFadd{-0.3834, -0.1278}{]} & \DIFadd{p=0.00144 }&
\DIFadd{q=0.00355 }\\
\DIFadd{Train QE }& \DIFadd{covertype }& \DIFadd{-0.3468 }& {[}\DIFadd{-1.407, 0.7137}{]} & \DIFadd{p=0.478 }&
\DIFadd{q=0.574 }& \DIFadd{-0.1341 }& {[}\DIFadd{-0.8722, 0.6039}{]} & \DIFadd{p=0.691 }& \DIFadd{q=0.744 }\\
\DIFadd{Train QE }& \DIFadd{kddcup99 }& \DIFadd{-5.804 }& {[}\DIFadd{-6.474, -5.134}{]} & \DIFadd{p=1.09e-08 }&
\DIFadd{q=4.56e-07 }& \DIFadd{-4.183 }& {[}\DIFadd{-5.557, -2.809}{]} & \DIFadd{p=7.18e-05 }& \DIFadd{q=0.000565 }\\
\DIFadd{Train QE }& \DIFadd{OVERALL }& \DIFadd{-2.241 }& {[}\DIFadd{-3.842, -0.6405}{]} & \DIFadd{p=0.0064 }& \DIFadd{q=NA }&
\DIFadd{-4.033 }& {[}\DIFadd{-5.707, -2.36}{]} & \DIFadd{p=4.69e-06 }& \DIFadd{q=NA }\\
\end{longtable}
}

\DIFaddend \textbf{Supplementary Table S8. Figure  \DIFaddbegin \DIFadd{15 }\DIFaddend deployment comparison percent
summary for tuned FloatSOM RNG versus  \DIFaddbegin \DIFadd{untuned }\DIFaddend hexagonal XPySOM across
\(QE_B\), \(QE_H\), and \(QE_T\).} Rows list per-dataset and
\texttt{GLOBAL\_OVERALL} entries with the plotted median percent change
and 95\% confidence interval.

\DIFaddbegin {\def\LTcaptype{none} %DIF >  do not increment counter
\begin{longtable}[]{@{}llll@{}}
\toprule\noalign{}
\DIFadd{metric }& \DIFadd{dataset }& \DIFadd{median \% change }& \DIFadd{95\% CI }\\
\midrule\noalign{}
\endhead
\bottomrule\noalign{}
\endlastfoot
\DIFadd{QE\_B }& \DIFadd{blobs }& \DIFadd{5.4588 }& {[}\DIFadd{3.8282, 7.0895}{]} \\
\DIFadd{QE\_B }& \DIFadd{circles }& \DIFadd{5.6542 }& {[}\DIFadd{4.552, 6.7563}{]} \\
\DIFadd{QE\_B }& \DIFadd{moons }& \DIFadd{4.4217 }& {[}\DIFadd{2.6294, 6.2141}{]} \\
\DIFadd{QE\_B }& \DIFadd{s\_curve }& \DIFadd{10.0599 }& {[}\DIFadd{9.2501, 10.8696}{]} \\
\DIFadd{QE\_B }& \DIFadd{swiss\_roll }& \DIFadd{19.7271 }& {[}\DIFadd{18.4881, 20.9661}{]} \\
\DIFadd{QE\_B }& \DIFadd{breast\_cancer }& \DIFadd{7.2273 }& {[}\DIFadd{6.3913, 8.0633}{]} \\
\DIFadd{QE\_B }& \DIFadd{wine }& \DIFadd{20.5163 }& {[}\DIFadd{19.3744, 21.6583}{]} \\
\DIFadd{QE\_B }& \DIFadd{iris }& \DIFadd{23.2197 }& {[}\DIFadd{20.569, 25.8704}{]} \\
\DIFadd{QE\_B }& \DIFadd{digits }& \DIFadd{8.3021 }& {[}\DIFadd{7.8283, 8.7758}{]} \\
\DIFadd{QE\_B }& \DIFadd{olivetti\_faces }& \DIFadd{10.9797 }& {[}\DIFadd{10.0517, 11.9077}{]} \\
\DIFadd{QE\_B }& \DIFadd{diabetes }& \DIFadd{9.2085 }& {[}\DIFadd{8.1874, 10.2296}{]} \\
\DIFadd{QE\_B }& \DIFadd{california\_housing }& \DIFadd{9.4344 }& {[}\DIFadd{9.0497, 9.8191}{]} \\
\DIFadd{QE\_B }& \DIFadd{covertype }& \DIFadd{30.3468 }& {[}\DIFadd{29.1115, 31.5821}{]} \\
\DIFadd{QE\_B }& \DIFadd{kddcup99 }& \DIFadd{38.2338 }& {[}\DIFadd{35.8971, 40.5705}{]} \\
\DIFadd{QE\_B }& \DIFadd{GLOBAL\_OVERALL }& \DIFadd{14.485 }& {[}\DIFadd{12.7836, 16.1865}{]} \\
\DIFadd{QE\_H }& \DIFadd{blobs }& \DIFadd{4.8943 }& {[}\DIFadd{3.3372, 6.4515}{]} \\
\DIFadd{QE\_H }& \DIFadd{circles }& \DIFadd{4.98 }& {[}\DIFadd{3.8583, 6.1016}{]} \\
\DIFadd{QE\_H }& \DIFadd{moons }& \DIFadd{3.8767 }& {[}\DIFadd{1.9564, 5.797}{]} \\
\DIFadd{QE\_H }& \DIFadd{s\_curve }& \DIFadd{9.4896 }& {[}\DIFadd{8.5416, 10.4377}{]} \\
\DIFadd{QE\_H }& \DIFadd{swiss\_roll }& \DIFadd{19.2924 }& {[}\DIFadd{17.9583, 20.6264}{]} \\
\DIFadd{QE\_H }& \DIFadd{breast\_cancer }& \DIFadd{0.1143 }& {[}\DIFadd{-0.6627, 0.8913}{]} \\
\DIFadd{QE\_H }& \DIFadd{wine }& \DIFadd{-1.8427 }& {[}\DIFadd{-3.2635, -0.4218}{]} \\
\DIFadd{QE\_H }& \DIFadd{iris }& \DIFadd{0.4853 }& {[}\DIFadd{-2.8955, 3.8661}{]} \\
\DIFadd{QE\_H }& \DIFadd{digits }& \DIFadd{4.9934 }& {[}\DIFadd{4.5072, 5.4796}{]} \\
\DIFadd{QE\_H }& \DIFadd{olivetti\_faces }& \DIFadd{2.6384 }& {[}\DIFadd{1.2928, 3.9839}{]} \\
\DIFadd{QE\_H }& \DIFadd{diabetes }& \DIFadd{1.3763 }& {[}\DIFadd{-0.3528, 3.1055}{]} \\
\DIFadd{QE\_H }& \DIFadd{california\_housing }& \DIFadd{8.5235 }& {[}\DIFadd{8.0055, 9.0414}{]} \\
\DIFadd{QE\_H }& \DIFadd{covertype }& \DIFadd{30.3196 }& {[}\DIFadd{29.0984, 31.5408}{]} \\
\DIFadd{QE\_H }& \DIFadd{kddcup99 }& \DIFadd{37.9291 }& {[}\DIFadd{35.5645, 40.2937}{]} \\
\DIFadd{QE\_H }& \DIFadd{GLOBAL\_OVERALL }& \DIFadd{9.0765 }& {[}\DIFadd{7.1203, 11.0326}{]} \\
\DIFadd{QE\_T }& \DIFadd{blobs }& \DIFadd{6.0226 }& {[}\DIFadd{4.2924, 7.7527}{]} \\
\DIFadd{QE\_T }& \DIFadd{circles }& \DIFadd{6.3297 }& {[}\DIFadd{5.1892, 7.4702}{]} \\
\DIFadd{QE\_T }& \DIFadd{moons }& \DIFadd{4.968 }& {[}\DIFadd{3.2744, 6.6616}{]} \\
\DIFadd{QE\_T }& \DIFadd{s\_curve }& \DIFadd{10.6309 }& {[}\DIFadd{9.9339, 11.3279}{]} \\
\DIFadd{QE\_T }& \DIFadd{swiss\_roll }& \DIFadd{20.1631 }& {[}\DIFadd{18.9915, 21.3347}{]} \\
\DIFadd{QE\_T }& \DIFadd{breast\_cancer }& \DIFadd{15.8475 }& {[}\DIFadd{14.5024, 17.1926}{]} \\
\DIFadd{QE\_T }& \DIFadd{wine }& \DIFadd{56.1875 }& {[}\DIFadd{54.54, 57.8351}{]} \\
\DIFadd{QE\_T }& \DIFadd{iris }& \DIFadd{63.2439 }& {[}\DIFadd{60.2553, 66.2326}{]} \\
\DIFadd{QE\_T }& \DIFadd{digits }& \DIFadd{11.8094 }& {[}\DIFadd{11.1554, 12.4634}{]} \\
\DIFadd{QE\_T }& \DIFadd{olivetti\_faces }& \DIFadd{21.029 }& {[}\DIFadd{19.7281, 22.3298}{]} \\
\DIFadd{QE\_T }& \DIFadd{diabetes }& \DIFadd{18.9602 }& {[}\DIFadd{18.1282, 19.7921}{]} \\
\DIFadd{QE\_T }& \DIFadd{california\_housing }& \DIFadd{10.3513 }& {[}\DIFadd{9.7384, 10.9642}{]} \\
\DIFadd{QE\_T }& \DIFadd{covertype }& \DIFadd{30.3738 }& {[}\DIFadd{29.1225, 31.6251}{]} \\
\DIFadd{QE\_T }& \DIFadd{kddcup99 }& \DIFadd{38.5352 }& {[}\DIFadd{36.2102, 40.8602}{]} \\
\DIFadd{QE\_T }& \DIFadd{GLOBAL\_OVERALL }& \DIFadd{22.4609 }& {[}\DIFadd{19.4613, 25.4605}{]} \\
\end{longtable}
}

\DIFaddend \textbf{Supplementary Table S9. Supplementary Figure  \DIFaddbegin \DIFadd{S7 }\DIFaddend deployment
comparison percent summary for tuned FloatSOM hexagonal versus  \DIFaddbegin \DIFadd{untuned
}\DIFaddend hexagonal XPySOM across \(QE_B\), \(QE_H\), and \(QE_T\).} Rows list
per-dataset and \texttt{GLOBAL\_OVERALL} entries with the plotted median
percent change and 95\% confidence interval.

\DIFaddbegin {\def\LTcaptype{none} %DIF >  do not increment counter
\begin{longtable}[]{@{}llll@{}}
\toprule\noalign{}
\DIFadd{metric }& \DIFadd{dataset }& \DIFadd{median \% change }& \DIFadd{95\% CI }\\
\midrule\noalign{}
\endhead
\bottomrule\noalign{}
\endlastfoot
\DIFadd{QE\_B }& \DIFadd{blobs }& \DIFadd{4.5418 }& {[}\DIFadd{3.0375, 6.0462}{]} \\
\DIFadd{QE\_B }& \DIFadd{circles }& \DIFadd{4.5998 }& {[}\DIFadd{3.6022, 5.5974}{]} \\
\DIFadd{QE\_B }& \DIFadd{moons }& \DIFadd{3.5895 }& {[}\DIFadd{2.051, 5.128}{]} \\
\DIFadd{QE\_B }& \DIFadd{s\_curve }& \DIFadd{9.7727 }& {[}\DIFadd{9.0814, 10.464}{]} \\
\DIFadd{QE\_B }& \DIFadd{swiss\_roll }& \DIFadd{17.6702 }& {[}\DIFadd{16.7526, 18.5878}{]} \\
\DIFadd{QE\_B }& \DIFadd{breast\_cancer }& \DIFadd{6.399 }& {[}\DIFadd{5.5895, 7.2084}{]} \\
\DIFadd{QE\_B }& \DIFadd{wine }& \DIFadd{15.0518 }& {[}\DIFadd{13.113, 16.9905}{]} \\
\DIFadd{QE\_B }& \DIFadd{iris }& \DIFadd{19.6158 }& {[}\DIFadd{17.7567, 21.475}{]} \\
\DIFadd{QE\_B }& \DIFadd{digits }& \DIFadd{7.7568 }& {[}\DIFadd{7.3648, 8.1489}{]} \\
\DIFadd{QE\_B }& \DIFadd{olivetti\_faces }& \DIFadd{10.367 }& {[}\DIFadd{9.5746, 11.1594}{]} \\
\DIFadd{QE\_B }& \DIFadd{diabetes }& \DIFadd{9.133 }& {[}\DIFadd{8.2954, 9.9707}{]} \\
\DIFadd{QE\_B }& \DIFadd{california\_housing }& \DIFadd{8.9193 }& {[}\DIFadd{8.4319, 9.4067}{]} \\
\DIFadd{QE\_B }& \DIFadd{covertype }& \DIFadd{29.6621 }& {[}\DIFadd{28.3242, 31}{]} \\
\DIFadd{QE\_B }& \DIFadd{kddcup99 }& \DIFadd{35.5053 }& {[}\DIFadd{34.0328, 36.9779}{]} \\
\DIFadd{QE\_B }& \DIFadd{GLOBAL\_OVERALL }& \DIFadd{13.0417 }& {[}\DIFadd{11.4631, 14.6204}{]} \\
\DIFadd{QE\_H }& \DIFadd{blobs }& \DIFadd{3.8869 }& {[}\DIFadd{2.4806, 5.2931}{]} \\
\DIFadd{QE\_H }& \DIFadd{circles }& \DIFadd{4.1388 }& {[}\DIFadd{3.1274, 5.1503}{]} \\
\DIFadd{QE\_H }& \DIFadd{moons }& \DIFadd{2.9919 }& {[}\DIFadd{1.3076, 4.6761}{]} \\
\DIFadd{QE\_H }& \DIFadd{s\_curve }& \DIFadd{9.2419 }& {[}\DIFadd{8.4744, 10.0095}{]} \\
\DIFadd{QE\_H }& \DIFadd{swiss\_roll }& \DIFadd{17.3207 }& {[}\DIFadd{16.3909, 18.2504}{]} \\
\DIFadd{QE\_H }& \DIFadd{breast\_cancer }& \DIFadd{-0.997 }& {[}\DIFadd{-1.683, -0.311}{]} \\
\DIFadd{QE\_H }& \DIFadd{wine }& \DIFadd{-2.8868 }& {[}\DIFadd{-4.4124, -1.3613}{]} \\
\DIFadd{QE\_H }& \DIFadd{iris }& \DIFadd{-0.4904 }& {[}\DIFadd{-2.1865, 1.2058}{]} \\
\DIFadd{QE\_H }& \DIFadd{digits }& \DIFadd{4.5103 }& {[}\DIFadd{4.0172, 5.0034}{]} \\
\DIFadd{QE\_H }& \DIFadd{olivetti\_faces }& \DIFadd{1.9861 }& {[}\DIFadd{0.769, 3.2032}{]} \\
\DIFadd{QE\_H }& \DIFadd{diabetes }& \DIFadd{2.245 }& {[}\DIFadd{1.1198, 3.3702}{]} \\
\DIFadd{QE\_H }& \DIFadd{california\_housing }& \DIFadd{8.0301 }& {[}\DIFadd{7.4144, 8.6458}{]} \\
\DIFadd{QE\_H }& \DIFadd{covertype }& \DIFadd{29.6665 }& {[}\DIFadd{28.3026, 31.0305}{]} \\
\DIFadd{QE\_H }& \DIFadd{kddcup99 }& \DIFadd{35.1253 }& {[}\DIFadd{33.6747, 36.5759}{]} \\
\DIFadd{QE\_H }& \DIFadd{GLOBAL\_OVERALL }& \DIFadd{8.1978 }& {[}\DIFadd{6.3289, 10.0667}{]} \\
\DIFadd{QE\_T }& \DIFadd{blobs }& \DIFadd{5.1961 }& {[}\DIFadd{3.5689, 6.8233}{]} \\
\DIFadd{QE\_T }& \DIFadd{circles }& \DIFadd{5.0606 }& {[}\DIFadd{4.0085, 6.1127}{]} \\
\DIFadd{QE\_T }& \DIFadd{moons }& \DIFadd{4.1898 }& {[}\DIFadd{2.7833, 5.5963}{]} \\
\DIFadd{QE\_T }& \DIFadd{s\_curve }& \DIFadd{10.3044 }& {[}\DIFadd{9.6672, 10.9416}{]} \\
\DIFadd{QE\_T }& \DIFadd{swiss\_roll }& \DIFadd{18.021 }& {[}\DIFadd{17.101, 18.9411}{]} \\
\DIFadd{QE\_T }& \DIFadd{breast\_cancer }& \DIFadd{15.3534 }& {[}\DIFadd{14.1184, 16.5884}{]} \\
\DIFadd{QE\_T }& \DIFadd{wine }& \DIFadd{43.6818 }& {[}\DIFadd{39.113, 48.2507}{]} \\
\DIFadd{QE\_T }& \DIFadd{iris }& \DIFadd{55.1557 }& {[}\DIFadd{50.9804, 59.331}{]} \\
\DIFadd{QE\_T }& \DIFadd{digits }& \DIFadd{11.2012 }& {[}\DIFadd{10.7183, 11.6841}{]} \\
\DIFadd{QE\_T }& \DIFadd{olivetti\_faces }& \DIFadd{20.4609 }& {[}\DIFadd{19.2364, 21.6853}{]} \\
\DIFadd{QE\_T }& \DIFadd{diabetes }& \DIFadd{17.7191 }& {[}\DIFadd{16.8369, 18.6013}{]} \\
\DIFadd{QE\_T }& \DIFadd{california\_housing }& \DIFadd{9.8167 }& {[}\DIFadd{9.2321, 10.4012}{]} \\
\DIFadd{QE\_T }& \DIFadd{covertype }& \DIFadd{29.6576 }& {[}\DIFadd{28.3443, 30.9708}{]} \\
\DIFadd{QE\_T }& \DIFadd{kddcup99 }& \DIFadd{35.8819 }& {[}\DIFadd{34.3431, 37.4206}{]} \\
\DIFadd{QE\_T }& \DIFadd{GLOBAL\_OVERALL }& \DIFadd{20.1214 }& {[}\DIFadd{17.5722, 22.6707}{]} \\
\end{longtable}
}

\DIFaddend \textbf{Supplementary Table S10. Supplementary Figure  \DIFaddbegin \DIFadd{S8 }\DIFaddend deployment
comparison percent summary for tuned FloatSOM MST versus  \DIFaddbegin \DIFadd{untuned
}\DIFaddend hexagonal XPySOM across \(QE_B\), \(QE_H\), and \(QE_T\).} Rows list
per-dataset and \texttt{GLOBAL\_OVERALL} entries with the plotted median
percent change and 95\% confidence interval.

\DIFaddbegin {\def\LTcaptype{none} %DIF >  do not increment counter
\begin{longtable}[]{@{}llll@{}}
\toprule\noalign{}
\DIFadd{metric }& \DIFadd{dataset }& \DIFadd{median \% change }& \DIFadd{95\% CI }\\
\midrule\noalign{}
\endhead
\bottomrule\noalign{}
\endlastfoot
\DIFadd{QE\_B }& \DIFadd{blobs }& \DIFadd{5.2243 }& {[}\DIFadd{3.6173, 6.8314}{]} \\
\DIFadd{QE\_B }& \DIFadd{circles }& \DIFadd{5.5236 }& {[}\DIFadd{4.4413, 6.6059}{]} \\
\DIFadd{QE\_B }& \DIFadd{moons }& \DIFadd{4.2266 }& {[}\DIFadd{2.4056, 6.0477}{]} \\
\DIFadd{QE\_B }& \DIFadd{s\_curve }& \DIFadd{9.968 }& {[}\DIFadd{9.2244, 10.7116}{]} \\
\DIFadd{QE\_B }& \DIFadd{swiss\_roll }& \DIFadd{19.5015 }& {[}\DIFadd{18.1673, 20.8356}{]} \\
\DIFadd{QE\_B }& \DIFadd{breast\_cancer }& \DIFadd{7.2275 }& {[}\DIFadd{6.4917, 7.9633}{]} \\
\DIFadd{QE\_B }& \DIFadd{wine }& \DIFadd{18.0952 }& {[}\DIFadd{16.4439, 19.7464}{]} \\
\DIFadd{QE\_B }& \DIFadd{iris }& \DIFadd{17.3842 }& {[}\DIFadd{13.8051, 20.9634}{]} \\
\DIFadd{QE\_B }& \DIFadd{digits }& \DIFadd{8.5168 }& {[}\DIFadd{7.9287, 9.105}{]} \\
\DIFadd{QE\_B }& \DIFadd{olivetti\_faces }& \DIFadd{11.9472 }& {[}\DIFadd{10.795, 13.0994}{]} \\
\DIFadd{QE\_B }& \DIFadd{diabetes }& \DIFadd{9.2961 }& {[}\DIFadd{8.2097, 10.3825}{]} \\
\DIFadd{QE\_B }& \DIFadd{california\_housing }& \DIFadd{9.474 }& {[}\DIFadd{9.0518, 9.8963}{]} \\
\DIFadd{QE\_B }& \DIFadd{covertype }& \DIFadd{32.2987 }& {[}\DIFadd{30.9976, 33.5997}{]} \\
\DIFadd{QE\_B }& \DIFadd{kddcup99 }& \DIFadd{40.8705 }& {[}\DIFadd{39.5821, 42.159}{]} \\
\DIFadd{QE\_B }& \DIFadd{GLOBAL\_OVERALL }& \DIFadd{14.2539 }& {[}\DIFadd{12.4871, 16.0207}{]} \\
\DIFadd{QE\_H }& \DIFadd{blobs }& \DIFadd{4.5705 }& {[}\DIFadd{3.0306, 6.1104}{]} \\
\DIFadd{QE\_H }& \DIFadd{circles }& \DIFadd{4.8064 }& {[}\DIFadd{3.6639, 5.9489}{]} \\
\DIFadd{QE\_H }& \DIFadd{moons }& \DIFadd{3.768 }& {[}\DIFadd{1.7919, 5.744}{]} \\
\DIFadd{QE\_H }& \DIFadd{s\_curve }& \DIFadd{9.4507 }& {[}\DIFadd{8.6251, 10.2763}{]} \\
\DIFadd{QE\_H }& \DIFadd{swiss\_roll }& \DIFadd{19.0489 }& {[}\DIFadd{17.593, 20.5048}{]} \\
\DIFadd{QE\_H }& \DIFadd{breast\_cancer }& \DIFadd{-0.7772 }& {[}\DIFadd{-1.3919, -0.1624}{]} \\
\DIFadd{QE\_H }& \DIFadd{wine }& \DIFadd{-2.2046 }& {[}\DIFadd{-3.9272, -0.4821}{]} \\
\DIFadd{QE\_H }& \DIFadd{iris }& \DIFadd{-0.872 }& {[}\DIFadd{-3.7238, 1.9798}{]} \\
\DIFadd{QE\_H }& \DIFadd{digits }& \DIFadd{5.2512 }& {[}\DIFadd{4.6436, 5.8587}{]} \\
\DIFadd{QE\_H }& \DIFadd{olivetti\_faces }& \DIFadd{2.5542 }& {[}\DIFadd{0.9211, 4.1873}{]} \\
\DIFadd{QE\_H }& \DIFadd{diabetes }& \DIFadd{1.059 }& {[}\DIFadd{-0.5135, 2.6316}{]} \\
\DIFadd{QE\_H }& \DIFadd{california\_housing }& \DIFadd{8.4673 }& {[}\DIFadd{7.9248, 9.0099}{]} \\
\DIFadd{QE\_H }& \DIFadd{covertype }& \DIFadd{32.2708 }& {[}\DIFadd{30.9787, 33.5629}{]} \\
\DIFadd{QE\_H }& \DIFadd{kddcup99 }& \DIFadd{40.3684 }& {[}\DIFadd{39.0453, 41.6915}{]} \\
\DIFadd{QE\_H }& \DIFadd{GLOBAL\_OVERALL }& \DIFadd{9.1258 }& {[}\DIFadd{7.0316, 11.2201}{]} \\
\DIFadd{QE\_T }& \DIFadd{blobs }& \DIFadd{5.8778 }& {[}\DIFadd{4.178, 7.5776}{]} \\
\DIFadd{QE\_T }& \DIFadd{circles }& \DIFadd{6.2415 }& {[}\DIFadd{5.1655, 7.3175}{]} \\
\DIFadd{QE\_T }& \DIFadd{moons }& \DIFadd{4.6869 }& {[}\DIFadd{3.0023, 6.3715}{]} \\
\DIFadd{QE\_T }& \DIFadd{s\_curve }& \DIFadd{10.486 }& {[}\DIFadd{9.7854, 11.1866}{]} \\
\DIFadd{QE\_T }& \DIFadd{swiss\_roll }& \DIFadd{19.9554 }& {[}\DIFadd{18.7303, 21.1804}{]} \\
\DIFadd{QE\_T }& \DIFadd{breast\_cancer }& \DIFadd{16.9227 }& {[}\DIFadd{15.7815, 18.0639}{]} \\
\DIFadd{QE\_T }& \DIFadd{wine }& \DIFadd{50.5129 }& {[}\DIFadd{47.2136, 53.8122}{]} \\
\DIFadd{QE\_T }& \DIFadd{iris }& \DIFadd{49.5872 }& {[}\DIFadd{43.4267, 55.7476}{]} \\
\DIFadd{QE\_T }& \DIFadd{digits }& \DIFadd{11.9794 }& {[}\DIFadd{11.277, 12.6817}{]} \\
\DIFadd{QE\_T }& \DIFadd{olivetti\_faces }& \DIFadd{23.2531 }& {[}\DIFadd{21.9339, 24.5724}{]} \\
\DIFadd{QE\_T }& \DIFadd{diabetes }& \DIFadd{19.5575 }& {[}\DIFadd{18.2819, 20.8331}{]} \\
\DIFadd{QE\_T }& \DIFadd{california\_housing }& \DIFadd{10.4893 }& {[}\DIFadd{9.9385, 11.04}{]} \\
\DIFadd{QE\_T }& \DIFadd{covertype }& \DIFadd{32.3264 }& {[}\DIFadd{31.0143, 33.6385}{]} \\
\DIFadd{QE\_T }& \DIFadd{kddcup99 }& \DIFadd{41.3715 }& {[}\DIFadd{40.0794, 42.6636}{]} \\
\DIFadd{QE\_T }& \DIFadd{GLOBAL\_OVERALL }& \DIFadd{21.6605 }& {[}\DIFadd{19.0539, 24.2672}{]} \\
\end{longtable}
}

\DIFaddend \textbf{Supplementary Table S11. Figure  \DIFaddbegin \DIFadd{14 }\DIFaddend topology runtime summary at
the largest common 8-GPU axis value for the dimension, sample, and grid
size scaling workloads.} Rows report the plotted 8-GPU mean runtimes for
hexagonal, MST, and RNG, together with the fastest and slowest topology
at that axis value and the maximum pairwise runtime spread.

\DIFaddbegin {\def\LTcaptype{none} %DIF >  do not increment counter
\begin{longtable}[]{@{}
  >{\raggedright\arraybackslash}p{(\linewidth - 14\tabcolsep) * \real{0.1250}}
  >{\raggedright\arraybackslash}p{(\linewidth - 14\tabcolsep) * \real{0.1250}}
  >{\raggedright\arraybackslash}p{(\linewidth - 14\tabcolsep) * \real{0.1250}}
  >{\raggedright\arraybackslash}p{(\linewidth - 14\tabcolsep) * \real{0.1250}}
  >{\raggedright\arraybackslash}p{(\linewidth - 14\tabcolsep) * \real{0.1250}}
  >{\raggedright\arraybackslash}p{(\linewidth - 14\tabcolsep) * \real{0.1250}}
  >{\raggedright\arraybackslash}p{(\linewidth - 14\tabcolsep) * \real{0.1250}}
  >{\raggedright\arraybackslash}p{(\linewidth - 14\tabcolsep) * \real{0.1250}}@{}}
\toprule\noalign{}
\begin{minipage}[b]{\linewidth}\raggedright
\DIFadd{axis
}\end{minipage} & \begin{minipage}[b]{\linewidth}\raggedright
\DIFadd{axis value
}\end{minipage} & \begin{minipage}[b]{\linewidth}\raggedright
\DIFadd{hexagonal s
}\end{minipage} & \begin{minipage}[b]{\linewidth}\raggedright
\DIFadd{MST s
}\end{minipage} & \begin{minipage}[b]{\linewidth}\raggedright
\DIFadd{RNG s
}\end{minipage} & \begin{minipage}[b]{\linewidth}\raggedright
\DIFadd{fastest
}\end{minipage} & \begin{minipage}[b]{\linewidth}\raggedright
\DIFadd{slowest
}\end{minipage} & \begin{minipage}[b]{\linewidth}\raggedright
\DIFadd{spread \%
}\end{minipage} \\
\midrule\noalign{}
\endhead
\bottomrule\noalign{}
\endlastfoot
\DIFadd{Dimension Scaling }& \DIFadd{5000 }& \DIFadd{404.17 }& \DIFadd{423.15 }& \DIFadd{409.94 }& \DIFadd{hexagonal }& \DIFadd{mst }&
\DIFadd{4.7 }\\
\DIFadd{Sample Scaling }& \DIFadd{1e+09 }& \DIFadd{363.58 }& \DIFadd{375.44 }& \DIFadd{369.41 }& \DIFadd{hexagonal }& \DIFadd{mst }&
\DIFadd{3.26 }\\
\DIFadd{Grid-Size Scaling }& \DIFadd{64 }& \DIFadd{32.54 }& \DIFadd{266.45 }& \DIFadd{880.83 }& \DIFadd{hexagonal }& \DIFadd{rng }&
\DIFadd{2606.61 }\\
\end{longtable}
}

\DIFadd{\textbf{Supplementary Table S12. Matched topology diagnostic paired
summaries across tuned and untuned profiles.} Rows report pooled paired
effects from the final matched random-seed topology diagnostic
benchmark. \texttt{Effect\ favoring\ comparator} is oriented so positive
values favor the comparator after applying each metric's directionality;
lower is better for QE, MTR, and dead-node fraction, while higher is
better for node utilization. Because these are pre-specified pooled
paired summaries rather than dataset-level test families, the table
reports raw paired-test p-values. \texttt{comp/ref/tie} gives the number
of comparator-favoring, reference-favoring, and tied matched pairs.
Metric suffixes denote holdout (\texttt{\_H}), train (\texttt{\_T}), and
balanced train-holdout (\texttt{\_B}) summaries. The tuned profile uses
the deployable tuned configurations reported in the main text.
}

{\def\LTcaptype{none} %DIF >  do not increment counter
\begin{longtable}[]{@{}
  >{\raggedright\arraybackslash}p{(\linewidth - 16\tabcolsep) * \real{0.0968}}
  >{\raggedright\arraybackslash}p{(\linewidth - 16\tabcolsep) * \real{0.0968}}
  >{\raggedright\arraybackslash}p{(\linewidth - 16\tabcolsep) * \real{0.0968}}
  >{\raggedleft\arraybackslash}p{(\linewidth - 16\tabcolsep) * \real{0.1290}}
  >{\raggedleft\arraybackslash}p{(\linewidth - 16\tabcolsep) * \real{0.1290}}
  >{\raggedright\arraybackslash}p{(\linewidth - 16\tabcolsep) * \real{0.0968}}
  >{\raggedleft\arraybackslash}p{(\linewidth - 16\tabcolsep) * \real{0.1290}}
  >{\raggedleft\arraybackslash}p{(\linewidth - 16\tabcolsep) * \real{0.1290}}
  >{\raggedright\arraybackslash}p{(\linewidth - 16\tabcolsep) * \real{0.0968}}@{}}
\toprule\noalign{}
\begin{minipage}[b]{\linewidth}\raggedright
\DIFadd{Comparison
}\end{minipage} & \begin{minipage}[b]{\linewidth}\raggedright
\DIFadd{Metric
}\end{minipage} & \begin{minipage}[b]{\linewidth}\raggedright
\DIFadd{Better
}\end{minipage} & \begin{minipage}[b]{\linewidth}\raggedleft
\DIFadd{n
}\end{minipage} & \begin{minipage}[b]{\linewidth}\raggedleft
\DIFadd{Effect favoring comparator
}\end{minipage} & \begin{minipage}[b]{\linewidth}\raggedright
\DIFadd{95\% CI
}\end{minipage} & \begin{minipage}[b]{\linewidth}\raggedleft
\DIFadd{dz
}\end{minipage} & \begin{minipage}[b]{\linewidth}\raggedleft
\DIFadd{raw p
}\end{minipage} & \begin{minipage}[b]{\linewidth}\raggedright
\DIFadd{comp/ref/tie
}\end{minipage} \\
\midrule\noalign{}
\endhead
\bottomrule\noalign{}
\endlastfoot
\DIFadd{Untuned hex vs untuned MST }& \DIFadd{QE\_H }& \DIFadd{lower }& \DIFadd{280 }& \DIFadd{0.1523 }& {[}\DIFadd{0.1110,
0.1937}{]} & \DIFadd{0.4330 }& \DIFadd{4.22e-12 }& \DIFadd{185/95/0 }\\
\DIFadd{Untuned hex vs untuned MST }& \DIFadd{QE\_T }& \DIFadd{lower }& \DIFadd{280 }& \DIFadd{0.3862 }& {[}\DIFadd{0.2780,
0.4944}{]} & \DIFadd{0.4199 }& \DIFadd{1.63e-11 }& \DIFadd{226/54/0 }\\
\DIFadd{Untuned hex vs untuned MST }& \DIFadd{QE\_B }& \DIFadd{lower }& \DIFadd{280 }& \DIFadd{0.2693 }& {[}\DIFadd{0.1975,
0.3410}{]} & \DIFadd{0.4415 }& \DIFadd{1.73e-12 }& \DIFadd{222/58/0 }\\
\DIFadd{Untuned hex vs untuned MST }& \DIFadd{MTR\_H }& \DIFadd{lower }& \DIFadd{280 }& \DIFadd{-0.6164 }&
{[}\DIFadd{-0.9036, -0.3292}{]} & \DIFadd{-0.2525 }& \DIFadd{3.24e-05 }& \DIFadd{125/155/0 }\\
\DIFadd{Untuned hex vs untuned MST }& \DIFadd{MTR\_T }& \DIFadd{lower }& \DIFadd{280 }& \DIFadd{0.2075 }& {[}\DIFadd{-0.0155,
0.4304}{]} & \DIFadd{0.1095 }& \DIFadd{0.068 }& \DIFadd{167/113/0 }\\
\DIFadd{Untuned hex vs untuned MST }& \DIFadd{MTR\_B }& \DIFadd{lower }& \DIFadd{280 }& \DIFadd{-0.2045 }&
{[}\DIFadd{-0.4382, 0.0293}{]} & \DIFadd{-0.1029 }& \DIFadd{0.086 }& \DIFadd{139/141/0 }\\
\DIFadd{Untuned hex vs untuned MST }& \DIFadd{Util\_H }& \DIFadd{higher }& \DIFadd{280 }& \DIFadd{0.0187 }&
{[}\DIFadd{0.0115, 0.0260}{]} & \DIFadd{0.3056 }& \DIFadd{5.88e-07 }& \DIFadd{110/73/97 }\\
\DIFadd{Untuned hex vs untuned MST }& \DIFadd{Util\_T }& \DIFadd{higher }& \DIFadd{280 }& \DIFadd{0.0414 }&
{[}\DIFadd{0.0345, 0.0483}{]} & \DIFadd{0.7049 }& \DIFadd{2.55e-26 }& \DIFadd{147/17/116 }\\
\DIFadd{Untuned hex vs untuned MST }& \DIFadd{Util\_B }& \DIFadd{higher }& \DIFadd{280 }& \DIFadd{0.0301 }&
{[}\DIFadd{0.0237, 0.0364}{]} & \DIFadd{0.5534 }& \DIFadd{5.62e-18 }& \DIFadd{144/46/90 }\\
\DIFadd{Untuned hex vs untuned MST }& \DIFadd{Dead\_H }& \DIFadd{lower }& \DIFadd{280 }& \DIFadd{0.0188 }& {[}\DIFadd{0.0115,
0.0260}{]} & \DIFadd{0.3056 }& \DIFadd{5.88e-07 }& \DIFadd{110/73/97 }\\
\DIFadd{Untuned hex vs untuned MST }& \DIFadd{Dead\_T }& \DIFadd{lower }& \DIFadd{280 }& \DIFadd{0.0414 }& {[}\DIFadd{0.0345,
0.0483}{]} & \DIFadd{0.7049 }& \DIFadd{2.55e-26 }& \DIFadd{147/17/116 }\\
\DIFadd{Untuned hex vs untuned MST }& \DIFadd{Dead\_B }& \DIFadd{lower }& \DIFadd{280 }& \DIFadd{0.0301 }& {[}\DIFadd{0.0237,
0.0364}{]} & \DIFadd{0.5534 }& \DIFadd{5.62e-18 }& \DIFadd{144/46/90 }\\
\DIFadd{Untuned hex vs untuned RNG }& \DIFadd{QE\_H }& \DIFadd{lower }& \DIFadd{280 }& \DIFadd{0.1523 }& {[}\DIFadd{0.1094,
0.1953}{]} & \DIFadd{0.4174 }& \DIFadd{2.10e-11 }& \DIFadd{206/74/0 }\\
\DIFadd{Untuned hex vs untuned RNG }& \DIFadd{QE\_T }& \DIFadd{lower }& \DIFadd{280 }& \DIFadd{0.3464 }& {[}\DIFadd{0.2480,
0.4447}{]} & \DIFadd{0.4143 }& \DIFadd{2.88e-11 }& \DIFadd{246/34/0 }\\
\DIFadd{Untuned hex vs untuned RNG }& \DIFadd{QE\_B }& \DIFadd{lower }& \DIFadd{280 }& \DIFadd{0.2494 }& {[}\DIFadd{0.1808,
0.3179}{]} & \DIFadd{0.4281 }& \DIFadd{6.99e-12 }& \DIFadd{241/39/0 }\\
\DIFadd{Untuned hex vs untuned RNG }& \DIFadd{MTR\_H }& \DIFadd{lower }& \DIFadd{280 }& \DIFadd{2.3952 }& {[}\DIFadd{2.1221,
2.6683}{]} & \DIFadd{1.0316 }& \DIFadd{6.30e-46 }& \DIFadd{256/24/0 }\\
\DIFadd{Untuned hex vs untuned RNG }& \DIFadd{MTR\_T }& \DIFadd{lower }& \DIFadd{280 }& \DIFadd{2.4838 }& {[}\DIFadd{2.2596,
2.7080}{]} & \DIFadd{1.3032 }& \DIFadd{3.18e-62 }& \DIFadd{272/8/0 }\\
\DIFadd{Untuned hex vs untuned RNG }& \DIFadd{MTR\_B }& \DIFadd{lower }& \DIFadd{280 }& \DIFadd{2.4395 }& {[}\DIFadd{2.2033,
2.6757}{]} & \DIFadd{1.2152 }& \DIFadd{5.16e-57 }& \DIFadd{269/11/0 }\\
\DIFadd{Untuned hex vs untuned RNG }& \DIFadd{Util\_H }& \DIFadd{higher }& \DIFadd{280 }& \DIFadd{0.0225 }&
{[}\DIFadd{0.0157, 0.0294}{]} & \DIFadd{0.3886 }& \DIFadd{3.65e-10 }& \DIFadd{125/55/100 }\\
\DIFadd{Untuned hex vs untuned RNG }& \DIFadd{Util\_T }& \DIFadd{higher }& \DIFadd{280 }& \DIFadd{0.0432 }&
{[}\DIFadd{0.0361, 0.0504}{]} & \DIFadd{0.7162 }& \DIFadd{5.62e-27 }& \DIFadd{157/10/113 }\\
\DIFadd{Untuned hex vs untuned RNG }& \DIFadd{Util\_B }& \DIFadd{higher }& \DIFadd{280 }& \DIFadd{0.0329 }&
{[}\DIFadd{0.0266, 0.0392}{]} & \DIFadd{0.6130 }& \DIFadd{3.69e-21 }& \DIFadd{157/30/93 }\\
\DIFadd{Untuned hex vs untuned RNG }& \DIFadd{Dead\_H }& \DIFadd{lower }& \DIFadd{280 }& \DIFadd{0.0225 }& {[}\DIFadd{0.0157,
0.0294}{]} & \DIFadd{0.3886 }& \DIFadd{3.65e-10 }& \DIFadd{125/55/100 }\\
\DIFadd{Untuned hex vs untuned RNG }& \DIFadd{Dead\_T }& \DIFadd{lower }& \DIFadd{280 }& \DIFadd{0.0432 }& {[}\DIFadd{0.0361,
0.0504}{]} & \DIFadd{0.7162 }& \DIFadd{5.62e-27 }& \DIFadd{157/10/113 }\\
\DIFadd{Untuned hex vs untuned RNG }& \DIFadd{Dead\_B }& \DIFadd{lower }& \DIFadd{280 }& \DIFadd{0.0329 }& {[}\DIFadd{0.0266,
0.0392}{]} & \DIFadd{0.6130 }& \DIFadd{3.69e-21 }& \DIFadd{157/30/93 }\\
\DIFadd{Untuned MST vs untuned RNG }& \DIFadd{QE\_H }& \DIFadd{lower }& \DIFadd{280 }& \DIFadd{-5.54e-06 }&
{[}\DIFadd{-0.0198, 0.0198}{]} & \DIFadd{-3.29e-05 }& \DIFadd{1.000 }& \DIFadd{138/142/0 }\\
\DIFadd{Untuned MST vs untuned RNG }& \DIFadd{QE\_T }& \DIFadd{lower }& \DIFadd{280 }& \DIFadd{-0.0398 }& {[}\DIFadd{-0.0616,
-0.0181}{]} & \DIFadd{-0.2159 }& \DIFadd{3.59e-04 }& \DIFadd{107/173/0 }\\
\DIFadd{Untuned MST vs untuned RNG }& \DIFadd{QE\_B }& \DIFadd{lower }& \DIFadd{280 }& \DIFadd{-0.0199 }& {[}\DIFadd{-0.0385,
-0.0013}{]} & \DIFadd{-0.1261 }& \DIFadd{0.036 }& \DIFadd{120/160/0 }\\
\DIFadd{Untuned MST vs untuned RNG }& \DIFadd{MTR\_H }& \DIFadd{lower }& \DIFadd{280 }& \DIFadd{3.0116 }& {[}\DIFadd{2.6987,
3.3245}{]} & \DIFadd{1.1324 }& \DIFadd{4.89e-52 }& \DIFadd{238/42/0 }\\
\DIFadd{Untuned MST vs untuned RNG }& \DIFadd{MTR\_T }& \DIFadd{lower }& \DIFadd{280 }& \DIFadd{2.2763 }& {[}\DIFadd{2.0287,
2.5240}{]} & \DIFadd{1.0815 }& \DIFadd{5.90e-49 }& \DIFadd{233/47/0 }\\
\DIFadd{Untuned MST vs untuned RNG }& \DIFadd{MTR\_B }& \DIFadd{lower }& \DIFadd{280 }& \DIFadd{2.6440 }& {[}\DIFadd{2.3810,
2.9070}{]} & \DIFadd{1.1826 }& \DIFadd{4.64e-55 }& \DIFadd{240/40/0 }\\
\DIFadd{Untuned MST vs untuned RNG }& \DIFadd{Util\_H }& \DIFadd{higher }& \DIFadd{280 }& \DIFadd{0.0038 }&
{[}\DIFadd{7.90e-05, 0.0075}{]} & \DIFadd{0.1201 }& \DIFadd{0.045 }& \DIFadd{93/64/123 }\\
\DIFadd{Untuned MST vs untuned RNG }& \DIFadd{Util\_T }& \DIFadd{higher }& \DIFadd{280 }& \DIFadd{0.0019 }&
{[}\DIFadd{-0.0012, 0.0050}{]} & \DIFadd{0.0712 }& \DIFadd{0.235 }& \DIFadd{82/58/140 }\\
\DIFadd{Untuned MST vs untuned RNG }& \DIFadd{Util\_B }& \DIFadd{higher }& \DIFadd{280 }& \DIFadd{0.0028 }&
{[}\DIFadd{-9.28e-05, 0.0058}{]} & \DIFadd{0.1139 }& \DIFadd{0.058 }& \DIFadd{101/69/110 }\\
\DIFadd{Untuned MST vs untuned RNG }& \DIFadd{Dead\_H }& \DIFadd{lower }& \DIFadd{280 }& \DIFadd{0.0038 }&
{[}\DIFadd{7.90e-05, 0.0075}{]} & \DIFadd{0.1201 }& \DIFadd{0.045 }& \DIFadd{93/64/123 }\\
\DIFadd{Untuned MST vs untuned RNG }& \DIFadd{Dead\_T }& \DIFadd{lower }& \DIFadd{280 }& \DIFadd{0.0019 }&
{[}\DIFadd{-0.0012, 0.0050}{]} & \DIFadd{0.0712 }& \DIFadd{0.235 }& \DIFadd{82/58/140 }\\
\DIFadd{Untuned MST vs untuned RNG }& \DIFadd{Dead\_B }& \DIFadd{lower }& \DIFadd{280 }& \DIFadd{0.0028 }&
{[}\DIFadd{-9.28e-05, 0.0058}{]} & \DIFadd{0.1139 }& \DIFadd{0.058 }& \DIFadd{101/69/110 }\\
\DIFadd{Tuned hex vs tuned MST }& \DIFadd{QE\_H }& \DIFadd{lower }& \DIFadd{280 }& \DIFadd{0.0404 }& {[}\DIFadd{0.0199,
0.0608}{]} & \DIFadd{0.2323 }& \DIFadd{1.27e-04 }& \DIFadd{199/81/0 }\\
\DIFadd{Tuned hex vs tuned MST }& \DIFadd{QE\_T }& \DIFadd{lower }& \DIFadd{280 }& \DIFadd{0.0872 }& {[}\DIFadd{0.0537,
0.1207}{]} & \DIFadd{0.3063 }& \DIFadd{5.54e-07 }& \DIFadd{246/34/0 }\\
\DIFadd{Tuned hex vs tuned MST }& \DIFadd{QE\_B }& \DIFadd{lower }& \DIFadd{280 }& \DIFadd{0.0638 }& {[}\DIFadd{0.0379,
0.0897}{]} & \DIFadd{0.2897 }& \DIFadd{2.07e-06 }& \DIFadd{226/54/0 }\\
\DIFadd{Tuned hex vs tuned MST }& \DIFadd{MTR\_H }& \DIFadd{lower }& \DIFadd{280 }& \DIFadd{22.2866 }& {[}\DIFadd{21.4521,
23.1212}{]} & \DIFadd{3.1417 }& \DIFadd{8.86e-147 }& \DIFadd{280/0/0 }\\
\DIFadd{Tuned hex vs tuned MST }& \DIFadd{MTR\_T }& \DIFadd{lower }& \DIFadd{280 }& \DIFadd{23.1240 }& {[}\DIFadd{22.2558,
23.9922}{]} & \DIFadd{3.1331 }& \DIFadd{1.77e-146 }& \DIFadd{280/0/0 }\\
\DIFadd{Tuned hex vs tuned MST }& \DIFadd{MTR\_B }& \DIFadd{lower }& \DIFadd{280 }& \DIFadd{22.7053 }& {[}\DIFadd{21.8739,
23.5368}{]} & \DIFadd{3.2126 }& \DIFadd{3.06e-149 }& \DIFadd{280/0/0 }\\
\DIFadd{Tuned hex vs tuned MST }& \DIFadd{Util\_H }& \DIFadd{higher }& \DIFadd{280 }& \DIFadd{0.0092 }& {[}\DIFadd{0.0055,
0.0129}{]} & \DIFadd{0.2905 }& \DIFadd{1.95e-06 }& \DIFadd{116/63/101 }\\
\DIFadd{Tuned hex vs tuned MST }& \DIFadd{Util\_T }& \DIFadd{higher }& \DIFadd{280 }& \DIFadd{0.0222 }& {[}\DIFadd{0.0185,
0.0260}{]} & \DIFadd{0.6975 }& \DIFadd{6.79e-26 }& \DIFadd{150/13/117 }\\
\DIFadd{Tuned hex vs tuned MST }& \DIFadd{Util\_B }& \DIFadd{higher }& \DIFadd{280 }& \DIFadd{0.0157 }& {[}\DIFadd{0.0125,
0.0189}{]} & \DIFadd{0.5798 }& \DIFadd{2.30e-19 }& \DIFadd{139/47/94 }\\
\DIFadd{Tuned hex vs tuned MST }& \DIFadd{Dead\_H }& \DIFadd{lower }& \DIFadd{280 }& \DIFadd{0.0092 }& {[}\DIFadd{0.0055,
0.0129}{]} & \DIFadd{0.2905 }& \DIFadd{1.95e-06 }& \DIFadd{116/63/101 }\\
\DIFadd{Tuned hex vs tuned MST }& \DIFadd{Dead\_T }& \DIFadd{lower }& \DIFadd{280 }& \DIFadd{0.0222 }& {[}\DIFadd{0.0185,
0.0260}{]} & \DIFadd{0.6975 }& \DIFadd{6.79e-26 }& \DIFadd{150/13/117 }\\
\DIFadd{Tuned hex vs tuned MST }& \DIFadd{Dead\_B }& \DIFadd{lower }& \DIFadd{280 }& \DIFadd{0.0157 }& {[}\DIFadd{0.0125,
0.0189}{]} & \DIFadd{0.5798 }& \DIFadd{2.30e-19 }& \DIFadd{140/47/93 }\\
\DIFadd{Tuned hex vs tuned RNG }& \DIFadd{QE\_H }& \DIFadd{lower }& \DIFadd{280 }& \DIFadd{0.0360 }& {[}\DIFadd{0.0186,
0.0533}{]} & \DIFadd{0.2434 }& \DIFadd{6.04e-05 }& \DIFadd{210/70/0 }\\
\DIFadd{Tuned hex vs tuned RNG }& \DIFadd{QE\_T }& \DIFadd{lower }& \DIFadd{280 }& \DIFadd{0.0671 }& {[}\DIFadd{0.0409,
0.0932}{]} & \DIFadd{0.3014 }& \DIFadd{8.25e-07 }& \DIFadd{245/35/0 }\\
\DIFadd{Tuned hex vs tuned RNG }& \DIFadd{QE\_B }& \DIFadd{lower }& \DIFadd{280 }& \DIFadd{0.0515 }& {[}\DIFadd{0.0318,
0.0712}{]} & \DIFadd{0.3077 }& \DIFadd{4.94e-07 }& \DIFadd{240/40/0 }\\
\DIFadd{Tuned hex vs tuned RNG }& \DIFadd{MTR\_H }& \DIFadd{lower }& \DIFadd{280 }& \DIFadd{25.4035 }& {[}\DIFadd{24.5013,
26.3057}{]} & \DIFadd{3.3124 }& \DIFadd{1.24e-152 }& \DIFadd{280/0/0 }\\
\DIFadd{Tuned hex vs tuned RNG }& \DIFadd{MTR\_T }& \DIFadd{lower }& \DIFadd{280 }& \DIFadd{25.9296 }& {[}\DIFadd{25.0046,
26.8545}{]} & \DIFadd{3.2979 }& \DIFadd{3.81e-152 }& \DIFadd{280/0/0 }\\
\DIFadd{Tuned hex vs tuned RNG }& \DIFadd{MTR\_B }& \DIFadd{lower }& \DIFadd{280 }& \DIFadd{25.6665 }& {[}\DIFadd{24.7693,
26.5638}{]} & \DIFadd{3.3652 }& \DIFadd{2.17e-154 }& \DIFadd{280/0/0 }\\
\DIFadd{Tuned hex vs tuned RNG }& \DIFadd{Util\_H }& \DIFadd{higher }& \DIFadd{280 }& \DIFadd{0.0117 }& {[}\DIFadd{0.0076,
0.0158}{]} & \DIFadd{0.3348 }& \DIFadd{5.07e-08 }& \DIFadd{110/59/111 }\\
\DIFadd{Tuned hex vs tuned RNG }& \DIFadd{Util\_T }& \DIFadd{higher }& \DIFadd{280 }& \DIFadd{0.0276 }& {[}\DIFadd{0.0226,
0.0326}{]} & \DIFadd{0.6496 }& \DIFadd{3.52e-23 }& \DIFadd{143/12/125 }\\
\DIFadd{Tuned hex vs tuned RNG }& \DIFadd{Util\_B }& \DIFadd{higher }& \DIFadd{280 }& \DIFadd{0.0197 }& {[}\DIFadd{0.0158,
0.0235}{]} & \DIFadd{0.6025 }& \DIFadd{1.38e-20 }& \DIFadd{138/40/102 }\\
\DIFadd{Tuned hex vs tuned RNG }& \DIFadd{Dead\_H }& \DIFadd{lower }& \DIFadd{280 }& \DIFadd{0.0117 }& {[}\DIFadd{0.0076,
0.0158}{]} & \DIFadd{0.3348 }& \DIFadd{5.07e-08 }& \DIFadd{110/59/111 }\\
\DIFadd{Tuned hex vs tuned RNG }& \DIFadd{Dead\_T }& \DIFadd{lower }& \DIFadd{280 }& \DIFadd{0.0276 }& {[}\DIFadd{0.0226,
0.0326}{]} & \DIFadd{0.6496 }& \DIFadd{3.52e-23 }& \DIFadd{143/12/125 }\\
\DIFadd{Tuned hex vs tuned RNG }& \DIFadd{Dead\_B }& \DIFadd{lower }& \DIFadd{280 }& \DIFadd{0.0197 }& {[}\DIFadd{0.0158,
0.0235}{]} & \DIFadd{0.6025 }& \DIFadd{1.38e-20 }& \DIFadd{138/40/102 }\\
\DIFadd{Tuned MST vs tuned RNG }& \DIFadd{QE\_H }& \DIFadd{lower }& \DIFadd{280 }& \DIFadd{-0.0044 }& {[}\DIFadd{-0.0218,
0.0130}{]} & \DIFadd{-0.0298 }& \DIFadd{0.619 }& \DIFadd{161/119/0 }\\
\DIFadd{Tuned MST vs tuned RNG }& \DIFadd{QE\_T }& \DIFadd{lower }& \DIFadd{280 }& \DIFadd{-0.0201 }& {[}\DIFadd{-0.0395,
-8.07e-04}{]} & \DIFadd{-0.1225 }& \DIFadd{0.041 }& \DIFadd{161/119/0 }\\
\DIFadd{Tuned MST vs tuned RNG }& \DIFadd{QE\_B }& \DIFadd{lower }& \DIFadd{280 }& \DIFadd{-0.0123 }& {[}\DIFadd{-0.0280,
0.0035}{]} & \DIFadd{-0.0916 }& \DIFadd{0.126 }& \DIFadd{167/113/0 }\\
\DIFadd{Tuned MST vs tuned RNG }& \DIFadd{MTR\_H }& \DIFadd{lower }& \DIFadd{280 }& \DIFadd{3.1169 }& {[}\DIFadd{2.8270,
3.4067}{]} & \DIFadd{1.2650 }& \DIFadd{5.69e-60 }& \DIFadd{254/26/0 }\\
\DIFadd{Tuned MST vs tuned RNG }& \DIFadd{MTR\_T }& \DIFadd{lower }& \DIFadd{280 }& \DIFadd{2.8056 }& {[}\DIFadd{2.5596,
3.0516}{]} & \DIFadd{1.3416 }& \DIFadd{1.81e-64 }& \DIFadd{257/23/0 }\\
\DIFadd{Tuned MST vs tuned RNG }& \DIFadd{MTR\_B }& \DIFadd{lower }& \DIFadd{280 }& \DIFadd{2.9612 }& {[}\DIFadd{2.7049,
3.2175}{]} & \DIFadd{1.3591 }& \DIFadd{1.75e-65 }& \DIFadd{260/20/0 }\\
\DIFadd{Tuned MST vs tuned RNG }& \DIFadd{Util\_H }& \DIFadd{higher }& \DIFadd{280 }& \DIFadd{0.0025 }& {[}\DIFadd{-3.62e-04,
0.0054}{]} & \DIFadd{0.1029 }& \DIFadd{0.086 }& \DIFadd{75/79/126 }\\
\DIFadd{Tuned MST vs tuned RNG }& \DIFadd{Util\_T }& \DIFadd{higher }& \DIFadd{280 }& \DIFadd{0.0054 }& {[}\DIFadd{0.0027,
0.0081}{]} & \DIFadd{0.2361 }& \DIFadd{9.90e-05 }& \DIFadd{69/51/160 }\\
\DIFadd{Tuned MST vs tuned RNG }& \DIFadd{Util\_B }& \DIFadd{higher }& \DIFadd{280 }& \DIFadd{0.0040 }& {[}\DIFadd{0.0017,
0.0062}{]} & \DIFadd{0.2073 }& \DIFadd{6.07e-04 }& \DIFadd{87/80/113 }\\
\DIFadd{Tuned MST vs tuned RNG }& \DIFadd{Dead\_H }& \DIFadd{lower }& \DIFadd{280 }& \DIFadd{0.0025 }& {[}\DIFadd{-3.62e-04,
0.0054}{]} & \DIFadd{0.1029 }& \DIFadd{0.086 }& \DIFadd{75/79/126 }\\
\DIFadd{Tuned MST vs tuned RNG }& \DIFadd{Dead\_T }& \DIFadd{lower }& \DIFadd{280 }& \DIFadd{0.0054 }& {[}\DIFadd{0.0027,
0.0081}{]} & \DIFadd{0.2361 }& \DIFadd{9.90e-05 }& \DIFadd{69/51/160 }\\
\DIFadd{Tuned MST vs tuned RNG }& \DIFadd{Dead\_B }& \DIFadd{lower }& \DIFadd{280 }& \DIFadd{0.0040 }& {[}\DIFadd{0.0017,
0.0062}{]} & \DIFadd{0.2073 }& \DIFadd{6.07e-04 }& \DIFadd{87/80/113 }\\
\DIFadd{Tuned vs untuned hex }& \DIFadd{QE\_H }& \DIFadd{lower }& \DIFadd{280 }& \DIFadd{0.1651 }& {[}\DIFadd{0.1227,
0.2075}{]} & \DIFadd{0.4580 }& \DIFadd{2.99e-13 }& \DIFadd{230/50/0 }\\
\DIFadd{Tuned vs untuned hex }& \DIFadd{QE\_T }& \DIFadd{lower }& \DIFadd{280 }& \DIFadd{0.6969 }& {[}\DIFadd{0.4891,
0.9047}{]} & \DIFadd{0.3945 }& \DIFadd{2.04e-10 }& \DIFadd{280/0/0 }\\
\DIFadd{Tuned vs untuned hex }& \DIFadd{QE\_B }& \DIFadd{lower }& \DIFadd{280 }& \DIFadd{0.4310 }& {[}\DIFadd{0.3117,
0.5503}{]} & \DIFadd{0.4249 }& \DIFadd{9.78e-12 }& \DIFadd{280/0/0 }\\
\DIFadd{Tuned vs untuned hex }& \DIFadd{MTR\_H }& \DIFadd{lower }& \DIFadd{280 }& \DIFadd{-24.2851 }& {[}\DIFadd{-25.2036,
-23.3666}{]} & \DIFadd{-3.1103 }& \DIFadd{1.12e-145 }& \DIFadd{0/280/0 }\\
\DIFadd{Tuned vs untuned hex }& \DIFadd{MTR\_T }& \DIFadd{lower }& \DIFadd{280 }& \DIFadd{-24.3956 }& {[}\DIFadd{-25.3383,
-23.4529}{]} & \DIFadd{-3.0444 }& \DIFadd{2.51e-143 }& \DIFadd{0/280/0 }\\
\DIFadd{Tuned vs untuned hex }& \DIFadd{MTR\_B }& \DIFadd{lower }& \DIFadd{280 }& \DIFadd{-24.3403 }& {[}\DIFadd{-25.2543,
-23.4263}{]} & \DIFadd{-3.1328 }& \DIFadd{1.82e-146 }& \DIFadd{0/280/0 }\\
\DIFadd{Tuned vs untuned hex }& \DIFadd{Util\_H }& \DIFadd{higher }& \DIFadd{280 }& \DIFadd{0.0087 }& {[}\DIFadd{0.0023,
0.0150}{]} & \DIFadd{0.1610 }& \DIFadd{0.007 }& \DIFadd{98/93/89 }\\
\DIFadd{Tuned vs untuned hex }& \DIFadd{Util\_T }& \DIFadd{higher }& \DIFadd{280 }& \DIFadd{0.0380 }& {[}\DIFadd{0.0321,
0.0439}{]} & \DIFadd{0.7533 }& \DIFadd{3.85e-29 }& \DIFadd{150/17/113 }\\
\DIFadd{Tuned vs untuned hex }& \DIFadd{Util\_B }& \DIFadd{higher }& \DIFadd{280 }& \DIFadd{0.0233 }& {[}\DIFadd{0.0180,
0.0287}{]} & \DIFadd{0.5149 }& \DIFadd{5.24e-16 }& \DIFadd{139/53/88 }\\
\DIFadd{Tuned vs untuned hex }& \DIFadd{Dead\_H }& \DIFadd{lower }& \DIFadd{280 }& \DIFadd{0.0087 }& {[}\DIFadd{0.0023,
0.0150}{]} & \DIFadd{0.1610 }& \DIFadd{0.007 }& \DIFadd{98/93/89 }\\
\DIFadd{Tuned vs untuned hex }& \DIFadd{Dead\_T }& \DIFadd{lower }& \DIFadd{280 }& \DIFadd{0.0380 }& {[}\DIFadd{0.0321,
0.0439}{]} & \DIFadd{0.7533 }& \DIFadd{3.85e-29 }& \DIFadd{150/17/113 }\\
\DIFadd{Tuned vs untuned hex }& \DIFadd{Dead\_B }& \DIFadd{lower }& \DIFadd{280 }& \DIFadd{0.0233 }& {[}\DIFadd{0.0180,
0.0287}{]} & \DIFadd{0.5149 }& \DIFadd{5.24e-16 }& \DIFadd{139/53/88 }\\
\DIFadd{Tuned vs untuned MST }& \DIFadd{QE\_H }& \DIFadd{lower }& \DIFadd{280 }& \DIFadd{0.0531 }& {[}\DIFadd{0.0319,
0.0744}{]} & \DIFadd{0.2939 }& \DIFadd{1.49e-06 }& \DIFadd{224/56/0 }\\
\DIFadd{Tuned vs untuned MST }& \DIFadd{QE\_T }& \DIFadd{lower }& \DIFadd{280 }& \DIFadd{0.3979 }& {[}\DIFadd{0.2681,
0.5276}{]} & \DIFadd{0.3607 }& \DIFadd{5.03e-09 }& \DIFadd{280/0/0 }\\
\DIFadd{Tuned vs untuned MST }& \DIFadd{QE\_B }& \DIFadd{lower }& \DIFadd{280 }& \DIFadd{0.2255 }& {[}\DIFadd{0.1546,
0.2963}{]} & \DIFadd{0.3744 }& \DIFadd{1.41e-09 }& \DIFadd{276/4/0 }\\
\DIFadd{Tuned vs untuned MST }& \DIFadd{MTR\_H }& \DIFadd{lower }& \DIFadd{280 }& \DIFadd{-1.3820 }& {[}\DIFadd{-1.6350,
-1.1290}{]} & \DIFadd{-0.6427 }& \DIFadd{8.56e-23 }& \DIFadd{55/225/0 }\\
\DIFadd{Tuned vs untuned MST }& \DIFadd{MTR\_T }& \DIFadd{lower }& \DIFadd{280 }& \DIFadd{-1.4791 }& {[}\DIFadd{-1.6529,
-1.3052}{]} & \DIFadd{-1.0006 }& \DIFadd{4.89e-44 }& \DIFadd{27/253/0 }\\
\DIFadd{Tuned vs untuned MST }& \DIFadd{MTR\_B }& \DIFadd{lower }& \DIFadd{280 }& \DIFadd{-1.4305 }& {[}\DIFadd{-1.6184,
-1.2427}{]} & \DIFadd{-0.8959 }& \DIFadd{1.14e-37 }& \DIFadd{39/241/0 }\\
\DIFadd{Tuned vs untuned MST }& \DIFadd{Util\_H }& \DIFadd{higher }& \DIFadd{280 }& \DIFadd{-8.93e-04 }& {[}\DIFadd{-0.0052,
0.0034}{]} & \DIFadd{-0.0247 }& \DIFadd{0.680 }& \DIFadd{75/100/105 }\\
\DIFadd{Tuned vs untuned MST }& \DIFadd{Util\_T }& \DIFadd{higher }& \DIFadd{280 }& \DIFadd{0.0189 }& {[}\DIFadd{0.0145,
0.0232}{]} & \DIFadd{0.5097 }& \DIFadd{9.49e-16 }& \DIFadd{121/37/122 }\\
\DIFadd{Tuned vs untuned MST }& \DIFadd{Util\_B }& \DIFadd{higher }& \DIFadd{280 }& \DIFadd{0.0090 }& {[}\DIFadd{0.0054,
0.0126}{]} & \DIFadd{0.2935 }& \DIFadd{1.55e-06 }& \DIFadd{103/79/98 }\\
\DIFadd{Tuned vs untuned MST }& \DIFadd{Dead\_H }& \DIFadd{lower }& \DIFadd{280 }& \DIFadd{-8.93e-04 }& {[}\DIFadd{-0.0052,
0.0034}{]} & \DIFadd{-0.0247 }& \DIFadd{0.680 }& \DIFadd{75/100/105 }\\
\DIFadd{Tuned vs untuned MST }& \DIFadd{Dead\_T }& \DIFadd{lower }& \DIFadd{280 }& \DIFadd{0.0189 }& {[}\DIFadd{0.0145,
0.0232}{]} & \DIFadd{0.5097 }& \DIFadd{9.49e-16 }& \DIFadd{121/37/122 }\\
\DIFadd{Tuned vs untuned MST }& \DIFadd{Dead\_B }& \DIFadd{lower }& \DIFadd{280 }& \DIFadd{0.0090 }& {[}\DIFadd{0.0054,
0.0126}{]} & \DIFadd{0.2935 }& \DIFadd{1.55e-06 }& \DIFadd{103/79/98 }\\
\DIFadd{Tuned vs untuned RNG }& \DIFadd{QE\_H }& \DIFadd{lower }& \DIFadd{280 }& \DIFadd{0.0487 }& {[}\DIFadd{0.0261,
0.0713}{]} & \DIFadd{0.2532 }& \DIFadd{3.07e-05 }& \DIFadd{223/57/0 }\\
\DIFadd{Tuned vs untuned RNG }& \DIFadd{QE\_T }& \DIFadd{lower }& \DIFadd{280 }& \DIFadd{0.4176 }& {[}\DIFadd{0.2918,
0.5434}{]} & \DIFadd{0.3905 }& \DIFadd{3.02e-10 }& \DIFadd{280/0/0 }\\
\DIFadd{Tuned vs untuned RNG }& \DIFadd{QE\_B }& \DIFadd{lower }& \DIFadd{280 }& \DIFadd{0.2331 }& {[}\DIFadd{0.1681,
0.2982}{]} & \DIFadd{0.4215 }& \DIFadd{1.39e-11 }& \DIFadd{280/0/0 }\\
\DIFadd{Tuned vs untuned RNG }& \DIFadd{MTR\_H }& \DIFadd{lower }& \DIFadd{280 }& \DIFadd{-1.2768 }& {[}\DIFadd{-1.4669,
-1.0867}{]} & \DIFadd{-0.7902 }& \DIFadd{2.56e-31 }& \DIFadd{35/245/0 }\\
\DIFadd{Tuned vs untuned RNG }& \DIFadd{MTR\_T }& \DIFadd{lower }& \DIFadd{280 }& \DIFadd{-0.9498 }& {[}\DIFadd{-1.0619,
-0.8378}{]} & \DIFadd{-0.9973 }& \DIFadd{7.77e-44 }& \DIFadd{21/259/0 }\\
\DIFadd{Tuned vs untuned RNG }& \DIFadd{MTR\_B }& \DIFadd{lower }& \DIFadd{280 }& \DIFadd{-1.1133 }& {[}\DIFadd{-1.2485,
-0.9781}{]} & \DIFadd{-0.9686 }& \DIFadd{4.33e-42 }& \DIFadd{28/252/0 }\\
\DIFadd{Tuned vs untuned RNG }& \DIFadd{Util\_H }& \DIFadd{higher }& \DIFadd{280 }& \DIFadd{-0.0021 }& {[}\DIFadd{-0.0064,
0.0021}{]} & \DIFadd{-0.0599 }& \DIFadd{0.317 }& \DIFadd{83/92/105 }\\
\DIFadd{Tuned vs untuned RNG }& \DIFadd{Util\_T }& \DIFadd{higher }& \DIFadd{280 }& \DIFadd{0.0224 }& {[}\DIFadd{0.0180,
0.0268}{]} & \DIFadd{0.6043 }& \DIFadd{1.10e-20 }& \DIFadd{134/19/127 }\\
\DIFadd{Tuned vs untuned RNG }& \DIFadd{Util\_B }& \DIFadd{higher }& \DIFadd{280 }& \DIFadd{0.0101 }& {[}\DIFadd{0.0067,
0.0136}{]} & \DIFadd{0.3465 }& \DIFadd{1.82e-08 }& \DIFadd{116/68/96 }\\
\DIFadd{Tuned vs untuned RNG }& \DIFadd{Dead\_H }& \DIFadd{lower }& \DIFadd{280 }& \DIFadd{-0.0021 }& {[}\DIFadd{-0.0064,
0.0021}{]} & \DIFadd{-0.0599 }& \DIFadd{0.317 }& \DIFadd{83/92/105 }\\
\DIFadd{Tuned vs untuned RNG }& \DIFadd{Dead\_T }& \DIFadd{lower }& \DIFadd{280 }& \DIFadd{0.0224 }& {[}\DIFadd{0.0180,
0.0268}{]} & \DIFadd{0.6043 }& \DIFadd{1.10e-20 }& \DIFadd{134/19/127 }\\
\DIFadd{Tuned vs untuned RNG }& \DIFadd{Dead\_B }& \DIFadd{lower }& \DIFadd{280 }& \DIFadd{0.0101 }& {[}\DIFadd{0.0067,
0.0136}{]} & \DIFadd{0.3465 }& \DIFadd{1.82e-08 }& \DIFadd{117/68/95 }\\
\end{longtable}
}

\DIFadd{\textbf{Supplementary Table S13. Full matched topology diagnostic means
by profile, dataset, and topology.} Rows report the mean value across
the final matched random seeds for each profile, dataset, topology, and
full-sampling setting. \texttt{untuned} denotes the untuned reference
profile using XPySOM-like default hyperparameters and \texttt{tuned}
denotes the fixed QE-tuned profile. Lower is better for QE, MTR, and
dead-node fraction; higher is better for node utilization. Metric
suffixes denote holdout (\texttt{\_H}), train (\texttt{\_T}), and
balanced train-holdout (\texttt{\_B}) summaries.
}

{\def\LTcaptype{none} %DIF >  do not increment counter
\begin{longtable}[]{@{}
  >{\raggedright\arraybackslash}p{(\linewidth - 30\tabcolsep) * \real{0.0492}}
  >{\raggedright\arraybackslash}p{(\linewidth - 30\tabcolsep) * \real{0.0492}}
  >{\raggedright\arraybackslash}p{(\linewidth - 30\tabcolsep) * \real{0.0492}}
  >{\raggedleft\arraybackslash}p{(\linewidth - 30\tabcolsep) * \real{0.0656}}
  >{\raggedleft\arraybackslash}p{(\linewidth - 30\tabcolsep) * \real{0.0656}}
  >{\raggedleft\arraybackslash}p{(\linewidth - 30\tabcolsep) * \real{0.0656}}
  >{\raggedleft\arraybackslash}p{(\linewidth - 30\tabcolsep) * \real{0.0656}}
  >{\raggedleft\arraybackslash}p{(\linewidth - 30\tabcolsep) * \real{0.0656}}
  >{\raggedleft\arraybackslash}p{(\linewidth - 30\tabcolsep) * \real{0.0656}}
  >{\raggedleft\arraybackslash}p{(\linewidth - 30\tabcolsep) * \real{0.0656}}
  >{\raggedleft\arraybackslash}p{(\linewidth - 30\tabcolsep) * \real{0.0656}}
  >{\raggedleft\arraybackslash}p{(\linewidth - 30\tabcolsep) * \real{0.0656}}
  >{\raggedleft\arraybackslash}p{(\linewidth - 30\tabcolsep) * \real{0.0656}}
  >{\raggedleft\arraybackslash}p{(\linewidth - 30\tabcolsep) * \real{0.0656}}
  >{\raggedleft\arraybackslash}p{(\linewidth - 30\tabcolsep) * \real{0.0656}}
  >{\raggedleft\arraybackslash}p{(\linewidth - 30\tabcolsep) * \real{0.0656}}@{}}
\toprule\noalign{}
\begin{minipage}[b]{\linewidth}\raggedright
\DIFadd{Profile
}\end{minipage} & \begin{minipage}[b]{\linewidth}\raggedright
\DIFadd{Dataset
}\end{minipage} & \begin{minipage}[b]{\linewidth}\raggedright
\DIFadd{Topology
}\end{minipage} & \begin{minipage}[b]{\linewidth}\raggedleft
\DIFadd{n
}\end{minipage} & \begin{minipage}[b]{\linewidth}\raggedleft
\DIFadd{QE\_H
}\end{minipage} & \begin{minipage}[b]{\linewidth}\raggedleft
\DIFadd{QE\_T
}\end{minipage} & \begin{minipage}[b]{\linewidth}\raggedleft
\DIFadd{QE\_B
}\end{minipage} & \begin{minipage}[b]{\linewidth}\raggedleft
\DIFadd{MTR\_H
}\end{minipage} & \begin{minipage}[b]{\linewidth}\raggedleft
\DIFadd{MTR\_T
}\end{minipage} & \begin{minipage}[b]{\linewidth}\raggedleft
\DIFadd{MTR\_B
}\end{minipage} & \begin{minipage}[b]{\linewidth}\raggedleft
\DIFadd{Util\_H
}\end{minipage} & \begin{minipage}[b]{\linewidth}\raggedleft
\DIFadd{Util\_T
}\end{minipage} & \begin{minipage}[b]{\linewidth}\raggedleft
\DIFadd{Util\_B
}\end{minipage} & \begin{minipage}[b]{\linewidth}\raggedleft
\DIFadd{Dead\_H
}\end{minipage} & \begin{minipage}[b]{\linewidth}\raggedleft
\DIFadd{Dead\_T
}\end{minipage} & \begin{minipage}[b]{\linewidth}\raggedleft
\DIFadd{Dead\_B
}\end{minipage} \\
\midrule\noalign{}
\endhead
\bottomrule\noalign{}
\endlastfoot
\DIFadd{untuned }& \DIFadd{blobs }& \DIFadd{hexagonal }& \DIFadd{20 }& \DIFadd{0.1321 }& \DIFadd{0.1313 }& \DIFadd{0.1317 }& \DIFadd{3.5132 }&
\DIFadd{3.5212 }& \DIFadd{3.5172 }& \DIFadd{1.0000 }& \DIFadd{1.0000 }& \DIFadd{1.0000 }& \DIFadd{0.0000 }& \DIFadd{0.0000 }& \DIFadd{0.0000 }\\
\DIFadd{untuned }& \DIFadd{blobs }& \DIFadd{mst }& \DIFadd{20 }& \DIFadd{0.1341 }& \DIFadd{0.1324 }& \DIFadd{0.1332 }& \DIFadd{5.2326 }& \DIFadd{5.1119
}& \DIFadd{5.1723 }& \DIFadd{1.0000 }& \DIFadd{1.0000 }& \DIFadd{1.0000 }& \DIFadd{0.0000 }& \DIFadd{0.0000 }& \DIFadd{0.0000 }\\
\DIFadd{untuned }& \DIFadd{blobs }& \DIFadd{rng }& \DIFadd{20 }& \DIFadd{0.1325 }& \DIFadd{0.1307 }& \DIFadd{0.1316 }& \DIFadd{2.7535 }& \DIFadd{2.7181
}& \DIFadd{2.7358 }& \DIFadd{1.0000 }& \DIFadd{1.0000 }& \DIFadd{1.0000 }& \DIFadd{0.0000 }& \DIFadd{0.0000 }& \DIFadd{0.0000 }\\
\DIFadd{untuned }& \DIFadd{breast\_cancer }& \DIFadd{hexagonal }& \DIFadd{20 }& \DIFadd{2.6474 }& \DIFadd{2.2132 }& \DIFadd{2.4303 }&
\DIFadd{9.9298 }& \DIFadd{6.7921 }& \DIFadd{8.3610 }& \DIFadd{0.7930 }& \DIFadd{0.9930 }& \DIFadd{0.8930 }& \DIFadd{0.2070 }& \DIFadd{0.0070 }&
\DIFadd{0.1070 }\\
\DIFadd{untuned }& \DIFadd{breast\_cancer }& \DIFadd{mst }& \DIFadd{20 }& \DIFadd{2.6624 }& \DIFadd{2.0732 }& \DIFadd{2.3678 }& \DIFadd{13.4902
}& \DIFadd{7.3257 }& \DIFadd{10.4079 }& \DIFadd{0.7445 }& \DIFadd{0.9935 }& \DIFadd{0.8690 }& \DIFadd{0.2555 }& \DIFadd{0.0065 }&
\DIFadd{0.1310 }\\
\DIFadd{untuned }& \DIFadd{breast\_cancer }& \DIFadd{rng }& \DIFadd{20 }& \DIFadd{2.6515 }& \DIFadd{2.0799 }& \DIFadd{2.3657 }& \DIFadd{7.7404
}& \DIFadd{4.1496 }& \DIFadd{5.9450 }& \DIFadd{0.7660 }& \DIFadd{0.9940 }& \DIFadd{0.8800 }& \DIFadd{0.2340 }& \DIFadd{0.0060 }&
\DIFadd{0.1200 }\\
\DIFadd{untuned }& \DIFadd{california\_housing }& \DIFadd{hexagonal }& \DIFadd{20 }& \DIFadd{0.7910 }& \DIFadd{0.7856 }&
\DIFadd{0.7883 }& \DIFadd{8.4626 }& \DIFadd{8.3360 }& \DIFadd{8.3993 }& \DIFadd{0.9995 }& \DIFadd{1.0000 }& \DIFadd{0.9998 }& \DIFadd{5.00e-04
}& \DIFadd{0.0000 }& \DIFadd{2.50e-04 }\\
\DIFadd{untuned }& \DIFadd{california\_housing }& \DIFadd{mst }& \DIFadd{20 }& \DIFadd{0.7600 }& \DIFadd{0.7451 }& \DIFadd{0.7526 }&
\DIFadd{7.8776 }& \DIFadd{7.6936 }& \DIFadd{7.7856 }& \DIFadd{0.9930 }& \DIFadd{1.0000 }& \DIFadd{0.9965 }& \DIFadd{0.0070 }& \DIFadd{0.0000 }&
\DIFadd{0.0035 }\\
\DIFadd{untuned }& \DIFadd{california\_housing }& \DIFadd{rng }& \DIFadd{20 }& \DIFadd{0.7646 }& \DIFadd{0.7522 }& \DIFadd{0.7584 }&
\DIFadd{4.5736 }& \DIFadd{4.4653 }& \DIFadd{4.5195 }& \DIFadd{0.9935 }& \DIFadd{1.0000 }& \DIFadd{0.9967 }& \DIFadd{0.0065 }& \DIFadd{0.0000 }&
\DIFadd{0.0033 }\\
\DIFadd{untuned }& \DIFadd{circles }& \DIFadd{hexagonal }& \DIFadd{20 }& \DIFadd{0.1611 }& \DIFadd{0.1590 }& \DIFadd{0.1601 }& \DIFadd{3.5646 }&
\DIFadd{3.5441 }& \DIFadd{3.5543 }& \DIFadd{1.0000 }& \DIFadd{1.0000 }& \DIFadd{1.0000 }& \DIFadd{0.0000 }& \DIFadd{0.0000 }& \DIFadd{0.0000 }\\
\DIFadd{untuned }& \DIFadd{circles }& \DIFadd{mst }& \DIFadd{20 }& \DIFadd{0.1653 }& \DIFadd{0.1620 }& \DIFadd{0.1636 }& \DIFadd{7.2056 }&
\DIFadd{6.9061 }& \DIFadd{7.0558 }& \DIFadd{1.0000 }& \DIFadd{1.0000 }& \DIFadd{1.0000 }& \DIFadd{0.0000 }& \DIFadd{0.0000 }& \DIFadd{0.0000 }\\
\DIFadd{untuned }& \DIFadd{circles }& \DIFadd{rng }& \DIFadd{20 }& \DIFadd{0.1627 }& \DIFadd{0.1599 }& \DIFadd{0.1613 }& \DIFadd{2.8684 }&
\DIFadd{2.8279 }& \DIFadd{2.8481 }& \DIFadd{1.0000 }& \DIFadd{1.0000 }& \DIFadd{1.0000 }& \DIFadd{0.0000 }& \DIFadd{0.0000 }& \DIFadd{0.0000 }\\
\DIFadd{untuned }& \DIFadd{covertype }& \DIFadd{hexagonal }& \DIFadd{20 }& \DIFadd{3.2499 }& \DIFadd{3.2515 }& \DIFadd{3.2507 }& \DIFadd{3.6170
}& \DIFadd{3.6102 }& \DIFadd{3.6136 }& \DIFadd{0.6585 }& \DIFadd{0.6635 }& \DIFadd{0.6610 }& \DIFadd{0.3415 }& \DIFadd{0.3365 }&
\DIFadd{0.3390 }\\
\DIFadd{untuned }& \DIFadd{covertype }& \DIFadd{mst }& \DIFadd{20 }& \DIFadd{2.4536 }& \DIFadd{2.4535 }& \DIFadd{2.4536 }& \DIFadd{2.3593 }&
\DIFadd{2.3616 }& \DIFadd{2.3604 }& \DIFadd{0.8460 }& \DIFadd{0.8470 }& \DIFadd{0.8465 }& \DIFadd{0.1540 }& \DIFadd{0.1530 }& \DIFadd{0.1535 }\\
\DIFadd{untuned }& \DIFadd{covertype }& \DIFadd{rng }& \DIFadd{20 }& \DIFadd{2.6068 }& \DIFadd{2.6067 }& \DIFadd{2.6068 }& \DIFadd{2.4747 }&
\DIFadd{2.4737 }& \DIFadd{2.4742 }& \DIFadd{0.8545 }& \DIFadd{0.8550 }& \DIFadd{0.8547 }& \DIFadd{0.1455 }& \DIFadd{0.1450 }& \DIFadd{0.1453 }\\
\DIFadd{untuned }& \DIFadd{diabetes }& \DIFadd{hexagonal }& \DIFadd{20 }& \DIFadd{1.6952 }& \DIFadd{1.3289 }& \DIFadd{1.5121 }& \DIFadd{13.1182
}& \DIFadd{7.7649 }& \DIFadd{10.4416 }& \DIFadd{0.6970 }& \DIFadd{0.9295 }& \DIFadd{0.8133 }& \DIFadd{0.3030 }& \DIFadd{0.0705 }&
\DIFadd{0.1867 }\\
\DIFadd{untuned }& \DIFadd{diabetes }& \DIFadd{mst }& \DIFadd{20 }& \DIFadd{1.6718 }& \DIFadd{1.2181 }& \DIFadd{1.4450 }& \DIFadd{14.5970 }&
\DIFadd{7.9549 }& \DIFadd{11.2760 }& \DIFadd{0.7015 }& \DIFadd{0.9725 }& \DIFadd{0.8370 }& \DIFadd{0.2985 }& \DIFadd{0.0275 }&
\DIFadd{0.1630 }\\
\DIFadd{untuned }& \DIFadd{diabetes }& \DIFadd{rng }& \DIFadd{20 }& \DIFadd{1.6803 }& \DIFadd{1.2362 }& \DIFadd{1.4583 }& \DIFadd{9.7368 }&
\DIFadd{4.1324 }& \DIFadd{6.9346 }& \DIFadd{0.7065 }& \DIFadd{0.9815 }& \DIFadd{0.8440 }& \DIFadd{0.2935 }& \DIFadd{0.0185 }& \DIFadd{0.1560 }\\
\DIFadd{untuned }& \DIFadd{digits }& \DIFadd{hexagonal }& \DIFadd{20 }& \DIFadd{4.5856 }& \DIFadd{4.2735 }& \DIFadd{4.4296 }& \DIFadd{7.6155 }&
\DIFadd{6.3931 }& \DIFadd{7.0043 }& \DIFadd{0.9205 }& \DIFadd{0.9665 }& \DIFadd{0.9435 }& \DIFadd{0.0795 }& \DIFadd{0.0335 }& \DIFadd{0.0565 }\\
\DIFadd{untuned }& \DIFadd{digits }& \DIFadd{mst }& \DIFadd{20 }& \DIFadd{4.4317 }& \DIFadd{3.9977 }& \DIFadd{4.2147 }& \DIFadd{7.2430 }& \DIFadd{5.5079
}& \DIFadd{6.3754 }& \DIFadd{0.9515 }& \DIFadd{0.9950 }& \DIFadd{0.9732 }& \DIFadd{0.0485 }& \DIFadd{0.0050 }& \DIFadd{0.0268 }\\
\DIFadd{untuned }& \DIFadd{digits }& \DIFadd{rng }& \DIFadd{20 }& \DIFadd{4.4481 }& \DIFadd{4.0470 }& \DIFadd{4.2475 }& \DIFadd{5.0653 }& \DIFadd{3.8881
}& \DIFadd{4.4767 }& \DIFadd{0.9625 }& \DIFadd{0.9980 }& \DIFadd{0.9803 }& \DIFadd{0.0375 }& \DIFadd{0.0020 }& \DIFadd{0.0198 }\\
\DIFadd{untuned }& \DIFadd{iris }& \DIFadd{hexagonal }& \DIFadd{20 }& \DIFadd{0.3657 }& \DIFadd{0.1823 }& \DIFadd{0.2740 }& \DIFadd{7.8578 }&
\DIFadd{5.3890 }& \DIFadd{6.6234 }& \DIFadd{0.3315 }& \DIFadd{0.6690 }& \DIFadd{0.5003 }& \DIFadd{0.6685 }& \DIFadd{0.3310 }& \DIFadd{0.4997 }\\
\DIFadd{untuned }& \DIFadd{iris }& \DIFadd{mst }& \DIFadd{20 }& \DIFadd{0.3609 }& \DIFadd{0.1138 }& \DIFadd{0.2374 }& \DIFadd{5.7900 }& \DIFadd{1.8655 }&
\DIFadd{3.8277 }& \DIFadd{0.3330 }& \DIFadd{0.7785 }& \DIFadd{0.5557 }& \DIFadd{0.6670 }& \DIFadd{0.2215 }& \DIFadd{0.4442 }\\
\DIFadd{untuned }& \DIFadd{iris }& \DIFadd{rng }& \DIFadd{20 }& \DIFadd{0.3563 }& \DIFadd{0.1149 }& \DIFadd{0.2356 }& \DIFadd{3.4344 }& \DIFadd{1.7062 }&
\DIFadd{2.5703 }& \DIFadd{0.3385 }& \DIFadd{0.7910 }& \DIFadd{0.5647 }& \DIFadd{0.6615 }& \DIFadd{0.2090 }& \DIFadd{0.4353 }\\
\DIFadd{untuned }& \DIFadd{kddcup99 }& \DIFadd{hexagonal }& \DIFadd{20 }& \DIFadd{0.4678 }& \DIFadd{0.4662 }& \DIFadd{0.4670 }& \DIFadd{3.1121
}& \DIFadd{3.1058 }& \DIFadd{3.1089 }& \DIFadd{0.8850 }& \DIFadd{0.9140 }& \DIFadd{0.8995 }& \DIFadd{0.1150 }& \DIFadd{0.0860 }&
\DIFadd{0.1005 }\\
\DIFadd{untuned }& \DIFadd{kddcup99 }& \DIFadd{mst }& \DIFadd{20 }& \DIFadd{0.3386 }& \DIFadd{0.3365 }& \DIFadd{0.3376 }& \DIFadd{1.7243 }&
\DIFadd{1.7225 }& \DIFadd{1.7234 }& \DIFadd{0.9760 }& \DIFadd{0.9770 }& \DIFadd{0.9765 }& \DIFadd{0.0240 }& \DIFadd{0.0230 }& \DIFadd{0.0235 }\\
\DIFadd{untuned }& \DIFadd{kddcup99 }& \DIFadd{rng }& \DIFadd{20 }& \DIFadd{0.3730 }& \DIFadd{0.3712 }& \DIFadd{0.3721 }& \DIFadd{2.9179 }&
\DIFadd{2.9130 }& \DIFadd{2.9155 }& \DIFadd{0.9270 }& \DIFadd{0.9310 }& \DIFadd{0.9290 }& \DIFadd{0.0730 }& \DIFadd{0.0690 }& \DIFadd{0.0710 }\\
\DIFadd{untuned }& \DIFadd{moons }& \DIFadd{hexagonal }& \DIFadd{20 }& \DIFadd{0.1355 }& \DIFadd{0.1346 }& \DIFadd{0.1351 }& \DIFadd{3.1842 }&
\DIFadd{3.1787 }& \DIFadd{3.1815 }& \DIFadd{1.0000 }& \DIFadd{1.0000 }& \DIFadd{1.0000 }& \DIFadd{0.0000 }& \DIFadd{0.0000 }& \DIFadd{0.0000 }\\
\DIFadd{untuned }& \DIFadd{moons }& \DIFadd{mst }& \DIFadd{20 }& \DIFadd{0.1376 }& \DIFadd{0.1356 }& \DIFadd{0.1366 }& \DIFadd{4.8419 }& \DIFadd{4.7384
}& \DIFadd{4.7902 }& \DIFadd{1.0000 }& \DIFadd{1.0000 }& \DIFadd{1.0000 }& \DIFadd{0.0000 }& \DIFadd{0.0000 }& \DIFadd{0.0000 }\\
\DIFadd{untuned }& \DIFadd{moons }& \DIFadd{rng }& \DIFadd{20 }& \DIFadd{0.1357 }& \DIFadd{0.1342 }& \DIFadd{0.1349 }& \DIFadd{2.6343 }& \DIFadd{2.6125
}& \DIFadd{2.6234 }& \DIFadd{1.0000 }& \DIFadd{1.0000 }& \DIFadd{1.0000 }& \DIFadd{0.0000 }& \DIFadd{0.0000 }& \DIFadd{0.0000 }\\
\DIFadd{untuned }& \DIFadd{olivetti\_faces }& \DIFadd{hexagonal }& \DIFadd{20 }& \DIFadd{42.6773 }& \DIFadd{34.2471 }& \DIFadd{38.4622
}& \DIFadd{8.9338 }& \DIFadd{5.0171 }& \DIFadd{6.9754 }& \DIFadd{0.6115 }& \DIFadd{0.8175 }& \DIFadd{0.7145 }& \DIFadd{0.3885 }& \DIFadd{0.1825
}& \DIFadd{0.2855 }\\
\DIFadd{untuned }& \DIFadd{olivetti\_faces }& \DIFadd{mst }& \DIFadd{20 }& \DIFadd{41.6810 }& \DIFadd{30.7001 }& \DIFadd{36.1905 }&
\DIFadd{10.5162 }& \DIFadd{4.4246 }& \DIFadd{7.4704 }& \DIFadd{0.6285 }& \DIFadd{0.8775 }& \DIFadd{0.7530 }& \DIFadd{0.3715 }& \DIFadd{0.1225 }&
\DIFadd{0.2470 }\\
\DIFadd{untuned }& \DIFadd{olivetti\_faces }& \DIFadd{rng }& \DIFadd{20 }& \DIFadd{41.5030 }& \DIFadd{30.9796 }& \DIFadd{36.2413 }&
\DIFadd{6.9887 }& \DIFadd{2.7613 }& \DIFadd{4.8750 }& \DIFadd{0.6325 }& \DIFadd{0.8845 }& \DIFadd{0.7585 }& \DIFadd{0.3675 }& \DIFadd{0.1155 }&
\DIFadd{0.2415 }\\
\DIFadd{untuned }& \DIFadd{s\_curve }& \DIFadd{hexagonal }& \DIFadd{20 }& \DIFadd{0.3334 }& \DIFadd{0.3319 }& \DIFadd{0.3327 }& \DIFadd{7.8608
}& \DIFadd{7.8535 }& \DIFadd{7.8571 }& \DIFadd{0.9990 }& \DIFadd{1.0000 }& \DIFadd{0.9995 }& \DIFadd{0.0010 }& \DIFadd{0.0000 }&
\DIFadd{5.00e-04 }\\
\DIFadd{untuned }& \DIFadd{s\_curve }& \DIFadd{mst }& \DIFadd{20 }& \DIFadd{0.3193 }& \DIFadd{0.3155 }& \DIFadd{0.3174 }& \DIFadd{7.0170 }&
\DIFadd{6.7520 }& \DIFadd{6.8845 }& \DIFadd{1.0000 }& \DIFadd{1.0000 }& \DIFadd{1.0000 }& \DIFadd{0.0000 }& \DIFadd{0.0000 }& \DIFadd{0.0000 }\\
\DIFadd{untuned }& \DIFadd{s\_curve }& \DIFadd{rng }& \DIFadd{20 }& \DIFadd{0.3211 }& \DIFadd{0.3180 }& \DIFadd{0.3195 }& \DIFadd{2.9213 }&
\DIFadd{2.8719 }& \DIFadd{2.8966 }& \DIFadd{1.0000 }& \DIFadd{1.0000 }& \DIFadd{1.0000 }& \DIFadd{0.0000 }& \DIFadd{0.0000 }& \DIFadd{0.0000 }\\
\DIFadd{untuned }& \DIFadd{swiss\_roll }& \DIFadd{hexagonal }& \DIFadd{20 }& \DIFadd{0.2902 }& \DIFadd{0.2889 }& \DIFadd{0.2896 }&
\DIFadd{8.7591 }& \DIFadd{8.7911 }& \DIFadd{8.7751 }& \DIFadd{0.9325 }& \DIFadd{0.9385 }& \DIFadd{0.9355 }& \DIFadd{0.0675 }& \DIFadd{0.0615 }&
\DIFadd{0.0645 }\\
\DIFadd{untuned }& \DIFadd{swiss\_roll }& \DIFadd{mst }& \DIFadd{20 }& \DIFadd{0.2725 }& \DIFadd{0.2695 }& \DIFadd{0.2710 }& \DIFadd{9.4331 }&
\DIFadd{9.1974 }& \DIFadd{9.3152 }& \DIFadd{0.9240 }& \DIFadd{0.9265 }& \DIFadd{0.9252 }& \DIFadd{0.0760 }& \DIFadd{0.0735 }& \DIFadd{0.0747 }\\
\DIFadd{untuned }& \DIFadd{swiss\_roll }& \DIFadd{rng }& \DIFadd{20 }& \DIFadd{0.2721 }& \DIFadd{0.2695 }& \DIFadd{0.2708 }& \DIFadd{2.8013 }&
\DIFadd{2.7562 }& \DIFadd{2.7788 }& \DIFadd{0.9645 }& \DIFadd{0.9655 }& \DIFadd{0.9650 }& \DIFadd{0.0355 }& \DIFadd{0.0345 }& \DIFadd{0.0350 }\\
\DIFadd{untuned }& \DIFadd{wine }& \DIFadd{hexagonal }& \DIFadd{20 }& \DIFadd{1.9354 }& \DIFadd{1.1334 }& \DIFadd{1.5344 }& \DIFadd{7.9569 }&
\DIFadd{3.7968 }& \DIFadd{5.8769 }& \DIFadd{0.3885 }& \DIFadd{0.7155 }& \DIFadd{0.5520 }& \DIFadd{0.6115 }& \DIFadd{0.2845 }& \DIFadd{0.4480 }\\
\DIFadd{untuned }& \DIFadd{wine }& \DIFadd{mst }& \DIFadd{20 }& \DIFadd{1.9459 }& \DIFadd{0.8674 }& \DIFadd{1.4066 }& \DIFadd{8.7875 }& \DIFadd{2.6270 }&
\DIFadd{5.7073 }& \DIFadd{0.3810 }& \DIFadd{0.8185 }& \DIFadd{0.5998 }& \DIFadd{0.6190 }& \DIFadd{0.1815 }& \DIFadd{0.4003 }\\
\DIFadd{untuned }& \DIFadd{wine }& \DIFadd{rng }& \DIFadd{20 }& \DIFadd{1.9270 }& \DIFadd{0.8784 }& \DIFadd{1.4027 }& \DIFadd{7.0421 }& \DIFadd{2.0440 }&
\DIFadd{4.5430 }& \DIFadd{0.3865 }& \DIFadd{0.8120 }& \DIFadd{0.5992 }& \DIFadd{0.6135 }& \DIFadd{0.1880 }& \DIFadd{0.4008 }\\
\DIFadd{tuned }& \DIFadd{blobs }& \DIFadd{hexagonal }& \DIFadd{20 }& \DIFadd{0.1284 }& \DIFadd{0.1258 }& \DIFadd{0.1271 }& \DIFadd{27.4818 }&
\DIFadd{27.5131 }& \DIFadd{27.4974 }& \DIFadd{1.0000 }& \DIFadd{1.0000 }& \DIFadd{1.0000 }& \DIFadd{0.0000 }& \DIFadd{0.0000 }&
\DIFadd{0.0000 }\\
\DIFadd{tuned }& \DIFadd{blobs }& \DIFadd{mst }& \DIFadd{20 }& \DIFadd{0.1285 }& \DIFadd{0.1255 }& \DIFadd{0.1270 }& \DIFadd{6.7794 }& \DIFadd{6.6317 }&
\DIFadd{6.7055 }& \DIFadd{1.0000 }& \DIFadd{1.0000 }& \DIFadd{1.0000 }& \DIFadd{0.0000 }& \DIFadd{0.0000 }& \DIFadd{0.0000 }\\
\DIFadd{tuned }& \DIFadd{blobs }& \DIFadd{rng }& \DIFadd{20 }& \DIFadd{0.1283 }& \DIFadd{0.1254 }& \DIFadd{0.1269 }& \DIFadd{3.6686 }& \DIFadd{3.6021 }&
\DIFadd{3.6353 }& \DIFadd{1.0000 }& \DIFadd{1.0000 }& \DIFadd{1.0000 }& \DIFadd{0.0000 }& \DIFadd{0.0000 }& \DIFadd{0.0000 }\\
\DIFadd{tuned }& \DIFadd{breast\_cancer }& \DIFadd{hexagonal }& \DIFadd{20 }& \DIFadd{2.6571 }& \DIFadd{1.9210 }& \DIFadd{2.2891 }&
\DIFadd{35.9418 }& \DIFadd{32.7551 }& \DIFadd{34.3485 }& \DIFadd{0.7165 }& \DIFadd{0.9910 }& \DIFadd{0.8538 }& \DIFadd{0.2835 }& \DIFadd{0.0090
}& \DIFadd{0.1462 }\\
\DIFadd{tuned }& \DIFadd{breast\_cancer }& \DIFadd{mst }& \DIFadd{20 }& \DIFadd{2.6625 }& \DIFadd{1.8673 }& \DIFadd{2.2649 }& \DIFadd{15.7822 }&
\DIFadd{10.7114 }& \DIFadd{13.2468 }& \DIFadd{0.6955 }& \DIFadd{0.9925 }& \DIFadd{0.8440 }& \DIFadd{0.3045 }& \DIFadd{0.0075 }&
\DIFadd{0.1560 }\\
\DIFadd{tuned }& \DIFadd{breast\_cancer }& \DIFadd{rng }& \DIFadd{20 }& \DIFadd{2.6530 }& \DIFadd{1.8688 }& \DIFadd{2.2609 }& \DIFadd{10.5715 }&
\DIFadd{6.2888 }& \DIFadd{8.4302 }& \DIFadd{0.7120 }& \DIFadd{0.9985 }& \DIFadd{0.8553 }& \DIFadd{0.2880 }& \DIFadd{0.0015 }& \DIFadd{0.1447 }\\
\DIFadd{tuned }& \DIFadd{california\_housing }& \DIFadd{hexagonal }& \DIFadd{20 }& \DIFadd{0.7416 }& \DIFadd{0.7197 }& \DIFadd{0.7307
}& \DIFadd{36.2493 }& \DIFadd{36.0668 }& \DIFadd{36.1580 }& \DIFadd{0.9835 }& \DIFadd{1.0000 }& \DIFadd{0.9918 }& \DIFadd{0.0165 }&
\DIFadd{0.0000 }& \DIFadd{0.0083 }\\
\DIFadd{tuned }& \DIFadd{california\_housing }& \DIFadd{mst }& \DIFadd{20 }& \DIFadd{0.7384 }& \DIFadd{0.7150 }& \DIFadd{0.7267 }&
\DIFadd{10.4011 }& \DIFadd{10.1778 }& \DIFadd{10.2894 }& \DIFadd{0.9800 }& \DIFadd{1.0000 }& \DIFadd{0.9900 }& \DIFadd{0.0200 }& \DIFadd{0.0000
}& \DIFadd{0.0100 }\\
\DIFadd{tuned }& \DIFadd{california\_housing }& \DIFadd{rng }& \DIFadd{20 }& \DIFadd{0.7389 }& \DIFadd{0.7162 }& \DIFadd{0.7275 }&
\DIFadd{6.3594 }& \DIFadd{6.2047 }& \DIFadd{6.2821 }& \DIFadd{0.9790 }& \DIFadd{1.0000 }& \DIFadd{0.9895 }& \DIFadd{0.0210 }& \DIFadd{0.0000 }&
\DIFadd{0.0105 }\\
\DIFadd{tuned }& \DIFadd{circles }& \DIFadd{hexagonal }& \DIFadd{20 }& \DIFadd{0.1560 }& \DIFadd{0.1526 }& \DIFadd{0.1543 }& \DIFadd{31.0386 }&
\DIFadd{30.9675 }& \DIFadd{31.0030 }& \DIFadd{1.0000 }& \DIFadd{1.0000 }& \DIFadd{1.0000 }& \DIFadd{0.0000 }& \DIFadd{0.0000 }&
\DIFadd{0.0000 }\\
\DIFadd{tuned }& \DIFadd{circles }& \DIFadd{mst }& \DIFadd{20 }& \DIFadd{0.1566 }& \DIFadd{0.1525 }& \DIFadd{0.1545 }& \DIFadd{9.9500 }& \DIFadd{9.5644
}& \DIFadd{9.7572 }& \DIFadd{1.0000 }& \DIFadd{1.0000 }& \DIFadd{1.0000 }& \DIFadd{0.0000 }& \DIFadd{0.0000 }& \DIFadd{0.0000 }\\
\DIFadd{tuned }& \DIFadd{circles }& \DIFadd{rng }& \DIFadd{20 }& \DIFadd{0.1562 }& \DIFadd{0.1523 }& \DIFadd{0.1542 }& \DIFadd{3.7441 }& \DIFadd{3.6653
}& \DIFadd{3.7047 }& \DIFadd{1.0000 }& \DIFadd{1.0000 }& \DIFadd{1.0000 }& \DIFadd{0.0000 }& \DIFadd{0.0000 }& \DIFadd{0.0000 }\\
\DIFadd{tuned }& \DIFadd{covertype }& \DIFadd{hexagonal }& \DIFadd{20 }& \DIFadd{2.2611 }& \DIFadd{2.2604 }& \DIFadd{2.2608 }& \DIFadd{27.3364
}& \DIFadd{27.3276 }& \DIFadd{27.3320 }& \DIFadd{0.7955 }& \DIFadd{0.7965 }& \DIFadd{0.7960 }& \DIFadd{0.2045 }& \DIFadd{0.2035 }&
\DIFadd{0.2040 }\\
\DIFadd{tuned }& \DIFadd{covertype }& \DIFadd{mst }& \DIFadd{20 }& \DIFadd{2.2210 }& \DIFadd{2.2207 }& \DIFadd{2.2209 }& \DIFadd{3.1564 }&
\DIFadd{3.1564 }& \DIFadd{3.1564 }& \DIFadd{0.8725 }& \DIFadd{0.8715 }& \DIFadd{0.8720 }& \DIFadd{0.1275 }& \DIFadd{0.1285 }& \DIFadd{0.1280 }\\
\DIFadd{tuned }& \DIFadd{covertype }& \DIFadd{rng }& \DIFadd{20 }& \DIFadd{2.1854 }& \DIFadd{2.1835 }& \DIFadd{2.1844 }& \DIFadd{2.8206 }&
\DIFadd{2.8178 }& \DIFadd{2.8192 }& \DIFadd{0.8865 }& \DIFadd{0.8870 }& \DIFadd{0.8867 }& \DIFadd{0.1135 }& \DIFadd{0.1130 }& \DIFadd{0.1132 }\\
\DIFadd{tuned }& \DIFadd{diabetes }& \DIFadd{hexagonal }& \DIFadd{20 }& \DIFadd{1.6791 }& \DIFadd{1.0986 }& \DIFadd{1.3888 }& \DIFadd{41.7611 }&
\DIFadd{38.9399 }& \DIFadd{40.3505 }& \DIFadd{0.6960 }& \DIFadd{0.9800 }& \DIFadd{0.8380 }& \DIFadd{0.3040 }& \DIFadd{0.0200 }&
\DIFadd{0.1620 }\\
\DIFadd{tuned }& \DIFadd{diabetes }& \DIFadd{mst }& \DIFadd{20 }& \DIFadd{1.6773 }& \DIFadd{1.0906 }& \DIFadd{1.3840 }& \DIFadd{16.6558 }&
\DIFadd{9.5850 }& \DIFadd{13.1204 }& \DIFadd{0.6835 }& \DIFadd{0.9995 }& \DIFadd{0.8415 }& \DIFadd{0.3165 }& \DIFadd{5.00e-04 }&
\DIFadd{0.1585 }\\
\DIFadd{tuned }& \DIFadd{diabetes }& \DIFadd{rng }& \DIFadd{20 }& \DIFadd{1.6769 }& \DIFadd{1.0967 }& \DIFadd{1.3868 }& \DIFadd{12.2733 }&
\DIFadd{6.2114 }& \DIFadd{9.2424 }& \DIFadd{0.6825 }& \DIFadd{0.9970 }& \DIFadd{0.8397 }& \DIFadd{0.3175 }& \DIFadd{0.0030 }& \DIFadd{0.1603 }\\
\DIFadd{tuned }& \DIFadd{digits }& \DIFadd{hexagonal }& \DIFadd{20 }& \DIFadd{4.4170 }& \DIFadd{3.8544 }& \DIFadd{4.1357 }& \DIFadd{36.8784 }&
\DIFadd{36.8120 }& \DIFadd{36.8452 }& \DIFadd{0.9050 }& \DIFadd{0.9735 }& \DIFadd{0.9393 }& \DIFadd{0.0950 }& \DIFadd{0.0265 }&
\DIFadd{0.0607 }\\
\DIFadd{tuned }& \DIFadd{digits }& \DIFadd{mst }& \DIFadd{20 }& \DIFadd{4.3875 }& \DIFadd{3.8143 }& \DIFadd{4.1009 }& \DIFadd{7.6847 }& \DIFadd{6.6297 }&
\DIFadd{7.1572 }& \DIFadd{0.9415 }& \DIFadd{0.9970 }& \DIFadd{0.9692 }& \DIFadd{0.0585 }& \DIFadd{0.0030 }& \DIFadd{0.0307 }\\
\DIFadd{tuned }& \DIFadd{digits }& \DIFadd{rng }& \DIFadd{20 }& \DIFadd{4.3877 }& \DIFadd{3.8152 }& \DIFadd{4.1014 }& \DIFadd{6.3308 }& \DIFadd{5.2171 }&
\DIFadd{5.7740 }& \DIFadd{0.9365 }& \DIFadd{0.9955 }& \DIFadd{0.9660 }& \DIFadd{0.0635 }& \DIFadd{0.0045 }& \DIFadd{0.0340 }\\
\DIFadd{tuned }& \DIFadd{iris }& \DIFadd{hexagonal }& \DIFadd{20 }& \DIFadd{0.3650 }& \DIFadd{0.0874 }& \DIFadd{0.2262 }& \DIFadd{30.7106 }&
\DIFadd{24.6950 }& \DIFadd{27.7028 }& \DIFadd{0.3215 }& \DIFadd{0.7235 }& \DIFadd{0.5225 }& \DIFadd{0.6785 }& \DIFadd{0.2765 }&
\DIFadd{0.4775 }\\
\DIFadd{tuned }& \DIFadd{iris }& \DIFadd{mst }& \DIFadd{20 }& \DIFadd{0.3688 }& \DIFadd{0.0814 }& \DIFadd{0.2251 }& \DIFadd{5.6300 }& \DIFadd{2.0910 }&
\DIFadd{3.8605 }& \DIFadd{0.3220 }& \DIFadd{0.7600 }& \DIFadd{0.5410 }& \DIFadd{0.6780 }& \DIFadd{0.2400 }& \DIFadd{0.4590 }\\
\DIFadd{tuned }& \DIFadd{iris }& \DIFadd{rng }& \DIFadd{20 }& \DIFadd{0.3552 }& \DIFadd{0.0611 }& \DIFadd{0.2081 }& \DIFadd{4.5694 }& \DIFadd{1.8345 }&
\DIFadd{3.2020 }& \DIFadd{0.3315 }& \DIFadd{0.8020 }& \DIFadd{0.5668 }& \DIFadd{0.6685 }& \DIFadd{0.1980 }& \DIFadd{0.4332 }\\
\DIFadd{tuned }& \DIFadd{kddcup99 }& \DIFadd{hexagonal }& \DIFadd{20 }& \DIFadd{0.3047 }& \DIFadd{0.2982 }& \DIFadd{0.3014 }& \DIFadd{4.9441 }&
\DIFadd{4.9339 }& \DIFadd{4.9390 }& \DIFadd{0.9540 }& \DIFadd{0.9695 }& \DIFadd{0.9617 }& \DIFadd{0.0460 }& \DIFadd{0.0305 }& \DIFadd{0.0382 }\\
\DIFadd{tuned }& \DIFadd{kddcup99 }& \DIFadd{mst }& \DIFadd{20 }& \DIFadd{0.2893 }& \DIFadd{0.2826 }& \DIFadd{0.2860 }& \DIFadd{2.0398 }& \DIFadd{2.0344
}& \DIFadd{2.0371 }& \DIFadd{0.9620 }& \DIFadd{0.9780 }& \DIFadd{0.9700 }& \DIFadd{0.0380 }& \DIFadd{0.0220 }& \DIFadd{0.0300 }\\
\DIFadd{tuned }& \DIFadd{kddcup99 }& \DIFadd{rng }& \DIFadd{20 }& \DIFadd{0.2934 }& \DIFadd{0.2881 }& \DIFadd{0.2908 }& \DIFadd{1.8173 }& \DIFadd{1.8160
}& \DIFadd{1.8166 }& \DIFadd{0.9625 }& \DIFadd{0.9740 }& \DIFadd{0.9682 }& \DIFadd{0.0375 }& \DIFadd{0.0260 }& \DIFadd{0.0318 }\\
\DIFadd{tuned }& \DIFadd{moons }& \DIFadd{hexagonal }& \DIFadd{20 }& \DIFadd{0.1317 }& \DIFadd{0.1292 }& \DIFadd{0.1304 }& \DIFadd{28.7197 }&
\DIFadd{28.5455 }& \DIFadd{28.6326 }& \DIFadd{1.0000 }& \DIFadd{1.0000 }& \DIFadd{1.0000 }& \DIFadd{0.0000 }& \DIFadd{0.0000 }&
\DIFadd{0.0000 }\\
\DIFadd{tuned }& \DIFadd{moons }& \DIFadd{mst }& \DIFadd{20 }& \DIFadd{0.1315 }& \DIFadd{0.1288 }& \DIFadd{0.1301 }& \DIFadd{6.9319 }& \DIFadd{6.7384 }&
\DIFadd{6.8351 }& \DIFadd{1.0000 }& \DIFadd{1.0000 }& \DIFadd{1.0000 }& \DIFadd{0.0000 }& \DIFadd{0.0000 }& \DIFadd{0.0000 }\\
\DIFadd{tuned }& \DIFadd{moons }& \DIFadd{rng }& \DIFadd{20 }& \DIFadd{0.1315 }& \DIFadd{0.1285 }& \DIFadd{0.1300 }& \DIFadd{3.5884 }& \DIFadd{3.5222 }&
\DIFadd{3.5553 }& \DIFadd{1.0000 }& \DIFadd{1.0000 }& \DIFadd{1.0000 }& \DIFadd{0.0000 }& \DIFadd{0.0000 }& \DIFadd{0.0000 }\\
\DIFadd{tuned }& \DIFadd{olivetti\_faces }& \DIFadd{hexagonal }& \DIFadd{20 }& \DIFadd{41.7747 }& \DIFadd{27.3023 }& \DIFadd{34.5385 }&
\DIFadd{43.2448 }& \DIFadd{41.6179 }& \DIFadd{42.4314 }& \DIFadd{0.6285 }& \DIFadd{0.9430 }& \DIFadd{0.7857 }& \DIFadd{0.3715 }& \DIFadd{0.0570
}& \DIFadd{0.2142 }\\
\DIFadd{tuned }& \DIFadd{olivetti\_faces }& \DIFadd{mst }& \DIFadd{20 }& \DIFadd{41.3166 }& \DIFadd{26.3826 }& \DIFadd{33.8496 }&
\DIFadd{11.8723 }& \DIFadd{6.8992 }& \DIFadd{9.3857 }& \DIFadd{0.6380 }& \DIFadd{0.9770 }& \DIFadd{0.8075 }& \DIFadd{0.3620 }& \DIFadd{0.0230 }&
\DIFadd{0.1925 }\\
\DIFadd{tuned }& \DIFadd{olivetti\_faces }& \DIFadd{rng }& \DIFadd{20 }& \DIFadd{41.4380 }& \DIFadd{26.7705 }& \DIFadd{34.1043 }&
\DIFadd{8.7106 }& \DIFadd{4.3695 }& \DIFadd{6.5400 }& \DIFadd{0.6360 }& \DIFadd{0.9640 }& \DIFadd{0.8000 }& \DIFadd{0.3640 }& \DIFadd{0.0360 }&
\DIFadd{0.2000 }\\
\DIFadd{tuned }& \DIFadd{s\_curve }& \DIFadd{hexagonal }& \DIFadd{20 }& \DIFadd{0.3070 }& \DIFadd{0.3016 }& \DIFadd{0.3043 }& \DIFadd{27.9434 }&
\DIFadd{27.9687 }& \DIFadd{27.9560 }& \DIFadd{1.0000 }& \DIFadd{1.0000 }& \DIFadd{1.0000 }& \DIFadd{0.0000 }& \DIFadd{0.0000 }&
\DIFadd{0.0000 }\\
\DIFadd{tuned }& \DIFadd{s\_curve }& \DIFadd{mst }& \DIFadd{20 }& \DIFadd{0.3060 }& \DIFadd{0.3005 }& \DIFadd{0.3032 }& \DIFadd{8.9435 }& \DIFadd{8.6914
}& \DIFadd{8.8175 }& \DIFadd{1.0000 }& \DIFadd{1.0000 }& \DIFadd{1.0000 }& \DIFadd{0.0000 }& \DIFadd{0.0000 }& \DIFadd{0.0000 }\\
\DIFadd{tuned }& \DIFadd{s\_curve }& \DIFadd{rng }& \DIFadd{20 }& \DIFadd{0.3059 }& \DIFadd{0.3005 }& \DIFadd{0.3032 }& \DIFadd{3.9253 }& \DIFadd{3.8508
}& \DIFadd{3.8881 }& \DIFadd{1.0000 }& \DIFadd{1.0000 }& \DIFadd{1.0000 }& \DIFadd{0.0000 }& \DIFadd{0.0000 }& \DIFadd{0.0000 }\\
\DIFadd{tuned }& \DIFadd{swiss\_roll }& \DIFadd{hexagonal }& \DIFadd{20 }& \DIFadd{0.2472 }& \DIFadd{0.2430 }& \DIFadd{0.2451 }&
\DIFadd{32.9706 }& \DIFadd{32.9121 }& \DIFadd{32.9413 }& \DIFadd{0.9665 }& \DIFadd{0.9665 }& \DIFadd{0.9665 }& \DIFadd{0.0335 }& \DIFadd{0.0335
}& \DIFadd{0.0335 }\\
\DIFadd{tuned }& \DIFadd{swiss\_roll }& \DIFadd{mst }& \DIFadd{20 }& \DIFadd{0.2437 }& \DIFadd{0.2394 }& \DIFadd{0.2415 }& \DIFadd{9.2023 }&
\DIFadd{8.9054 }& \DIFadd{9.0538 }& \DIFadd{0.9955 }& \DIFadd{0.9955 }& \DIFadd{0.9955 }& \DIFadd{0.0045 }& \DIFadd{0.0045 }& \DIFadd{0.0045 }\\
\DIFadd{tuned }& \DIFadd{swiss\_roll }& \DIFadd{rng }& \DIFadd{20 }& \DIFadd{0.2427 }& \DIFadd{0.2388 }& \DIFadd{0.2407 }& \DIFadd{3.6234 }&
\DIFadd{3.5686 }& \DIFadd{3.5960 }& \DIFadd{0.9920 }& \DIFadd{0.9920 }& \DIFadd{0.9920 }& \DIFadd{0.0080 }& \DIFadd{0.0080 }& \DIFadd{0.0080 }\\
\DIFadd{tuned }& \DIFadd{wine }& \DIFadd{hexagonal }& \DIFadd{20 }& \DIFadd{1.9856 }& \DIFadd{0.6770 }& \DIFadd{1.3313 }& \DIFadd{32.2560 }&
\DIFadd{27.5770 }& \DIFadd{29.9165 }& \DIFadd{0.3710 }& \DIFadd{0.7955 }& \DIFadd{0.5833 }& \DIFadd{0.6290 }& \DIFadd{0.2045 }&
\DIFadd{0.4168 }\\
\DIFadd{tuned }& \DIFadd{wine }& \DIFadd{mst }& \DIFadd{20 }& \DIFadd{1.9633 }& \DIFadd{0.5493 }& \DIFadd{1.2563 }& \DIFadd{10.4343 }& \DIFadd{3.0798 }&
\DIFadd{6.7570 }& \DIFadd{0.3760 }& \DIFadd{0.8790 }& \DIFadd{0.6275 }& \DIFadd{0.6240 }& \DIFadd{0.1210 }& \DIFadd{0.3725 }\\
\DIFadd{tuned }& \DIFadd{wine }& \DIFadd{rng }& \DIFadd{20 }& \DIFadd{1.9598 }& \DIFadd{0.4868 }& \DIFadd{1.2233 }& \DIFadd{9.8250 }& \DIFadd{2.6490 }&
\DIFadd{6.2370 }& \DIFadd{0.3835 }& \DIFadd{0.9160 }& \DIFadd{0.6498 }& \DIFadd{0.6165 }& \DIFadd{0.0840 }& \DIFadd{0.3502 }\\
\end{longtable}
}

\DIFadd{\textbf{Supplementary Table S14. Matched local CuPy and Ray streaming
execution-path diagnostics.} Values are paired differences computed as
Ray streaming minus local CuPy for matched dataset, seed, topology, and
tuned configuration. Each row summarizes the paired differences across
14 datasets and 10 seeds. The table reports the mean paired difference,
standard deviation of the paired difference, 95\% confidence interval,
and Benjamini-Hochberg adjusted q-value.
}

{\def\LTcaptype{none} %DIF >  do not increment counter
\begin{longtable}[]{@{}
  >{\raggedright\arraybackslash}p{(\linewidth - 10\tabcolsep) * \real{0.1364}}
  >{\raggedright\arraybackslash}p{(\linewidth - 10\tabcolsep) * \real{0.1364}}
  >{\raggedleft\arraybackslash}p{(\linewidth - 10\tabcolsep) * \real{0.1818}}
  >{\raggedleft\arraybackslash}p{(\linewidth - 10\tabcolsep) * \real{0.1818}}
  >{\raggedleft\arraybackslash}p{(\linewidth - 10\tabcolsep) * \real{0.1818}}
  >{\raggedleft\arraybackslash}p{(\linewidth - 10\tabcolsep) * \real{0.1818}}@{}}
\toprule\noalign{}
\begin{minipage}[b]{\linewidth}\raggedright
\DIFadd{Topology
}\end{minipage} & \begin{minipage}[b]{\linewidth}\raggedright
\DIFadd{Metric
}\end{minipage} & \begin{minipage}[b]{\linewidth}\raggedleft
\DIFadd{Mean difference
}\end{minipage} & \begin{minipage}[b]{\linewidth}\raggedleft
\DIFadd{SD difference
}\end{minipage} & \begin{minipage}[b]{\linewidth}\raggedleft
\DIFadd{95\% CI
}\end{minipage} & \begin{minipage}[b]{\linewidth}\raggedleft
\DIFadd{BH q
}\end{minipage} \\
\midrule\noalign{}
\endhead
\bottomrule\noalign{}
\endlastfoot
\DIFadd{hexagonal }& \DIFadd{Balanced QE }& \DIFadd{3.08e-05 }& \DIFadd{3.00e-04 }& {[}\DIFadd{-1.90e-05,
8.05e-05}{]} & \DIFadd{0.7840 }\\
\DIFadd{hexagonal }& \DIFadd{Balanced MTR }& \DIFadd{-0.0121 }& \DIFadd{0.1526 }& {[}\DIFadd{-0.0374, 0.0132}{]} &
\DIFadd{0.7840 }\\
\DIFadd{hexagonal }& \DIFadd{Balanced node utilization }& \DIFadd{-1.07e-04 }& \DIFadd{0.0013 }&
{[}\DIFadd{-3.17e-04, 1.03e-04}{]} & \DIFadd{0.7840 }\\
\DIFadd{hexagonal }& \DIFadd{Balanced dead-node fraction }& \DIFadd{1.07e-04 }& \DIFadd{0.0013 }&
{[}\DIFadd{-1.03e-04, 3.17e-04}{]} & \DIFadd{0.7840 }\\
\DIFadd{MST }& \DIFadd{Balanced QE }& \DIFadd{-2.96e-04 }& \DIFadd{0.0043 }& {[}\DIFadd{-0.0010, 4.15e-04}{]} &
\DIFadd{0.7840 }\\
\DIFadd{MST }& \DIFadd{Balanced MTR }& \DIFadd{-0.0153 }& \DIFadd{0.5061 }& {[}\DIFadd{-0.0991, 0.0685}{]} &
\DIFadd{0.8317 }\\
\DIFadd{MST }& \DIFadd{Balanced node utilization }& \DIFadd{3.21e-04 }& \DIFadd{0.0052 }& {[}\DIFadd{-5.40e-04,
0.0012}{]} & \DIFadd{0.7840 }\\
\DIFadd{MST }& \DIFadd{Balanced dead-node fraction }& \DIFadd{-3.21e-04 }& \DIFadd{0.0052 }& {[}\DIFadd{-0.0012,
5.40e-04}{]} & \DIFadd{0.7840 }\\
\DIFadd{RNG }& \DIFadd{Balanced QE }& \DIFadd{0.0049 }& \DIFadd{0.1260 }& {[}\DIFadd{-0.0160, 0.0258}{]} & \DIFadd{0.8294 }\\
\DIFadd{RNG }& \DIFadd{Balanced MTR }& \DIFadd{-0.0823 }& \DIFadd{0.5164 }& {[}\DIFadd{-0.1679, 0.0032}{]} &
\DIFadd{0.7840 }\\
\DIFadd{RNG }& \DIFadd{Balanced node utilization }& \DIFadd{-1.07e-04 }& \DIFadd{0.0035 }& {[}\DIFadd{-6.90e-04,
4.76e-04}{]} & \DIFadd{0.8317 }\\
\DIFadd{RNG }& \DIFadd{Balanced dead-node fraction }& \DIFadd{1.07e-04 }& \DIFadd{0.0035 }& {[}\DIFadd{-4.76e-04,
6.90e-04}{]} & \DIFadd{0.8317 }\\
\end{longtable}
}

\DIFadd{\textbf{Supplementary Table S15. Post hoc matched cross-topology
balanced-QE contrasts across the initial-radius sweep.} For each
topology pair and radius, seed-level log \(QE_B\) ratios were averaged
within each dataset and then tested across the 14 equally weighted
dataset effects. \texttt{QE\ ratio} is the exponentiated mean log ratio
(comparator/reference), so values below 1 and positive improvement
values favor the comparator. Confidence intervals are exponentiated 95\%
paired \(t\) intervals. Benjamini--Hochberg q-values adjust across all
21 topology-pair--radius tests.
}

{\def\LTcaptype{none} %DIF >  do not increment counter
\begin{longtable}[]{@{}
  >{\raggedright\arraybackslash}p{(\linewidth - 16\tabcolsep) * \real{0.0882}}
  >{\raggedleft\arraybackslash}p{(\linewidth - 16\tabcolsep) * \real{0.1176}}
  >{\raggedleft\arraybackslash}p{(\linewidth - 16\tabcolsep) * \real{0.1176}}
  >{\raggedleft\arraybackslash}p{(\linewidth - 16\tabcolsep) * \real{0.1176}}
  >{\raggedleft\arraybackslash}p{(\linewidth - 16\tabcolsep) * \real{0.1176}}
  >{\raggedright\arraybackslash}p{(\linewidth - 16\tabcolsep) * \real{0.0882}}
  >{\raggedleft\arraybackslash}p{(\linewidth - 16\tabcolsep) * \real{0.1176}}
  >{\raggedleft\arraybackslash}p{(\linewidth - 16\tabcolsep) * \real{0.1176}}
  >{\raggedleft\arraybackslash}p{(\linewidth - 16\tabcolsep) * \real{0.1176}}@{}}
\toprule\noalign{}
\begin{minipage}[b]{\linewidth}\raggedright
\DIFadd{Contrast
}\end{minipage} & \begin{minipage}[b]{\linewidth}\raggedleft
\DIFadd{Radius
}\end{minipage} & \begin{minipage}[b]{\linewidth}\raggedleft
\DIFadd{n datasets
}\end{minipage} & \begin{minipage}[b]{\linewidth}\raggedleft
\DIFadd{n seed pairs
}\end{minipage} & \begin{minipage}[b]{\linewidth}\raggedleft
\DIFadd{QE ratio
}\end{minipage} & \begin{minipage}[b]{\linewidth}\raggedright
\DIFadd{95\% CI
}\end{minipage} & \begin{minipage}[b]{\linewidth}\raggedleft
\DIFadd{Improvement
}\end{minipage} & \begin{minipage}[b]{\linewidth}\raggedleft
\DIFadd{raw p
}\end{minipage} & \begin{minipage}[b]{\linewidth}\raggedleft
\DIFadd{BH q
}\end{minipage} \\
\midrule\noalign{}
\endhead
\bottomrule\noalign{}
\endlastfoot
\DIFadd{MST vs Hex }& \DIFadd{0.5 }& \DIFadd{14 }& \DIFadd{280 }& \DIFadd{0.9969 }& {[}\DIFadd{0.9876, 1.0062}{]} & \DIFadd{0.31\% }&
\DIFadd{0.4801 }& \DIFadd{0.5930 }\\
\DIFadd{RNG vs Hex }& \DIFadd{0.5 }& \DIFadd{14 }& \DIFadd{280 }& \DIFadd{0.9880 }& {[}\DIFadd{0.9726, 1.0036}{]} & \DIFadd{1.20\% }&
\DIFadd{0.1193 }& \DIFadd{0.1927 }\\
\DIFadd{RNG vs MST }& \DIFadd{0.5 }& \DIFadd{14 }& \DIFadd{280 }& \DIFadd{0.9911 }& {[}\DIFadd{0.9803, 1.0020}{]} & \DIFadd{0.89\% }&
\DIFadd{0.1019 }& \DIFadd{0.1784 }\\
\DIFadd{MST vs Hex }& \DIFadd{0.75 }& \DIFadd{14 }& \DIFadd{280 }& \DIFadd{0.9992 }& {[}\DIFadd{0.9880, 1.0105}{]} & \DIFadd{0.08\% }&
\DIFadd{0.8745 }& \DIFadd{0.8745 }\\
\DIFadd{RNG vs Hex }& \DIFadd{0.75 }& \DIFadd{14 }& \DIFadd{280 }& \DIFadd{0.9930 }& {[}\DIFadd{0.9760, 1.0103}{]} & \DIFadd{0.70\% }&
\DIFadd{0.3961 }& \DIFadd{0.5199 }\\
\DIFadd{RNG vs MST }& \DIFadd{0.75 }& \DIFadd{14 }& \DIFadd{280 }& \DIFadd{0.9938 }& {[}\DIFadd{0.9870, 1.0007}{]} & \DIFadd{0.62\% }&
\DIFadd{0.0751 }& \DIFadd{0.1434 }\\
\DIFadd{MST vs Hex }& \DIFadd{1.027 }& \DIFadd{14 }& \DIFadd{280 }& \DIFadd{0.9849 }& {[}\DIFadd{0.9732, 0.9968}{]} & \DIFadd{1.51\% }&
\DIFadd{0.0172 }& \DIFadd{0.0361 }\\
\DIFadd{RNG vs Hex }& \DIFadd{1.027 }& \DIFadd{14 }& \DIFadd{280 }& \DIFadd{0.9815 }& {[}\DIFadd{0.9674, 0.9957}{]} & \DIFadd{1.85\% }&
\DIFadd{0.0150 }& \DIFadd{0.0351 }\\
\DIFadd{RNG vs MST }& \DIFadd{1.027 }& \DIFadd{14 }& \DIFadd{280 }& \DIFadd{0.9965 }& {[}\DIFadd{0.9847, 1.0084}{]} & \DIFadd{0.35\% }&
\DIFadd{0.5329 }& \DIFadd{0.5943 }\\
\DIFadd{MST vs Hex }& \DIFadd{1.5 }& \DIFadd{14 }& \DIFadd{280 }& \DIFadd{0.9670 }& {[}\DIFadd{0.9458, 0.9887}{]} & \DIFadd{3.30\% }&
\DIFadd{0.0062 }& \DIFadd{0.0162 }\\
\DIFadd{RNG vs Hex }& \DIFadd{1.5 }& \DIFadd{14 }& \DIFadd{280 }& \DIFadd{0.9633 }& {[}\DIFadd{0.9412, 0.9859}{]} & \DIFadd{3.67\% }&
\DIFadd{0.0041 }& \DIFadd{0.0141 }\\
\DIFadd{RNG vs MST }& \DIFadd{1.5 }& \DIFadd{14 }& \DIFadd{280 }& \DIFadd{0.9962 }& {[}\DIFadd{0.9834, 1.0092}{]} & \DIFadd{0.38\% }&
\DIFadd{0.5377 }& \DIFadd{0.5943 }\\
\DIFadd{MST vs Hex }& \DIFadd{2 }& \DIFadd{14 }& \DIFadd{280 }& \DIFadd{0.9536 }& {[}\DIFadd{0.9263, 0.9817}{]} & \DIFadd{4.64\% }&
\DIFadd{0.0036 }& \DIFadd{0.0141 }\\
\DIFadd{RNG vs Hex }& \DIFadd{2 }& \DIFadd{14 }& \DIFadd{280 }& \DIFadd{0.9516 }& {[}\DIFadd{0.9248, 0.9791}{]} & \DIFadd{4.84\% }&
\DIFadd{0.0024 }& \DIFadd{0.0141 }\\
\DIFadd{RNG vs MST }& \DIFadd{2 }& \DIFadd{14 }& \DIFadd{280 }& \DIFadd{0.9979 }& {[}\DIFadd{0.9870, 1.0090}{]} & \DIFadd{0.21\% }&
\DIFadd{0.6871 }& \DIFadd{0.7214 }\\
\DIFadd{MST vs Hex }& \DIFadd{3 }& \DIFadd{14 }& \DIFadd{280 }& \DIFadd{0.9243 }& {[}\DIFadd{0.8793, 0.9716}{]} & \DIFadd{7.57\% }&
\DIFadd{0.0047 }& \DIFadd{0.0141 }\\
\DIFadd{RNG vs Hex }& \DIFadd{3 }& \DIFadd{14 }& \DIFadd{280 }& \DIFadd{0.9340 }& {[}\DIFadd{0.8989, 0.9705}{]} & \DIFadd{6.60\% }&
\DIFadd{0.0020 }& \DIFadd{0.0141 }\\
\DIFadd{RNG vs MST }& \DIFadd{3 }& \DIFadd{14 }& \DIFadd{280 }& \DIFadd{1.0105 }& {[}\DIFadd{0.9943, 1.0269}{]} & \DIFadd{-1.05\% }&
\DIFadd{0.1873 }& \DIFadd{0.2809 }\\
\DIFadd{MST vs Hex }& \DIFadd{5 }& \DIFadd{14 }& \DIFadd{280 }& \DIFadd{0.8842 }& {[}\DIFadd{0.8224, 0.9506}{]} & \DIFadd{11.58\% }&
\DIFadd{0.0028 }& \DIFadd{0.0141 }\\
\DIFadd{RNG vs Hex }& \DIFadd{5 }& \DIFadd{14 }& \DIFadd{280 }& \DIFadd{0.8919 }& {[}\DIFadd{0.8416, 0.9452}{]} & \DIFadd{10.81\% }&
\DIFadd{0.0009 }& \DIFadd{0.0141 }\\
\DIFadd{RNG vs MST }& \DIFadd{5 }& \DIFadd{14 }& \DIFadd{280 }& \DIFadd{1.0088 }& {[}\DIFadd{0.9923, 1.0255}{]} & \DIFadd{-0.88\% }&
\DIFadd{0.2731 }& \DIFadd{0.3824 }\\
\end{longtable}
}

\DIFaddend \subsection{12. Supplementary Figures (End
Matter)}\label{supplementary-figures-end-matter}

\pandocbounded{\includesvg[keepaspectratio]{assets_manual/figures/supp_fig_s1.svg}}
\emph{Supplementary Figure S1. \DIFaddbegin \DIFadd{Untuned }\DIFaddend FloatSOM \DIFaddbegin \DIFadd{MST }\DIFaddend versus \DIFaddbegin \DIFadd{hexagonal
}\DIFaddend XPySOM . Panels A-C report paired \(QE\) effects for \(QE_B\), \(QE_H\),
and \(QE_T\). Panel D reports dataset-level median runtime deltas
against dataset size, where each numbered dot is the median matched-seed
value of \texttt{FloatSOM\ time\ -\ XPySOM\ time}; negative values favor
FloatSOM and positive values favor XPySOM. The point numbers map to
Supplementary Table S4. Forest whiskers denote \DIFaddbegin \DIFadd{exact distribution-free
}\DIFaddend 95\%  intervals  \DIFaddbegin \DIFadd{for }\DIFaddend the  paired  \DIFaddbegin \DIFadd{Wilcoxon location estimate}\DIFaddend .}

\pandocbounded{\includesvg[keepaspectratio]{assets_manual/figures/supp_fig_s2.svg}}
\emph{Supplementary Figure S2. \DIFaddbegin \DIFadd{Untuned }\DIFaddend FloatSOM \DIFaddbegin \DIFadd{RNG }\DIFaddend versus \DIFaddbegin \DIFadd{hexagonal
}\DIFaddend XPySOM . Panels A-C report paired \(QE\) effects for \(QE_B\), \(QE_H\),
and \(QE_T\). Panel D reports dataset-level median runtime deltas
against dataset size, where each numbered dot is the median matched-seed
value of \texttt{FloatSOM\ time\ -\ XPySOM\ time}; negative values favor
FloatSOM and positive values favor XPySOM. The point numbers map to
Supplementary Table S5. Forest whiskers denote \DIFaddbegin \DIFadd{exact distribution-free
}\DIFaddend 95\%  intervals  \DIFaddbegin \DIFadd{for }\DIFaddend the  paired  \DIFaddbegin \DIFadd{Wilcoxon location estimate}\DIFaddend .}

\pandocbounded{\includesvg[keepaspectratio]{assets_manual/figures/supp_fig_s3.svg}}
\emph{Supplementary Figure S3.  \DIFaddbegin \DIFadd{Hexagonal FloatSOM--XPySOM implementation
}\DIFaddend calibration . Panels A-C report paired \(QE\) effects for \(QE_B\),
\(QE_H\), and \(QE_T\). Panel D reports dataset-level median runtime
deltas against dataset size, where each numbered dot is the median
matched-seed value of \texttt{FloatSOM\ time\ -\ XPySOM\ time}; negative
values favor FloatSOM and positive values favor XPySOM. The point
numbers map to Supplementary Table S6. Forest whiskers denote \DIFaddbegin \DIFadd{exact
distribution-free }\DIFaddend 95\%  intervals  \DIFaddbegin \DIFadd{for }\DIFaddend the  paired  \DIFaddbegin \DIFadd{Wilcoxon location
estimate}\DIFaddend .}

\pandocbounded{\includesvg[keepaspectratio]{assets_manual/figures/supp_fig_s4.svg}}
\emph{Supplementary Figure S4.  \DIFaddbegin \DIFadd{Topology sensitivity analyses }\DIFaddend under full
sampling .  \DIFaddbegin \DIFadd{A: hexagonal versus MST; B: hexagonal versus RNG; C: MST
versus RNG. Within each row, columns }\DIFaddend report \DIFaddbegin \DIFadd{matched top-\(k\) }\DIFaddend paired
 \DIFaddbegin \DIFadd{sensitivity analyses }\DIFaddend for \(QE_B\), \(QE_H\), and \(QE_T\)  \DIFaddbegin \DIFadd{at
\(k\in\{1,3,5,10\}\); directional labels indicate which topology is
favored}\DIFaddend . }

\pandocbounded{\includesvg[keepaspectratio]{assets_manual/figures/supp_fig_s5.svg}}
\emph{Supplementary Figure S5.  \DIFaddbegin \DIFadd{Tuned }\DIFaddend versus  \DIFaddbegin \DIFadd{untuned fixed-configuration
comparisons by }\DIFaddend topology .  \DIFaddbegin \DIFadd{A: hexagonal; B: MST; C: RNG. Within each row,
columns }\DIFaddend report  \(QE_B\), \(QE_H\), and \(QE_T\) \DIFaddbegin \DIFadd{for 14 datasets and 20
seeds under matched dataset--seed--topology--split keys}\DIFaddend . \DIFaddbegin \DIFadd{Positive values
indicate that the tuned configuration outperforms the untuned reference;
profile definitions are given in Section 4.3.1.}\DIFaddend }

\pandocbounded{\includesvg[keepaspectratio]{assets_manual/figures/supp_fig_s6.svg}}
\emph{Supplementary Figure S6.  \DIFaddbegin \DIFadd{Full GPU-count scaling context for MST
and hexagonal }\DIFaddend under \DIFaddbegin \DIFadd{matched }\DIFaddend full sampling  \DIFaddbegin \DIFadd{settings}\DIFaddend .  A-C \DIFaddbegin \DIFadd{: MST dimension}\DIFaddend ,
 \DIFaddbegin \DIFadd{sample}\DIFaddend , and  \DIFaddbegin \DIFadd{grid-size scaling; D-F: hexagonal dimension, sample, and
grid-size scaling}\DIFaddend . \DIFaddbegin \DIFadd{Curves correspond to \(G\in\{1,2,4,8\}\) GPUs and
report mean wall-clock runtime (s) with \(\pm 1\) standard-deviation
error bars across \(n=3\) repeated runs per configuration.}\DIFaddend }

\pandocbounded{\includesvg[keepaspectratio]{assets_manual/figures/supp_fig_s7.svg}}
\emph{Supplementary Figure S7.  \DIFaddbegin \DIFadd{Deployment comparison of untuned
hexagonal XPySOM }\DIFaddend versus  \DIFaddbegin \DIFadd{tuned FloatSOM hexagonal}\DIFaddend .  \DIFaddbegin \DIFadd{A: \(QE_B\); B:
\(QE_H\); C: \(QE_T\) under }\DIFaddend the matched  \DIFaddbegin \DIFadd{dataset/seed comparison keys.
Positive values indicate tuned FloatSOM hexagonal outperforms untuned
hexagonal XPySOM. Forest whiskers denote 95\% }\DIFaddend paired  \DIFaddbegin \DIFadd{\(t\)-test
confidence intervals around the mean paired effect. Per-dataset }\DIFaddend and
 \DIFaddbegin \DIFadd{\texttt{GLOBAL\_OVERALL} panel summaries are listed in Supplementary
Table S9}\DIFaddend .}

\pandocbounded{\includesvg[keepaspectratio]{assets_manual/figures/supp_fig_s8.svg}}
\emph{Supplementary Figure S8.  \DIFaddbegin \DIFadd{Deployment }\DIFaddend comparison  \DIFaddbegin \DIFadd{of untuned
}\DIFaddend hexagonal  \DIFaddbegin \DIFadd{XPySOM versus tuned FloatSOM MST. A: }\DIFaddend \(QE_B\) \DIFaddbegin \DIFadd{; B: }\DIFaddend \(QE_H\) \DIFaddbegin \DIFadd{; C:
}\DIFaddend \(QE_T\) under the matched  \DIFaddbegin \DIFadd{dataset/seed comparison }\DIFaddend keys. Positive values
indicate  tuned  \DIFaddbegin \DIFadd{FloatSOM MST }\DIFaddend outperforms  untuned  \DIFaddbegin \DIFadd{hexagonal XPySOM}\DIFaddend . Forest
whiskers denote 95\% paired \(t\)-test confidence intervals around the
mean paired effect. \DIFaddbegin \DIFadd{Per-dataset and \texttt{GLOBAL\_OVERALL} panel
summaries are listed in Supplementary Table S10.}\DIFaddend }

\end{document}
